\chapter{Birational and carrier reflection}
\index{reflection!birational}\index{reflection!carrier}
\label{chap:bj-reflection}

\section{Orientation}
\label{sec:ch39-orientation}

\subsection{Pullback is formal; reflection is not}
\label{sec:ch39-what}

Let $\beta:Y\to X$ be a projective birational morphism of normal integral
projective varieties over the trivially valued field $k$, and let
$\theta\in N^1(X)_{\mathbf R}$.  Pulling a $\theta$-psh function back
along $\beta$ is a matter of composing with a continuous map: it moves an
approximating net of Fubini--Study functions to an approximating net of
Fubini--Study functions, it preserves every order relation, and it costs
nothing.  We prove it in \Cref{thm:ch39-pullback} and then never think about
it again.

The reverse operation is the subject of this chapter, and it is not formal.
Given $\psi\in\PSH_Y(\beta^*\theta)$, one asks for
$\chi\in\PSH_X(\theta)$ with $\beta^*\chi=\psi$.  Three obstacles have to
be cleared, and they are genuinely different from one another.

\emph{First, the candidate has to be produced.}  The map
$\beta^{\An}:Y^{\An}\to X^{\An}$ is surjective, so a function on
$X^{\An}$ whose pullback is $\psi$ is unique if it exists, and its values
are forced: $\chi(x)=\psi(w)$ for any $w$ with $\beta^{\An}(w)=x$.  But
this is a definition only if $\psi$ is constant on the fibres of
$\beta^{\An}$, and the fibres are large.  Over a point of $X^{\An}$ whose
support is a closed point $p$ blown up by $\beta$, the fibre is the whole
analytification of the exceptional divisor above $p$
(\Cref{prop:ch39-fibres}).  Fibrewise constancy is therefore a real
restriction; it is \emph{false} for a general homogeneous psh function on
$Y$ in a class that is not pulled back from $X$
(\Cref{ex:ch39-not-fibre-constant}), and its validity in a pulled-back class
is a theorem, not a definition (\Cref{cor:ch39-fibre-constant}).

\emph{Second, the class has to be repaired.}  A $\beta^*\theta$-psh
function on $Y$ is a decreasing limit of Fubini--Study functions whose
classes converge to $\beta^*\theta$.  Those classes live in $N^1(Y)$ and
there is no reason for them to lie in the image of the injection
$\beta^*:N^1(X)\hookrightarrow N^1(Y)$; on a one-point blowup of a surface
they will in general have a nonzero coefficient along the exceptional class.
An approximating stage on $Y$ is therefore \emph{not} the pullback of
anything on $X$, and the passage downstairs must first cancel the drift.
The cancellation is performed in \Cref{lem:ch39-interpolation} by an exact
identity in $N^1$, not by an estimate; the reason it must be exact is
recorded in \Cref{warn:ch38-drift-not-absorbed}.

\emph{Third, the monotonicity has to survive the repair.}  The only limit
operation available in this chapter is the decreasing one
(\Cref{warn:ch39-firewall}).  The corrected data must therefore still be
ordered, and the interpolation that cancels the class drift must not destroy
the order in which the b-divisors increase.  That this is possible at all
depends on one inequality, $\beta^*F-G\ge0$, built into the fixed-pole
construction of \Cref{lem:ch39-fixed-pole}, and on the choice of the
interpolation parameter as a square root, \eqref{eq:bj-reflection-delta}.

The chapter proves reflection twice.  The homogeneous statement,
\Cref{thm:bj-homogeneous-reflection}, is unconditional: pullback is a
bijection
\[
 \beta^*:\PSH_{\mathrm{hom},X}(\theta)
 \xrightarrow{\ \sim\ }\PSH_{\mathrm{hom},Y}(\beta^*\theta)
\]
for every real class $\theta$, with no positivity hypothesis on
$\theta$ and no envelope hypothesis anywhere.  The inhomogeneous
statement, \Cref{thm:bj-carrier-reflection}, is conditional in a precise and
usable way: it holds for those $u\in\PSH_Y(\beta^*\theta)$ that are
dominated by a homogeneous function of the same class up to an additive
constant.  Only in \Cref{sec:ch39-descent}, where the envelope property is
assumed \emph{on the source} $Y$, is such a dominating function
manufactured for every $u$ at once.  The same source hypothesis is then
used a second time to produce the psh upper limit of an increasing net.
These are the chapter's only two uses of an envelope property, and in both
cases it is a displayed hypothesis, never an imported theorem.

\subsection{The firewall}
\label{sec:ch39-firewall}

\Cref{chap:bj-cartier-realization} stated a firewall in
\Cref{warn:ch38-firewall} because that chapter is one of the inputs of
\Cref{thm:bj-final-envelope}.  The present chapter is another such input,
and the danger is sharper here: \Cref{chap:bj-envelope-statement} cites
\Cref{thm:bj-birational-ep-descent} in its proof map, and
\Cref{chap:bj-alterations} cites it in the proof of the normal case.  Any
appeal here to the envelope theorem, however indirect, would close a loop.

\begin{warning}{The envelope theorem is not available in this chapter}
\label{warn:ch39-firewall}
Nothing in \Cref{chap:bj-reflection} may use, directly or transitively:
\begin{enumerate}
\item \Cref{thm:bj-final-envelope}, in any of its formulations;
\item the envelope property $\operatorname{EP}$ of
 \Cref{def:bj-envelope-property} as a \emph{theorem} for any class on any
 variety --- it appears below only as an explicit hypothesis, displayed in
 the statement of every result that carries it;
\item closure of $\PSH_X(\theta)$ under suprema, regularized or not, and
 in particular any appeal to an increasing net of psh functions having a
 psh limit;
\item the obstacle envelope $P_\theta$ of
 \Cref{def:bj-obstacle-envelope} regarded as an operator with values in
 $\PSH_X(\theta)$;
\item \Cref{thm:ch35-cartier-psh-criterion} in its direction
 $\Leftarrow$, and the consequences
 \Cref{thm:bj-nef-bdivisor-correspondence} and
 \Cref{cor:ch35-monotone-approximation} that rest on it;
\item any statement of \Cref{chap:bj-envelope-consequences}.
\end{enumerate}
Every limit taken below is a decreasing limit, and every supremum that
appears is either a supremum over a \emph{finite} set, or a supremum over a
compact parameter interval which is immediately rewritten as a decreasing
limit by \Cref{lem:bj-compact-legendre-window}, or the canonical-carrier
supremum \eqref{eq:bj-canonical-carrier}, which is psh precisely because the
envelope property has been \emph{assumed} at that point, or the
increasing-net supremum \eqref{eq:bj-source-net-limit}, whose psh extension
is supplied by the same source hypothesis through
\Cref{prop:bj-increasing-net-criterion}.
\end{warning}

Two clarifications keep the firewall usable rather than paralysing.

First, \Cref{prop:bj-increasing-net-criterion} --- the equivalence between
the family form and the increasing-net form of the envelope property --- is
\emph{not} forbidden.  It is a formal equivalence between two statements
about a fixed pair $(X,\theta)$; its proof in
\Cref{chap:bj-envelope-statement} uses divisorial determination and the
directedness of the set of finite subsets of an index set, and it does not
use \Cref{thm:bj-final-envelope}.  We use it in
\Cref{thm:bj-birational-ep-descent} to convert an assumed envelope property
upstairs into an assumed-free conclusion downstairs.

Second, the homogeneous Siu decomposition
\Cref{thm:ch35-homogeneous-siu-general} is not forbidden either.  Its
statement contains the sentence ``no envelope hypothesis is required'', and
\Cref{chap:bj-homogeneity-bdivisors} proves the strict case in full and
imports only the borderline case; the warning printed after it in that
chapter distinguishes it from the \emph{inhomogeneous} Siu statement, which
does need the envelope property and which we never use.
\Cref{ex:ch39-reflection-borderline} shows that the borderline case really
does occur in the application below, so the distinction is not academic.

\begin{remark}[Why the firewall is not vacuous here]
\label{rem:ch39-firewall-not-vacuous}
A reader who knows the finished theory will observe that once
\Cref{thm:bj-final-envelope} is available, \Cref{thm:bj-carrier-reflection}
becomes a triviality: every $u\in\PSH_Y(\beta^*\theta)$ then has the
canonical carrier of \Cref{lem:bj-canonical-carrier}, and the carrier
hypothesis is automatic.  This is exactly the reason the hypothesis has to
be carried explicitly through the present chapter.  The logical order is
\[
 \begin{gathered}
 \text{homogeneous reflection}
 \ \Rightarrow\
 \text{carrier reflection}
 \ \Rightarrow\
 \operatorname{EP}\text{-descent}\\
 \Rightarrow\ \cdots\ \Rightarrow\
 \Cref{thm:bj-final-envelope}
 \ \Rightarrow\
 \text{carriers exist always},
 \end{gathered}
\]
and reading the last implication backwards into the second is the standard
way to produce a circular proof of the envelope theorem.  The same trap in
its b-divisor form is described in
\Cref{rem:ch38-firewall-not-vacuous}.
\end{remark}

\subsection{The five tasks, and where each is discharged}
\label{sec:ch39-payoff}

\begin{enumerate}
\item\label{itm:ch39-task:hom}
 \emph{Homogeneous birational reflection.}  For every
 $\theta\in N^1(X)_{\mathbf R}$, pullback is a bijection from
 $\PSH_{\mathrm{hom},X}(\theta)$ onto
 $\PSH_{\mathrm{hom},Y}(\beta^*\theta)$.  This is
 \Cref{thm:bj-homogeneous-reflection}, proved in
 \Cref{sec:ch39-homogeneous} out of the realization theorem
 \Cref{thm:bj-real-cartier-realization}, the fixed-pole data
 \Cref{lem:ch39-fixed-pole}, and the exact drift identity
 \eqref{eq:bj-reflection-total}.  No positivity of $\theta$ is assumed.

\item\label{itm:ch39-task:carrier}
 \emph{Legendre carrier reflection.}  If $u\in\PSH_Y(\beta^*\theta)$
 admits a homogeneous carrier in the same class, then $u$ is a pullback.
 This is \Cref{thm:bj-carrier-reflection}, proved in
 \Cref{sec:ch39-legendre}.  The mechanism is the Legendre decomposition
 \Cref{thm:homogeneous-legendre-decomposition}: the carrier truncates the
 unbounded parameter interval to a compact window
 (\Cref{lem:ch35-carrier-tail}), and the compact window is converted from a
 supremum into a decreasing limit by the summable grid correction of
 \Cref{lem:bj-compact-legendre-window}.

\item\label{itm:ch39-task:determination}
 \emph{Equality on divisorial points and psh determination.}  Both proofs
 end by upgrading an identity on $X^{\divis}=Y^{\divis}$ to an identity on
 the whole analytification.  \Cref{sec:ch39-determination} isolates this
 step, states the exact hypotheses under which it is legitimate
 (\Cref{prop:ch39-determination-package}), and shows by
 \Cref{ex:ch39-usc-not-determined-birational} and
 \Cref{ex:ch39-two-usc-extensions} that dropping either hypothesis makes it
 false --- and false precisely at the proper-support points.

\item\label{itm:ch39-task:descent}
 \emph{Descent of the envelope property.}  If
 $\operatorname{EP}_Y(\beta^*\theta)$ holds then so does
 $\operatorname{EP}_X(\theta)$.  This is
 \Cref{thm:bj-birational-ep-descent}, proved in \Cref{sec:ch39-descent}
 through the canonical carrier \Cref{lem:bj-canonical-carrier}.

\item\label{itm:ch39-task:diagram}
 \emph{The exact functorial diagram.}  \Cref{sec:ch39-diagram} draws the
 square relating the four psh spaces and the implication chain governing
 the envelope property, with every arrow labelled as an equivalence, an
 injection, or a one-way implication, and with the direction that is
 \emph{not} proved marked as such.
\end{enumerate}

\subsection{Standing hypotheses and notation}
\label{sec:ch39-hypotheses}

Throughout the chapter, unless explicitly relaxed,
\[
 \begin{gathered}
 k=\bar k\ \text{trivially valued},\qquad
 X,\ Y\ \text{normal integral projective over }k,\\
 \beta:Y\longrightarrow X\ \text{projective birational},\qquad
 \theta\in N^1(X)_{\mathbf R}\ \text{arbitrary}.
 \end{gathered}
\tag{39.1}\label{eq:ch39-standing}
\]
We write
\[
                         \beta:Y\longrightarrow X
\tag{39.2}\label{eq:bj-reflection-map}
\]
for the map under discussion and $\beta^{\An}:Y^{\An}\to X^{\An}$ for the
induced map of analytifications.  Since $\beta$ is birational, the
function fields agree, $k(Y)=k(X)$, and this identification is used
without comment; it is what makes $X^{\divis}$ and $Y^{\divis}$ the same
set, and Cartier b-divisors over $X$ and over $Y$ the same objects
(\Cref{prop:ch38-evaluation-functorial}\,\cref{itm:ch38-fun:cofinal}).

The word \emph{class} always means a class in $N^1$ with real
coefficients.  We emphasise, following the correction recorded in
\Cref{def:bj-psh} and \Cref{lem:ch34-psh-nonempty}, that
\[
 \PSH_X(\theta)\ne\varnothing
 \ \Longrightarrow\ \theta\in\operatorname{Psef}(X),
\]
and that nefness of $\theta$ is equivalent to the constants being
$\theta$-psh and is therefore available only when a \emph{bounded} psh
function is in play.  Homogeneous psh functions are never bounded unless
they vanish (\Cref{prop:ch35-homogeneous-normalization} and the warning
after it), so \emph{no result of this chapter may assume that $\theta$ or
$\beta^*\theta$ is nef}.  This is not a stylistic scruple: the classes to
which \Cref{chap:bj-alterations} applies
\Cref{thm:bj-birational-ep-descent} are arbitrary real classes on a normal
variety, pulled back along an alteration, and there is no mechanism in that
argument for making them nef.

Because of this, the closure calculus of \Cref{thm:bj-basic-closure} ---
which already holds for arbitrary numerical classes, with no positivity
hypothesis --- is supplemented in \Cref{lem:ch39-closure-any-class} by the
monotonicity of $\PSH_X(\theta)$ in $\theta$, which that theorem does not
contain, and it is the supplemented form which is used below.  As in
\Cref{chap:bj-cartier-realization} we abbreviate by \emph{PSH2} the
decreasing-limit axiom, \Cref{lem:ch39-closure-any-class}\,(4): if
$\theta_j\to\theta$ and $\psi_j\in\PSH_X(\theta_j)$ is a decreasing net
whose pointwise limit is generically finite, that limit is
$\theta$-psh.  The fixed-class case, \Cref{thm:bj-basic-closure}\,(4), is
its specialization $\theta_j=\theta$.  Its two hypotheses --- a controlled class and a
generically finite limit --- are verified explicitly at every application.

\subsection{Dependencies}
\label{sec:ch39-dependencies}

From \Cref{chap:trivial-analytification}: the support and centre maps
\Cref{thm:ch32-support-center-maps}, the partition by support
\Cref{prop:ch32-support-partition}, the scaling action
\Cref{prop:ch32-scaling-action}, the divisorial points
\Cref{def:ch32-divisorial}, and the maximum principle
\Cref{prop:ch32-maximum-principle}.

From \Cref{chap:bj-fs-pl-psh}: the definition \Cref{def:bj-psh}, the
structure of a Fubini--Study presentation
\Cref{prop:ch34-generic-finiteness}, the Veronese normalization
\Cref{lem:ch34-veronese}, the existence thresholds
\Cref{lem:ch34-psh-nonempty}, and the closure calculus
\Cref{thm:bj-basic-closure}.

From \Cref{chap:bj-homogeneity-bdivisors}: the scaling calculus
\Cref{lem:ch35-action-basic}, the normalization
\Cref{prop:ch35-homogeneous-normalization}, the comparison theorem
\Cref{thm:ch35-divisorial-comparison} with its corollary
\Cref{cor:ch35-divisorial-determination}, the Legendre decomposition
\Cref{thm:homogeneous-legendre-decomposition} with the slice trichotomy
\Cref{prop:ch35-slice-trichotomy}, the carrier tail estimate
\Cref{lem:ch35-carrier-tail}, and the homogeneous Siu decomposition
\Cref{thm:ch35-homogeneous-siu} and
\Cref{thm:ch35-homogeneous-siu-general}.

From \Cref{chap:bj-cartier-realization}: the evaluation pairing
\Cref{prop:ch38-evaluation-welldefined} and its functoriality
\Cref{prop:ch38-evaluation-functorial}, the notion of nef-total class
\Cref{def:bj-nef-total}, the activation lemma \Cref{lem:ch38-activation},
the realization theorem \Cref{thm:bj-real-cartier-realization}, the
approximation package \Cref{cor:ch38-gf-regularization-bdivisor}, and the
partial pullback statement \Cref{cor:ch38-realization-pullback} whose gap
this chapter closes (\Cref{rem:ch38-pullback-forward}).

From \Cref{chap:bj-envelope-history} and
\Cref{chap:bj-envelope-statement}: the two definitions
\Cref{def:bj-envelope-property} and \Cref{def:bj-obstacle-envelope}, the
notion of a bounded-above family \Cref{def:bj-bounded-family}, and the
equivalence \Cref{prop:bj-increasing-net-criterion}.  Nothing else from
\Cref{chap:bj-envelope-statement} is used, and \Cref{thm:bj-final-envelope}
is used nowhere.

From the appendices: the extended-real ledger of
\Cref{sec:extended-arithmetic} and \Cref{lem:extended-real-safe}; the net
vocabulary of \Cref{app:nets}, in particular the Willard subnet
\Cref{def:appE-willard-subnet}; the upper-semicontinuous calculus
\Cref{prop:appE-usc-majorant} and \Cref{lem:appE-usc-calculus}; and the
birational input \Cref{thm:appB-admissible-domination} used inside
\Cref{lem:ch38-activation}.

\section{The birational dictionary between the two analytifications}
\label{sec:ch39-analytifications}

\subsection{The map on analytifications}
\label{sec:ch39-beta-an}

\begin{lemma}
\label{lem:ch39-beta-an}
Let $\beta:Y\to X$ be as in \eqref{eq:ch39-standing}.  Then the induced
map $\beta^{\An}:Y^{\An}\to X^{\An}$, $w\mapsto w\circ\beta^{\#}$,
satisfies:
\begin{enumerate}
\item\label{itm:ch39-ban:cont} $\beta^{\An}$ is continuous and
 surjective, and $Y^{\An}$, $X^{\An}$ being compact Hausdorff, it is a
 closed map and a topological quotient map;
\item\label{itm:ch39-ban:supp} $s(\beta^{\An}(w))=\overline{\beta(s(w))}$
 for every $w\in Y^{\An}$; in particular $\beta^{\An}$ maps points of
 full support to points of full support and $\beta^{\An}(v_{Y,\triv})
 =v_{X,\triv}$;
\item\label{itm:ch39-ban:val} $\beta^{\An}$ restricts to a bijection
 $Y^{\val}\to X^{\val}$ and to the identity on
 $Y^{\divis}=X^{\divis}$, under the identification $k(Y)=k(X)$;
\item\label{itm:ch39-ban:scaling} $\beta^{\An}(tw)=t\,\beta^{\An}(w)$ for
 every $t>0$, so $\beta^{\An}$ is equivariant for the scaling action of
 \eqref{eq:valuation-scaling};
\item\label{itm:ch39-ban:sections} for every nonzero
 $s\in H^0(X,mL)$, $\log|\beta^*s|=\beta^*\log|s|$ as functions
 $Y^{\An}\to[-\infty,0]$, and for every nonzero coherent ideal
 $\mathfrak a\subseteq\cO_X$,
 $\log|\mathfrak a\cdot\cO_Y|=\beta^*\log|\mathfrak a|$.
\end{enumerate}
\end{lemma}

\begin{proof}
We prove \cref{itm:ch39-ban:supp}--\cref{itm:ch39-ban:val} first and deduce
\cref{itm:ch39-ban:cont} from them.

\cref{itm:ch39-ban:supp}  By definition $f\in\mathfrak p_{\beta^{\An}(w)}$
if and only if $w(\beta^{\#}f)=+\infty$, that is,
$\beta^{\#}f\in\mathfrak p_w$, that is, $f$ vanishes on $s(w)$, that
is, $f$ vanishes on $\overline{\beta(s(w))}$.  If $s(w)=Y$, then
$\overline{\beta(Y)}=X$ since $\beta$ is dominant; and
$v_{Y,\triv}$ is the trivial valuation of $k(Y)=k(X)$, which pushes
forward to the trivial valuation of $k(X)$.

\cref{itm:ch39-ban:val}  By \cref{itm:ch39-ban:supp} the map sends
$Y^{\val}$ into $X^{\val}$.  Both sets are, by
\eqref{eq:triv-full-support-locus}, the set of real valuations of the common
field $k(X)=k(Y)$, trivial on $k$, admitting a centre; and a valuation
of $k(X)$ has a centre on $X$ if and only if it has one on $Y$,
because $X$ and $Y$ are both proper over $k$ and the valuative
criterion of properness \Cref{prop:ch32-valuation-criterion} applies to
each.  Under this identification $\beta^{\An}$ is the identity.  A
divisorial point of $Y$ is a positive multiple of $\ord_E$ for a prime
divisor $E$ on some model over $Y$, and models over $Y$ are cofinal
among models over $X$ by
\Cref{prop:ch38-evaluation-functorial}\,\cref{itm:ch38-fun:cofinal}; hence
$Y^{\divis}=X^{\divis}$ as subsets of the set of valuations of $k(X)$.

\cref{itm:ch39-ban:cont}  Continuity of the analytification of a morphism is
part of the construction \Cref{def:trivial-analytification}, and
$X^{\An}$, $Y^{\An}$ are compact Hausdorff by
\Cref{thm:triv-berkovich-compactness}.  Hence
$\beta^{\An}(Y^{\An})$ is compact, therefore closed in $X^{\An}$; by
\cref{itm:ch39-ban:val} it contains $X^{\divis}$, which is dense in
$X^{\An}$ by \Cref{thm:triv-divisorial-density}.  Therefore
$\beta^{\An}$ is surjective.  A continuous map between compact Hausdorff
spaces is closed, and a surjective closed continuous map is a quotient map.

\cref{itm:ch39-ban:scaling}  $(tw)\circ\beta^{\#}=t\,(w\circ\beta^{\#})$
because scaling acts by multiplying values, \eqref{eq:valuation-scaling}.

\cref{itm:ch39-ban:sections}  For $w\in Y^{\An}$,
$\log|\beta^*s|(w)=-w(\beta^*s)=-(\beta^{\An}w)(s)=\log|s|(\beta^{\An}w)$,
by \eqref{eq:section-log-value} and the definition of $\beta^{\An}$; the
computation is valid in $[-\infty,0]$ with no subtraction.  The ideal
statement follows by applying this to a finite set of local generators and
using \eqref{eq:ch34-log-ideal}: the ideal generated in $\cO_Y$ by the
pullbacks of generators of $\mathfrak a$ is $\mathfrak a\cdot\cO_Y$.
\end{proof}

\begin{remark}[What surjectivity buys, and what it does not]
\label{rem:ch39-surjectivity}
Surjectivity of $\beta^{\An}$ is used constantly below, in exactly two
ways.  It makes $\beta^*$ injective on functions, so that an identity
$\beta^*\chi_1=\beta^*\chi_2$ forces $\chi_1=\chi_2$, and it makes
$\beta^*$ reflect inequalities: $\beta^*\chi_1\le\beta^*\chi_2$ implies
$\chi_1\le\chi_2$ pointwise.  Both are immediate and both are valid for
arbitrary extended-real functions, with no positivity or continuity
hypothesis.  What surjectivity does \emph{not} give is a way to define
$\chi$ from $\psi$; for that one needs $\psi$ to be constant on fibres,
which is the content of \Cref{cor:ch39-fibre-constant} and is available only
after the reflection theorem has been proved.
\end{remark}

\subsection{Fibres, and where they are large}
\label{sec:ch39-fibres}

\begin{proposition}
\label{prop:ch39-fibres}
Let $\beta:Y\to X$ be as in \eqref{eq:ch39-standing}.
\begin{enumerate}
\item\label{itm:ch39-fib:val} If $x\in X^{\val}$, that is, if $s(x)=X$,
 then $(\beta^{\An})^{-1}(x)$ is a single point.
\item\label{itm:ch39-fib:triv} For every closed subscheme $Z\subseteq X$,
\[
 (\beta^{\An})^{-1}\bigl(\{x\in X^{\An}:s(x)\subseteq Z\}\bigr)
 =\{w\in Y^{\An}\ :\ s(w)\subseteq\beta^{-1}(Z)\}
 =\bigl(\beta^{-1}(Z)\bigr)^{\An},
\tag{39.3}\label{eq:ch39-fibre-trivial}
\]
 a closed subset of $Y^{\An}$.
\item\label{itm:ch39-fib:point} If $p\in X$ is a closed point, with
 set-theoretic fibre $F_p:=\beta^{-1}(p)$, then
\[
 (\beta^{\An})^{-1}(v_{p,\triv})=F_p^{\An},
\]
 which is a nonempty closed subset of $Y^{\An}$, all of whose points have
 proper support, and which is infinite as soon as $\dim F_p>0$.
\end{enumerate}
\end{proposition}

\begin{proof}
\cref{itm:ch39-fib:val}  By
\Cref{lem:ch39-beta-an}\,\cref{itm:ch39-ban:val} the map $\beta^{\An}$
restricts to a bijection $Y^{\val}\to X^{\val}$, so it suffices to show
that a point $w\in Y^{\An}$ with $\beta^{\An}(w)\in X^{\val}$ has full
support.  By \Cref{lem:ch39-beta-an}\,\cref{itm:ch39-ban:supp} we have
$\overline{\beta(s(w))}=X$.  Now $\beta$ is proper, so
$\beta(s(w))$ is closed and hence equals $X$; since $\beta$ is
birational, $\dim s(w)\ge\dim X=\dim Y$, and $s(w)$ is an integral
closed subscheme of the integral scheme $Y$, so $s(w)=Y$.

\cref{itm:ch39-fib:triv}  For $w\in Y^{\An}$ and $x=\beta^{\An}(w)$,
\Cref{lem:ch39-beta-an}\,\cref{itm:ch39-ban:supp} gives
$s(x)=\overline{\beta(s(w))}$, which is contained in $Z$ if and only if
$\beta(s(w))\subseteq Z$, that is, if and only if
$s(w)\subseteq\beta^{-1}(Z)$.  The second equality in
\eqref{eq:ch39-fibre-trivial} and the closedness are
\Cref{prop:ch32-support-partition}\,\cref{itm:ch32-sp:closed}.

\cref{itm:ch39-fib:point}  Take $Z=\{p\}$ in \cref{itm:ch39-fib:triv}.  The
only point of $X^{\An}$ with support $\{p\}$ is $v_{p,\triv}$, by
\Cref{prop:ch32-support-partition}\,\cref{itm:ch32-sp:triv}, so the left-hand
side of \eqref{eq:ch39-fibre-trivial} is exactly
$(\beta^{\An})^{-1}(v_{p,\triv})$, and the right-hand side is
$F_p^{\An}$.  It is nonempty by surjectivity of $\beta^{\An}$,
\Cref{lem:ch39-beta-an}\,\cref{itm:ch39-ban:cont}.  Every point $w$ of it
has $s(w)\subseteq F_p\subsetneq Y$, hence proper support in the sense of
\Cref{def:ch32-proper-support}.  If $\dim F_p>0$, then $F_p$ has
infinitely many closed points $q$, and the points $v_{q,\triv}\in
F_p^{\An}$ are pairwise distinct by
\Cref{prop:ch32-support-partition}.
\end{proof}

\begin{example}[The fibre over a blown-up point of the plane]
\label{ex:ch39-fibre-plane}
Let $X=\mathbf P^2_k$ with homogeneous coordinates $x,y,z$, let
$p=[0:0:1]$, and let $\beta:Y=\operatorname{Bl}_pX\to X$ be the blowup
with exceptional prime $E\cong\mathbf P^1$.  By
\Cref{prop:ch39-fibres}\,\cref{itm:ch39-fib:point},
\[
 (\beta^{\An})^{-1}(v_{p,\triv})=E^{\An}\cong(\mathbf P^1_k)^{\An},
\]
which by \Cref{thm:ch32-p1-classification} is the cone over the set of
closed points of $E$, with vertex $v_{E,\triv}$ and one ray
$\{t\,\ord_q:t\ge0\}$ for each closed point $q\in E$ --- here
$\ord_q$ is computed inside $E$, so these are \emph{not} points of
$Y^{\divis}$; they have support contained in $E$, hence proper support
in $Y$.  Every point of this fibre has proper support.  Concretely, the
fibre consists of $v_{E,\triv}$, of $v_{q,\triv}$ for $q\in E$ a
closed point, and of the interpolating semivaluations along each ray.

Meanwhile the divisorial point $\ord_E\in Y^{\divis}$ does \emph{not}
lie in this fibre: it has full support on $Y$, and
$\beta^{\An}(\ord_E)=\ord_E\in X^{\divis}$, the divisorial valuation of
$k(X)$ with centre $p$.  This is the distinction that
\Cref{lem:ch39-beta-an}\,\cref{itm:ch39-ban:val} and
\Cref{prop:ch39-fibres}\,\cref{itm:ch39-fib:val} record: the exceptional
divisor of $\beta$ contributes a \emph{point of $X^{\divis}$}, while
the exceptional fibre of $\beta^{\An}$ sits entirely in the
proper-support locus.  Confusing the two is the commonest source of error
in birational arguments over a trivially valued field.
\end{example}

\subsection{Closure properties in a class that need not be nef}
\label{sec:ch39-closure}

\Cref{thm:bj-basic-closure} already covers arbitrary numerical classes, and
by \Cref{lem:ch34-psh-nonempty} its assertions are vacuous unless the
relevant classes are pseudoeffective.  What it does not contain is the
monotonicity of $\PSH_X(\theta)$ in $\theta$.  That extra statement is
needed here since a homogeneous psh function can exist in a class that is
pseudoeffective but not nef, and since \Cref{chap:bj-alterations} applies
\Cref{thm:bj-birational-ep-descent} to an arbitrary $\theta$.  We first
prove the elementary monotonicity of the Fubini--Study spaces and then
package it with the closure theorem in the form used below.

\begin{lemma}
\label{lem:ch39-hgf-monotone}
Let $L,M\in\Pic(X)_{\mathbf Q}$ be such that $M-L$ is semiample, that
is, such that $r(M-L)$ is a globally generated line bundle for some
$r\in\mathbf Z_{>0}$.  Then, for every additive subgroup
$\Lambda\subseteq\mathbf R$,
\[
 \cH^{\mathrm{gf}}_\Lambda(L)\subseteq\cH^{\mathrm{gf}}_\Lambda(M),
 \qquad
 \cH_\Lambda(L)\subseteq\cH_\Lambda(M).
\tag{39.4}\label{eq:ch39-hgf-monotone}
\]
\end{lemma}

\begin{proof}
Let $\phi=\frac1m\max_j\{\log|s_j|+\lambda_j\}$ with
$0\ne s_j\in H^0(X,mL)$, as in \eqref{eq:gf-fs-presentation}.  Replacing
$m$ by a multiple by \Cref{lem:ch34-veronese}, which changes neither
$\phi$ nor the weights, we may assume that $r\mid m$, that $mL$ and
$mM$ are honest line bundles, and that $m(M-L)$ is globally generated.
Choose global sections $u_1,\dots,u_N$ generating $m(M-L)$.  Then
$s_ju_k\in H^0(X,mM)$ and $s_ju_k\ne0$, the sheaf $m(M-L)$ being a
line bundle on the integral scheme $X$ and $u_k\ne0$.  Since the
$u_k$ generate, the ideal they define in any local frame is
$\cO_X$, so by \eqref{eq:ch34-log-ideal}
\[
 \max_k\log|u_k|=\log|\cO_X|=0
 \qquad\text{on }X^{\An}.
\]
Therefore, using that a finite maximum in $\Rbar$ distributes over the
sum with $\log|s_j|$ --- an addition of elements of $[-\infty,0]$, never
a subtraction ---
\[
 \frac1m\max_{j,k}\{\log|s_ju_k|+\lambda_j\}
 =\frac1m\max_j\Bigl\{\log|s_j|+\lambda_j+\max_k\log|u_k|\Bigr\}
 =\phi .
\]
The system $(m,(s_ju_k),(\lambda_j))$ is a weighted section system for
$M$ with weights in $\Lambda$, and it is generically finite because the
$s_ju_k$ do not all vanish identically: some $s_j$ does not, and
$u_k\ne0$ while $X$ is integral.  Hence
$\phi\in\cH^{\mathrm{gf}}_\Lambda(M)$.  If moreover the $s_j$ have no
common zero, then neither do the $s_ju_k$: at any point some $s_j$ is
nonzero, and some $u_k$ is nonzero since the $u_k$ generate.
This gives the second inclusion.
\end{proof}

\begin{warning}{The inclusion is not induced by an inclusion of classes}
\label{warn:ch39-hgf-monotone}
\Cref{lem:ch39-hgf-monotone} requires $M-L$ \emph{semiample}, not merely
nef, because its proof needs sections, and it is a statement about a fixed
function belonging to two different Fubini--Study spaces --- not about
comparing the two spaces as subsets of a common one indexed by classes.
The corresponding statement at the level of $\PSH$, namely
$\PSH_X(\theta)\subseteq\PSH_X(\theta+\eta)$ for $\eta$ nef, is
\Cref{lem:ch39-closure-any-class}\,(5) below, and it does hold for merely
nef $\eta$, because there the semiampleness is recovered by an
arbitrarily small ample perturbation which is then absorbed in the limit of
classes.
\end{warning}

\begin{lemma}
\label{lem:ch39-closure-any-class}
Let $X$ be a projective $k$-variety and let
$\theta,\theta'\in N^1(X)_{\mathbf R}$ be arbitrary.
\begin{enumerate}
\item If $\phi\in\PSH_X(\theta)$ and $c\in\mathbf R$, then
 $\phi+c\in\PSH_X(\theta)$.
\item If $\phi,\psi\in\PSH_X(\theta)$, then
 $\max\{\phi,\psi\}\in\PSH_X(\theta)$.
\item If $\phi\in\PSH_X(\theta)$, $\psi\in\PSH_X(\theta')$ and
 $s,t\ge0$ are not both $0$, then
  $s\phi+t\psi\in\PSH_X(s\theta+t\theta')$, the sum being formed in
 $\Rbar$; a summand whose coefficient is $0$ is omitted, and is never
 interpreted as $0\cdot(-\infty)$.
\item If $\theta_j\to\theta$ in $N^1(X)_{\mathbf R}$, if
 $\psi_j\in\PSH_X(\theta_j)$ is a decreasing net, and if
 $\psi:=\lim_j\psi_j$ is generically finite, then
 $\psi\in\PSH_X(\theta)$.  In particular this holds for a constant net of
 classes, and the set $\{\theta:\psi\in\PSH_X(\theta)\}$ is closed for
 every fixed $\psi$.
\item If $\phi\in\PSH_X(\theta)$ and $\eta\in\Nef(X)$, then
 $\phi\in\PSH_X(\theta+\eta)$.
\end{enumerate}
\end{lemma}

\begin{proof}
Clauses (1)--(4) are precisely
\Cref{thm:bj-basic-closure}\,(1)--(5), with its fixed-class and
varying-class decreasing-limit clauses combined in (4).  That theorem is
already stated for arbitrary numerical classes, and its proof records the
generic-finiteness hypothesis needed by PSH2.  Thus only (5) is new here.

(5)  Let $(\phi_i)$ approximate $\phi$ in classes $\alpha_i\to\theta$.
Fix a very ample $A$.  For $n\ge1$ choose a
$\mathbf Q$-line bundle $N_n$ whose class is ample and satisfies
$\|c_1(N_n)-\eta\|\le1/n$.  This is possible because
$\eta+\varepsilon c_1(A)$ is ample for every $\varepsilon>0$, these classes
converge to $\eta$ as $\varepsilon\downarrow0$, and rational classes are
dense in the open ample cone and are represented by $\mathbf Q$-line
bundles.
Since $N_n$ is ample it is semiample, so
\Cref{lem:ch39-hgf-monotone} gives
$\phi_i\in\cH^{\mathrm{gf}}_{\mathbf R}(L_i+N_n)$.  Indexing by the pairs
$(i,n)\in I\times\mathbf Z_{>0}$ with the product order, we obtain a
decreasing net --- decreasing in $i$ and constant in $n$ --- of
generically finite Fubini--Study functions with classes
$\alpha_i+c_1(N_n)\to\theta+\eta$ and limit $\phi$.  Hence
$\phi\in\PSH_X(\theta+\eta)$.
\end{proof}

\begin{remark}[What was and was not generalized]
\label{rem:ch39-closure-scope}
\Cref{lem:ch39-closure-any-class} restates clauses (1)--(5) of
\Cref{thm:bj-basic-closure} --- which are already proved there for arbitrary
numerical classes, that theorem carrying no positivity hypothesis --- merging
its (4) and~(5) into a single clause~(4), and adds the new clause~(5), the
monotonicity of $\PSH_X$ in the class.
It does \emph{not} generalize anything about suprema: there is still no
statement about increasing families, and \Cref{warn:ch39-firewall} remains
in force.  Nor does it assert that $\PSH_X(\theta)$ is nonempty; by
\Cref{lem:ch34-psh-nonempty} that requires $\theta$ pseudoeffective, and
in the pseudoeffective boundary case it may still fail to contain any
bounded function.  Every application below of ``PSH2'' is an application of
clause (4).
\end{remark}

\subsection{Pullback of psh functions}
\label{sec:ch39-pullback}

\begin{theorem}
\label{thm:ch39-pullback}
Let $\beta:Y\to X$ be as in \eqref{eq:ch39-standing} and let
$\theta\in N^1(X)_{\mathbf R}$.  Write
$\beta^*u:=u\circ\beta^{\An}$.  Then:
\begin{enumerate}
\item\label{itm:ch39-pb:fs} $\beta^*$ maps
 $\cH^{\mathrm{gf}}_\Lambda(L)$ into
 $\cH^{\mathrm{gf}}_\Lambda(\beta^*L)$ and $\cH_\Lambda(L)$ into
 $\cH_\Lambda(\beta^*L)$, for every $\mathbf Q$-line bundle $L$ on
 $X$ and every additive subgroup $\Lambda\subseteq\mathbf R$;
\item\label{itm:ch39-pb:psh} $\beta^*$ maps $\PSH_X(\theta)$ into
 $\PSH_Y(\beta^*\theta)$ and
 $\PSH_{\mathrm{hom},X}(\theta)$ into
 $\PSH_{\mathrm{hom},Y}(\beta^*\theta)$;
\item\label{itm:ch39-pb:order} $\beta^*$ is injective and reflects the
 order: for arbitrary functions $u_1,u_2$ on $X^{\An}$,
 $\beta^*u_1\le\beta^*u_2$ if and only if $u_1\le u_2$;
\item\label{itm:ch39-pb:ops} $\beta^*$ commutes with the addition of
 real constants and with arbitrary suprema and infima, preserves upper
 semicontinuity, and commutes with the scaling action:
 $\beta^*(t\cdot u)=t\cdot\beta^*u$ for $t>0$.
 In particular $\beta^*$ commutes with the Legendre slices,
 $\beta^*(\widehat u_\lambda)=\widehat{(\beta^*u)}_\lambda$;
\item\label{itm:ch39-pb:sup} $\sup_{Y^{\An}}\beta^*u=\sup_{X^{\An}}u$ and
 $\beta^*u(v_{Y,\triv})=u(v_{X,\triv})$.
\end{enumerate}
\end{theorem}

\begin{proof}
\cref{itm:ch39-pb:fs}  Let
$\phi=\frac1m\max_j\{\log|s_j|+\lambda_j\}$ with
$0\ne s_j\in H^0(X,mL)$.  Then $\beta^*s_j\in H^0(Y,m\beta^*L)$ is
nonzero, because $\beta$ is dominant and $Y$ is integral, and by
\Cref{lem:ch39-beta-an}\,\cref{itm:ch39-ban:sections},
\[
 \beta^*\phi=\frac1m\max_j\{\log|\beta^*s_j|+\lambda_j\},
\]
a weighted section system for $\beta^*L$ with the same weights.  Generic
finiteness holds because $Y$ is integral and some $\beta^*s_j$ is
nonzero.  If the $s_j$ have no common zero on $X$, then the
$\beta^*s_j$ have no common zero on $Y$, because
$\beta$ is surjective.

\cref{itm:ch39-pb:psh}  Let $u\in\PSH_X(\theta)$ with approximating net
$u_i\in\cH^{\mathrm{gf}}_{\mathbf R}(L_i)$, $c_1(L_i)\to\theta$,
$u_i\searrow u$.  By \cref{itm:ch39-pb:fs},
$\beta^*u_i\in\cH^{\mathrm{gf}}_{\mathbf R}(\beta^*L_i)$.  The net
$(\beta^*u_i)$ is decreasing, being the composition of a decreasing net
with a fixed map, and it converges pointwise to $\beta^*u$.  Its classes
satisfy $\beta^*c_1(L_i)\to\beta^*\theta$, the map
$\beta^*:N^1(X)_{\mathbf R}\to N^1(Y)_{\mathbf R}$ being linear on a
finite-dimensional space, hence continuous.  Finally $\beta^*u$ is
generically finite: $Y$ is integral, and
$\beta^*u(v_{Y,\triv})=u(v_{X,\triv})>-\infty$ by
\Cref{lem:ch39-beta-an}\,\cref{itm:ch39-ban:supp} and
\Cref{prop:ch32-maximum-principle}\,\cref{itm:ch32-mp:gf}.  Hence
\Cref{def:bj-psh} applies to $\beta^*u$ with the class
$\beta^*\theta$.  Homogeneity is preserved by
\cref{itm:ch39-pb:ops}.

\cref{itm:ch39-pb:order}  Immediate from surjectivity of $\beta^{\An}$,
\Cref{lem:ch39-beta-an}\,\cref{itm:ch39-ban:cont}: for $x\in X^{\An}$
choose $w$ with $\beta^{\An}(w)=x$ and evaluate.

\cref{itm:ch39-pb:ops}  Composition with a fixed map commutes with all
pointwise lattice operations and with the addition of constants.  For the
scaling action, \Cref{lem:ch39-beta-an}\,\cref{itm:ch39-ban:scaling} gives
\[
 \beta^*(t\cdot u)(w)=(t\cdot u)(\beta^{\An}w)
 =t\,u\bigl(t^{-1}\beta^{\An}w\bigr)
 =t\,u\bigl(\beta^{\An}(t^{-1}w)\bigr)
 =(t\cdot\beta^*u)(w),
\]
and the statement about slices follows from
\eqref{eq:homogeneous-legendre-slice}, an infimum over $t>0$ of the
functions $t\cdot u-t\lambda$, together with the commutation of
$\beta^*$ with infima and with the subtraction of the finite constant
$t\lambda$.  Upper semicontinuity is preserved because $\beta^{\An}$ is
continuous: $\{\beta^*u\ge c\}=(\beta^{\An})^{-1}\{u\ge c\}$ is closed
whenever $\{u\ge c\}$ is.

\cref{itm:ch39-pb:sup}  The first identity is surjectivity; the second is
\Cref{lem:ch39-beta-an}\,\cref{itm:ch39-ban:supp}.
\end{proof}

\begin{warning}{Pullback is injective but not surjective for free}
\label{warn:ch39-pullback-not-onto}
\Cref{thm:ch39-pullback} gives an injection
$\beta^*:\PSH_X(\theta)\hookrightarrow\PSH_Y(\beta^*\theta)$ and nothing
more.  Its image is contained in the set of functions constant on the fibres
of $\beta^{\An}$, and by \Cref{prop:ch39-fibres} those fibres are large:
over the trivial semivaluation of a blown-up point they are the
analytification of the whole exceptional locus.  Constancy along such a
fibre is a nontrivial condition, and it is \emph{not} implied by
membership in $\PSH_Y(\eta)$ for a general $\eta$; see
\Cref{ex:ch39-not-fibre-constant}.  That it \emph{is} implied when
$\eta=\beta^*\theta$ and the function is homogeneous is the content of
\Cref{thm:bj-homogeneous-reflection}, and the proof of that theorem does not
proceed by verifying fibrewise constancy: it constructs the function
downstairs.
\end{warning}

\subsection{Numerical classes and total classes}
\label{sec:ch39-classes}

\begin{lemma}
\label{lem:ch39-total-class}
Let $B$ be a real Cartier b-divisor over $X$, determined on a model
$\pi:Z\to X$ that factors as $Z\xrightarrow{\ \rho\ }Y
\xrightarrow{\ \beta\ }X$.  Then, for every
$\theta\in N^1(X)_{\mathbf R}$, the total class of $B$ in the class
$\theta$ computed over $X$ and the total class of $B$ in the class
$\beta^*\theta$ computed over $Y$ coincide:
\[
 \pi^*\theta+B_Z=\rho^*(\beta^*\theta)+B_Z\in N^1(Z)_{\mathbf R}.
\tag{39.5}\label{eq:ch39-total-class}
\]
Consequently $B$ is nef-total in the class $\theta$ over $X$ if and
only if it is nef-total in the class $\beta^*\theta$ over $Y$, and
$B\le0$ over $X$ if and only if $B\le0$ over $Y$.
\end{lemma}

\begin{proof}
The identity $\pi=\beta\rho$ gives $\pi^*\theta=\rho^*\beta^*\theta$,
which is \eqref{eq:ch39-total-class}.  Nef-totality is by
\Cref{def:bj-nef-total} a property of the class
\eqref{eq:ch39-total-class} on a determination, and by
\Cref{lem:ch38-nef-total-intrinsic}\,\cref{itm:ch38-nt:indep} it does not
depend on which determination is used; models over $Y$ are cofinal among
models over $X$ by
\Cref{prop:ch38-evaluation-functorial}\,\cref{itm:ch38-fun:cofinal}.
Antieffectivity is the condition $B_W\le0$ for all models $W$, and the
two model categories have the same cofinal part, so the two conditions
agree; this is
\Cref{lem:ch38-determination-equivalence}\,\cref{itm:ch38-de:order}.
\end{proof}

\begin{proposition}
\label{prop:ch39-cartier-reflection}
Let $B\le0$ be a real Cartier b-divisor over $X$, determined on a model
factoring through $\beta$, and suppose that its total class in $\theta$
is nef.  Then the realizations of \Cref{thm:bj-real-cartier-realization}
over $X$ and over $Y$ exist and satisfy
\[
 \beta^*\chi^X_B=\chi^Y_B
 \qquad\text{on }Y^{\An}.
\tag{39.6}\label{eq:ch39-cartier-reflection}
\]
\end{proposition}

\begin{proof}
By \Cref{lem:ch39-total-class} the hypotheses of
\Cref{thm:bj-real-cartier-realization} hold over $X$ in the class
$\theta$ and over $Y$ in the class $\beta^*\theta$; let $\chi^X_B$
and $\chi^Y_B$ be the two realizations.  By
\Cref{thm:ch39-pullback}\,\cref{itm:ch39-pb:psh},
$\beta^*\chi^X_B\in\PSH_{\mathrm{hom},Y}(\beta^*\theta)$.  Both
$\beta^*\chi^X_B$ and $\chi^Y_B$ restrict to $\psi_B$ on
$Y^{\divis}=X^{\divis}$: the second by the theorem, the first because
$\beta^{\An}$ is the identity on divisorial points,
\Cref{lem:ch39-beta-an}\,\cref{itm:ch39-ban:val}, and
$\psi_B$ is the same function computed over $X$ or over $Y$ by
\Cref{prop:ch38-evaluation-functorial}\,\cref{itm:ch38-fun:psi}.  Both are
$\beta^*\theta$-psh, so
\Cref{cor:ch35-divisorial-determination}\,(2) gives
\eqref{eq:ch39-cartier-reflection}.
\end{proof}

\begin{remark}[Why \Cref{prop:ch39-cartier-reflection} is not the theorem]
\label{rem:ch39-cartier-not-enough}
\Cref{prop:ch39-cartier-reflection} settles reflection for the functions
$\psi_B$ attached to Cartier b-divisors, and it does so by a two-line
argument: the numerical hypothesis is literally the same upstairs and
downstairs, by \Cref{lem:ch39-total-class}, and divisorial determination
does the rest.  It even settles surjectivity of $\beta^*$ on the subset of
$\PSH_{\mathrm{hom},Y}(\beta^*\theta)$ consisting of realizations of
Cartier b-divisors.  This subset is however \emph{not} all of
$\PSH_{\mathrm{hom},Y}(\beta^*\theta)$: a general homogeneous psh function
corresponds to a Weil b-divisor which is Cartier on no single model, by
\Cref{rem:ch38-cartier-not-weil}, and is only a decreasing limit of Cartier
data.  The difficulty is that the approximating Cartier data on $Y$ come
with classes $\alpha_i\in N^1(Y)$ converging to $\beta^*\theta$ but not
lying in $\beta^*N^1(X)$, so \Cref{lem:ch39-total-class} does not apply to
the individual stages.  Removing that discrepancy, in a way compatible with
the order, is the entire content of \Cref{sec:ch39-homogeneous}.
\end{remark}

\section{Homogeneous reflection}
\label{sec:ch39-homogeneous}

\subsection{Reduction to a dominating model}
\label{sec:ch39-model-replacement}

The fixed-pole construction below produces its data only after replacing
$Y$ by a model dominating it.  The following observation makes that
replacement free of charge, and it is the first place where
\Cref{thm:ch39-pullback} earns its keep.

\begin{lemma}
\label{lem:ch39-model-replacement}
Let $\rho:Y''\to Y$ be a projective birational morphism with $Y''$
normal, and put $\beta'':=\beta\rho:Y''\to X$.  If
\[
 (\beta'')^*:\PSH_{\mathrm{hom},X}(\theta)
 \longrightarrow\PSH_{\mathrm{hom},Y''}\bigl((\beta'')^*\theta\bigr)
\]
is surjective, then so is
$\beta^*:\PSH_{\mathrm{hom},X}(\theta)\to
\PSH_{\mathrm{hom},Y}(\beta^*\theta)$.  The same implication holds with
$\PSH_{\mathrm{hom}}$ replaced by $\PSH$ throughout.
\end{lemma}

\begin{proof}
Let $\psi\in\PSH_{\mathrm{hom},Y}(\beta^*\theta)$.  By
\Cref{thm:ch39-pullback}\,\cref{itm:ch39-pb:psh} applied to $\rho$,
$\rho^*\psi\in\PSH_{\mathrm{hom},Y''}(\rho^*\beta^*\theta)
=\PSH_{\mathrm{hom},Y''}((\beta'')^*\theta)$.  By hypothesis there is
$\chi\in\PSH_{\mathrm{hom},X}(\theta)$ with
$(\beta'')^*\chi=\rho^*\psi$, that is,
$\rho^*(\beta^*\chi)=\rho^*\psi$.  Since $\rho^{\An}$ is surjective,
\Cref{thm:ch39-pullback}\,\cref{itm:ch39-pb:order} gives
$\beta^*\chi=\psi$.  The inhomogeneous case is verbatim the same.
\end{proof}

\subsection{The fixed-pole data over a birational base}
\label{sec:ch39-fixed-pole}

\begin{lemma}
\label{lem:ch39-fixed-pole}
Assume $\dim X>0$.  There exist a projective birational morphism
$\rho:Y''\to Y$ with $Y''$ normal, an effective Cartier divisor $F$ on
$X$ and an effective $\mathbf Q$-Cartier divisor $G$ on $Y''$ such
that, writing $\beta'':=\beta\rho$ and $\omega:=\theta+[F]$,
\[
 \begin{gathered}
 [F]\ \text{and}\ \omega\ \text{are ample},
 \qquad
 Q:=(\beta'')^*\omega-G\ \text{is ample},\\
 (\beta'')^*F=G+R
 \quad\text{with}\quad
 R\ge\tfrac12(\beta'')^*F\ge0 .
 \end{gathered}
\tag{39.7}\label{eq:bj-reflection-fixed-pole}
\]
No positivity hypothesis on $\theta$ is used, and no b-divisor enters the
statement.
\end{lemma}

\begin{proof}
This is the content of Steps~1--3 of the proof of
\Cref{lem:ch38-activation}, applied to the zero b-divisor $B=0$ on the
model $\beta:Y\to X$; we recall the construction because only items
\cref{itm:ch38-act:ample}, \cref{itm:ch38-act:dominated} and
\cref{itm:ch38-act:Qample} of that lemma are quoted here, and because we must
check that those three items are obtained without the nefness of the total
class --- a hypothesis which is not available to us, $\theta$ being an
arbitrary real class.  Indeed Steps~4 and~5 of that proof, which are the
only steps using the total class, establish exactly the two items
\cref{itm:ch38-act:total} and \cref{itm:ch38-act:pinef} that we do not use.

\emph{Step 1: a relatively ample exceptional divisor.}
The morphism $\beta$ is proper and is an isomorphism over the complement
of a closed subset $Z_0\subsetneq X$.  By
\Cref{thm:appB-admissible-domination} there is a blowup $X'\to X$ with
centre contained in $Z_0$ which factors through $Y$.  Let $Y''$ be the
normalization of $X'$; it is projective over $X$ because normalization
is finite, it is normal, and $\rho:Y''\to Y$ is projective birational.  If
$\mathcal I\subseteq\cO_X$ is the ideal blown up, then
$\mathcal I\cdot\cO_{Y''}=\cO_{Y''}(-G_0)$ for an effective Cartier
divisor $G_0$ with $-G_0$ relatively ample over $X$, by
\eqref{eq:appB-exceptional-ideal} and the construction
\eqref{eq:appB-normalized-blowup} of the normalized blowup; moreover
$\Supp G_0\subseteq(\beta'')^{-1}(V(\mathcal I))$ with
$V(\mathcal I)\subseteq Z_0\subsetneq X$.

\emph{Step 2: the activation divisor $F$.}
Set $Z:=Z_0\cup V(\mathcal I)$, a proper closed subset of $X$.  Let
$A$ be a very ample divisor on $X$.  Choose $m\in\mathbf Z_{>0}$ so
large that $\mathcal I_Z\otimes\cO_X(mA)$ is globally generated, which is
possible by Serre's theorem, and so large that $\theta+m\,c_1(A)$ is
ample, which is possible because $\Amp(X)$ is open and contains
$m\,c_1(A)+\theta$ for $m\gg0$.  Let $F\in|mA|$ be the divisor of
zeros of a nonzero global section of $\mathcal I_Z\otimes\cO_X(mA)$.  Then
$F$ is an effective Cartier divisor with $\Supp F\supseteq Z$, the class
$[F]=m\,c_1(A)$ is ample, and $\omega=\theta+[F]$ is ample.  This is the
first assertion of \eqref{eq:bj-reflection-fixed-pole}.

\emph{Step 3: normalizing $G$.}
By \Cref{lem:ch38-numerical-birational}\,\cref{itm:ch38-num:relative},
applied to $\beta'':Y''\to X$, to the effective Cartier divisor $G_0$
with $-G_0$ relatively ample, and to the ample class $\omega$, there is
$r_1\ge1$ with $r(\beta'')^*\omega-G_0$ ample for $r\ge r_1$; dividing
by $r$, the class $(\beta'')^*\omega-r^{-1}G_0$ is ample for
$r\ge r_1$.  Next, $\Supp G_0\subseteq(\beta'')^{-1}(V(\mathcal I))
\subseteq(\beta'')^{-1}(\Supp F)=\Supp(\beta'')^*F$, and $(\beta'')^*F$
has strictly positive coefficient along every component of its support;
since $G_0$ has finitely many components, there is $r_2\ge1$ with
$2r_2^{-1}G_0\le(\beta'')^*F$.  Put $r:=\max\{r_1,r_2\}$ and
$G:=r^{-1}G_0$, an effective $\mathbf Q$-Cartier divisor.  Then
$2G\le(\beta'')^*F$, so $R:=(\beta'')^*F-G\ge G\ge0$ and in fact
$R\ge\frac12(\beta'')^*F$, which is the second assertion; and
$Q=(\beta'')^*\omega-G$ is ample, which is the third.
\end{proof}

\begin{remark}[Each item of \eqref{eq:bj-reflection-fixed-pole}, and what
it is for]
\label{rem:ch39-fixed-pole-uses}
The ampleness of $\omega$ is what allows the whole reflected family to
live in a single space $\PSH_X(\omega)$, so that PSH2 applies with a fixed
class; without it the drift ledger of
\Cref{prop:ch38-drift-ledger} could not be oriented, and
\Cref{warn:ch38-drift-not-absorbed} explains why smallness of the drift
would not be a substitute.  The ampleness of $[F]$ is what makes
$\omega=\theta+[F]$ ample for an arbitrary, possibly very negative,
$\theta$; this is the only place where the arbitrariness of $\theta$ is
absorbed.  The inequality $(\beta'')^*F-G\ge0$ is what makes the
interpolation \eqref{eq:bj-reflection-Ei} order preserving, by
\eqref{eq:bj-reflection-Ei-order}; it is the single inequality on which the
monotonicity of the whole construction rests, and it is the reason
\Cref{lem:ch38-activation} was stated with the normalization
$2G\le\pi^*F$ rather than merely $G\ge0$.  The ampleness of $Q$ is
what allows a small multiple of the auxiliary ample class $A$ on $Y''$
to be subtracted without losing nefness, in the last line of
\eqref{eq:bj-reflection-total}.  Nothing else about $G$ is used; in
particular the relative ampleness of $-G$, which is essential inside
\Cref{thm:bj-real-cartier-realization}, plays no role in the present
chapter, because the realization theorem performs its own activation.
\end{remark}

\subsection{Orienting the class drift}
\label{sec:ch39-drift}

\begin{lemma}
\label{lem:ch39-drift-window}
Let $Y''$ be a normal projective variety, let $A$ be an ample class on
$Y''$, let $\gamma\in N^1(Y'')_{\mathbf R}$, and let
$(\alpha_i)_{i\in I}$ be a net in $N^1(Y'')_{\mathbf R}$ with
$\alpha_i\to\gamma$.  Then there are a directed set $J$, an
order-preserving cofinal map $J\to I$, $j\mapsto i(j)$, and a
nonincreasing map $J\to\mathbf Q_{>0}$, $j\mapsto\epsilon_j$, such that
\[
 \epsilon_j=\delta_j^2\ \text{with}\ \delta_j\in\mathbf Q\cap(0,1),
 \qquad
 \lim_j\delta_j=0,
 \qquad
 \gamma+\epsilon_jA-\alpha_{i(j)}\ \text{is ample for every }j.
\tag{39.8}\label{eq:ch39-drift-window}
\]
\end{lemma}

\begin{proof}
For each integer $n\ge2$ the class $n^{-2}A$ is ample, so by openness of
$\Amp(Y'')$ in the finite-dimensional space $N^1(Y'')_{\mathbf R}$ there
is $\varrho_n>0$ such that every class within distance $\varrho_n$ of
$n^{-2}A$ is ample.  Since $\alpha_i\to\gamma$, there is $i_n\in I$
with $\|\gamma-\alpha_i\|<\varrho_n$ for all $i\ge i_n$; hence
$\gamma+n^{-2}A-\alpha_i$ is ample for all $i\ge i_n$.  Put
\[
 J:=\{(i,n)\in I\times\mathbf Z_{\ge2}\ :\ i\ge i_n\},
\]
ordered by $(i,n)\le(i',n')$ if and only if $i\le i'$ and $n\le n'$.
This is directed: given $(i,n)$ and $(i',n')$, set
$n'':=\max\{n,n'\}$ and choose $i''\ge i,i',i_{n''}$, which is possible
because $I$ is directed; then $(i'',n'')\in J$ dominates both.  The
forgetful map $(i,n)\mapsto i$ is order preserving and cofinal --- given
$i\in I$, choose $n=2$ and $i''\ge i,i_2$ --- hence it defines a
Willard subnet in the sense of \Cref{def:appE-willard-subnet}.  Set
$\delta_{(i,n)}:=n^{-1}\in\mathbf Q\cap(0,1)$ and
$\epsilon_{(i,n)}:=n^{-2}=\delta_{(i,n)}^2$.  These are nonincreasing
along $J$, tend to $0$, and satisfy the ampleness in
\eqref{eq:ch39-drift-window} by construction.
\end{proof}

The class whose ampleness \Cref{lem:ch39-drift-window} arranges,
\[
                 (\beta'')^*\theta+\epsilon_jA-\alpha_j,
\tag{39.9}\label{eq:bj-reflection-class-correction}
\]
is the \emph{class correction}: it measures how much auxiliary positivity on
$Y''$ has to be spent to dominate the drift $\alpha_j\to(\beta'')^*\theta$.
Two features of \eqref{eq:bj-reflection-class-correction} are used below and
neither is a formality.  It is required to be \emph{ample}, not merely nef,
because it will be added to a semiample class and the sum must be nef for
every $j$ with no compactness argument available; and $\epsilon_j$ is
required to tend to $0$, because otherwise the auxiliary class $A$,
which lives on $Y''$ and is not pulled back from $X$, would survive into
the limit and the reflected function would have the wrong divisorial trace.

\begin{remark}[Why a square]
\label{rem:ch39-why-square}
The parameter $\delta_j=\sqrt{\epsilon_j}$ is used in
\eqref{eq:bj-reflection-Ei} to interpolate between the drifted data and the
fixed pole.  Two competing demands fix the exponent.  The interpolation
weight $\delta_j$ must tend to $0$, or else the limit b-divisor would
retain a contribution from $G$ and the reflected function would not have
the required divisorial trace.  On the other hand the correction
$t_j=(1-\delta_j)\epsilon_j/\delta_j$ subtracted from the ample class
$Q$ in \eqref{eq:bj-reflection-total} must stay bounded by the fixed
margin $c$ of $Q$; this forces $\epsilon_j=O(\delta_j)$.  Taking
$\delta_j=\sqrt{\epsilon_j}$ satisfies both with room to spare, since then
$t_j=(1-\delta_j)\delta_j\le\frac14$.  Any exponent
$\delta_j=\epsilon_j^{s}$ with $0<s<1$ would do; the square root is
chosen because it keeps $\delta_j$ rational when $\epsilon_j$ is the
square of a rational, which is what
\Cref{lem:ch39-drift-window} arranges and what
\Cref{lem:ch38-ordered-rational-dominated} needs downstream inside the
realization theorem.
\end{remark}

\subsection{The interpolated b-divisors}
\label{sec:ch39-interpolation}

We now assemble the data.  Fix, for the rest of this subsection, the
following:
\begin{itemize}
\item $\beta'':Y''\to X$, $F$, $G$, $\omega=\theta+[F]$,
 $R=(\beta'')^*F-G$ and $Q=(\beta'')^*\omega-G$ as in
 \Cref{lem:ch39-fixed-pole};
\item an ample class $A$ on $Y''$ and a rational $c>0$ with
 $Q-cA$ nef, which exists because $Q$ is ample;
\item a net of rational Cartier b-divisors $-D_i\le0$ over $Y''$ and
 $\mathbf Q$-line bundles $L_i$ on $Y''$ with
 $\alpha_i:=c_1(L_i)$, such that on a determination of $D_i$ the class
 $\pi^*\alpha_i-D_i$ is semiample, such that $D_i$ increases, and such
 that $\alpha_i\to(\beta'')^*\theta$;
\item the data $J$, $i(j)$, $\epsilon_j=\delta_j^2$ produced by
 \Cref{lem:ch39-drift-window} for $\gamma=(\beta'')^*\theta$, refined so
 that $\delta_j\le c$ for every $j$ --- possible because
 $\delta_j\to0$ and the tail of a net is cofinal.
\end{itemize}
To lighten the notation we write $D_j$ for $D_{i(j)}$ and
$\alpha_j$ for $\alpha_{i(j)}$.

\begin{lemma}
\label{lem:ch39-interpolation}
With the data above, define, on a common determination of $D_j$
dominating $Y''$,
\[
 P_j:=(\beta'')^*\theta+\epsilon_jA-D_j,
\tag{39.10}\label{eq:bj-reflection-Pi}
\]
\[
 t_j:=\frac{(1-\delta_j)\epsilon_j}{\delta_j}=(1-\delta_j)\delta_j\le c,
\tag{39.11}\label{eq:bj-reflection-delta}
\]
\[
 D'_j:=D_j+\overline F,
 \qquad
 E_j:=(1-\delta_j)D'_j+\delta_j\overline G ,
\tag{39.12}\label{eq:bj-reflection-Ei}
\]
where $\overline F$ and $\overline G$ are the Cartier b-divisors
determined by $F$ on $X$ and by $G$ on $Y''$.  Then:
\begin{enumerate}
\item\label{itm:ch39-int:nefP} $P_j$ is nef;
\item\label{itm:ch39-int:pos} $E_j\ge0$, so that $-E_j$ is an
 antieffective real Cartier b-divisor over $X$, determined on a model
 dominating $Y''$;
\item\label{itm:ch39-int:order} the family $(E_j)$ is nondecreasing:
\[
 E_{j'}-E_j
 =(1-\delta_{j'})(D'_{j'}-D'_j)
  +(\delta_j-\delta_{j'})(D'_j-\overline G)\ \ge\ 0
 \qquad(j\le j');
\tag{39.13}\label{eq:bj-reflection-Ei-order}
\]
\item\label{itm:ch39-int:total} the total class of $-E_j$ in the class
 $\omega$, computed over $X$, satisfies the exact identity
\[
 \begin{aligned}
 (\beta'')^*\omega-E_j
  &=(1-\delta_j)P_j+\delta_jQ-(1-\delta_j)\epsilon_jA\\
  &=(1-\delta_j)P_j+\delta_j(Q-t_jA),
 \end{aligned}
\tag{39.14}\label{eq:bj-reflection-total}
\]
and both summands in the last line are nef; in particular $-E_j$ is
nef-total in the class $\omega$;
\item\label{itm:ch39-int:limit} for every $v\in X^{\divis}$,
\[
 E_j(v)\longrightarrow D(v)+v(F)\in[0,+\infty],
 \qquad
 D(v):=\lim_jD_j(v)=\sup_jD_j(v).
\tag{39.15}\label{eq:bj-reflection-Ei-limit}
\]
\end{enumerate}
\end{lemma}

\begin{proof}
\cref{itm:ch39-int:nefP}  On a determination $\pi:Z\to Y''$ of $D_j$,
\[
 P_j=\pi^*\bigl((\beta'')^*\theta+\epsilon_jA-\alpha_j\bigr)
      +\bigl(\pi^*\alpha_j-D_j\bigr).
\]
The first summand is the pullback of an ample class, ample by
\eqref{eq:ch39-drift-window}, hence nef by
\Cref{lem:ch38-numerical-birational}\,\cref{itm:ch38-num:inj}; the second is
semiample by hypothesis, hence nef.  A sum of nef classes is nef.

\cref{itm:ch39-int:pos}  $D_j\ge0$, $\overline F\ge0$ and
$\overline G\ge0$, and $0<\delta_j<1$; a convex combination of
effective b-divisors is effective.  That $E_j$ is Cartier over $X$
follows from
\Cref{prop:ch38-evaluation-functorial}\,\cref{itm:ch38-fun:cofinal}: it is
determined on a determination of $D_j$, which is a model over $X$ as
well as over $Y''$.

\cref{itm:ch39-int:order}  Expand:
\[
 \begin{aligned}
 E_{j'}-E_j
 &=(1-\delta_{j'})D'_{j'}+\delta_{j'}\overline G
  -(1-\delta_j)D'_j-\delta_j\overline G\\
 &=(1-\delta_{j'})\bigl(D'_{j'}-D'_j\bigr)
  +(\delta_j-\delta_{j'})D'_j
  -(\delta_j-\delta_{j'})\overline G,
 \end{aligned}
\]
which is the displayed identity.  For $j\le j'$ we have
$D'_{j'}-D'_j=D_{j'}-D_j\ge0$ and
$\delta_j-\delta_{j'}\ge0$, the latter because $(\delta_j)$ is
nonincreasing; and
\[
 D'_j-\overline G=D_j+\overline F-\overline G
 =D_j+\overline R\ \ge\ 0
\]
by \eqref{eq:bj-reflection-fixed-pole}, the trace of $\overline F$ on
$Y''$ being $(\beta'')^*F=G+R$.  All three terms are effective, hence so
is the sum.

\cref{itm:ch39-int:total}  On a determination $\pi:Z\to Y''$ of $D_j$,
write $\varpi:=\beta''\pi:Z\to X$.  Then, using
$\varpi^*\omega=\pi^*(\beta'')^*\omega$,
$\overline F_Z=\varpi^*F=\pi^*(\beta'')^*F$ and
$\overline G_Z=\pi^*G$,
\[
 \begin{aligned}
 \varpi^*\omega-E_j
 &=\pi^*(\beta'')^*\omega
   -(1-\delta_j)\bigl(D_j+\pi^*(\beta'')^*F\bigr)-\delta_j\pi^*G\\
 &=(1-\delta_j)\bigl[\pi^*(\beta'')^*\omega-\pi^*(\beta'')^*F-D_j\bigr]
   +\delta_j\pi^*\bigl[(\beta'')^*\omega-G\bigr]\\
 &=(1-\delta_j)\bigl[\pi^*(\beta'')^*\theta-D_j\bigr]+\delta_j\pi^*Q\\
 &=(1-\delta_j)\bigl[P_j-\epsilon_j\pi^*A\bigr]+\delta_j\pi^*Q\\
 &=(1-\delta_j)P_j+\delta_j\pi^*Q-(1-\delta_j)\epsilon_j\pi^*A,
 \end{aligned}
\]
which is the first line of \eqref{eq:bj-reflection-total}, all classes being
understood on $Z$ and the pullbacks being suppressed in the display of the
statement.  The second line follows from
$(1-\delta_j)\epsilon_j=\delta_jt_j$, which is
\eqref{eq:bj-reflection-delta}.  Finally $P_j$ is nef by
\cref{itm:ch39-int:nefP}, and
\[
 Q-t_jA=(Q-cA)+(c-t_j)A
\]
is nef, being the sum of a nef class and a nonnegative multiple of an ample
class; here $t_j=(1-\delta_j)\delta_j\le\delta_j\le c$.  Pullbacks of nef
classes are nef, so $\varpi^*\omega-E_j$ is nef, which by
\Cref{def:bj-nef-total} is nef-totality of $-E_j$ in the class
$\omega$.

\cref{itm:ch39-int:limit}  Let $v\in X^{\divis}$.  All three of
$D_j(v)$, $v(F)$ and $v(G)$ are finite real numbers, by
\Cref{prop:ch38-evaluation-welldefined}\,\cref{itm:ch38-ev:finite}, and
\[
 E_j(v)=(1-\delta_j)\bigl(D_j(v)+v(F)\bigr)+\delta_jv(G)
\]
by linearity of the evaluation pairing,
\Cref{prop:ch38-evaluation-welldefined}\,\cref{itm:ch38-ev:linear}.  Here
$v(F)$ means the value of $\overline F$ at $v$, which is the same
number whether $\overline F$ is regarded over $X$ or over $Y''$, by
\Cref{prop:ch38-evaluation-functorial}\,\cref{itm:ch38-fun:psi}.  The net
$\bigl(D_j(v)\bigr)_j$ is nondecreasing, hence converges in
$[0,+\infty]$ to its supremum $D(v)$, while $\delta_j\to0$ and
$v(G)$ is a fixed real number.  If $D(v)<\infty$ the limit is
$D(v)+v(F)$ by ordinary arithmetic; if $D(v)=+\infty$ then
$(1-\delta_j)(D_j(v)+v(F))\to+\infty$ while $\delta_jv(G)\to0$, so the
limit is $+\infty=D(v)+v(F)$.  No indeterminate expression arises, because
each $E_j(v)$ is finite and only the limit may be infinite.
\end{proof}

\subsection{The theorem}
\label{sec:ch39-homogeneous-theorem}

\begin{theorem}
\label{thm:bj-homogeneous-reflection}
Let $\beta:Y\to X$ be as in \eqref{eq:ch39-standing} and let
$\theta\in N^1(X)_{\mathbf R}$ be arbitrary.  Pullback is a bijection
\[
 \beta^*:\PSH_{\mathrm{hom},X}(\theta)
       \xrightarrow{\ \sim\ }
       \PSH_{\mathrm{hom},Y}(\beta^*\theta).
\tag{39.16}\label{eq:bj-homogeneous-reflection}
\]
No envelope property is assumed.
\end{theorem}

\begin{proof}
If $\dim X=0$ then $\beta$ is an isomorphism and there is nothing to
prove; assume $\dim X>0$.

Pullback lands in the right space by
\Cref{thm:ch39-pullback}\,\cref{itm:ch39-pb:psh}, and it is injective by
\Cref{thm:ch39-pullback}\,\cref{itm:ch39-pb:order}, since $\beta^{\An}$ is
surjective.  Only surjectivity has to be proved.

\medskip
\noindent\emph{Step 0: normalization of the data.}
Let $\rho:Y''\to Y$, $F$, $G$, $\omega=\theta+[F]$,
$R=(\beta\rho)^*F-G$ and $Q=(\beta\rho)^*\omega-G$ be as in
\Cref{lem:ch39-fixed-pole}.  By \Cref{lem:ch39-model-replacement} it
suffices to prove that $(\beta\rho)^*$ is surjective onto
$\PSH_{\mathrm{hom},Y''}((\beta\rho)^*\theta)$.  We therefore rename, and
assume for the rest of the proof that
\[
 Y=Y'',\qquad \beta=\beta\rho,\qquad
 \text{$F$ on $X$ and $G$ on $Y$ satisfy
 \eqref{eq:bj-reflection-fixed-pole}.}
\]
Fix an ample class $A$ on $Y$ and a rational $c>0$ with $Q-cA$ nef,
which exists because $Q$ is ample.  Note that $F\ne0$, its class being
ample and $\dim X>0$.

\medskip
\noindent\emph{Step 1: regularize upstairs.}
Fix $\psi\in\PSH_{\mathrm{hom},Y}(\beta^*\theta)$.  Apply
\Cref{cor:ch38-gf-regularization-bdivisor} on $Y$ to $\psi$ and the
class $\beta^*\theta$.  It provides a directed set $I$, rational Cartier
b-divisors $C_i=-D_i\le0$ over $Y$, and $\mathbf Q$-line bundles
$L_i$ on $Y$ with $\alpha_i:=c_1(L_i)$, such that
\[
 D_i\ge0,\qquad
 D_i\le D_{i'}\ \ (i\le i'),\qquad
 \alpha_i\longrightarrow\beta^*\theta,
\tag{39.17}\label{eq:bj-homogeneous-regularization-data}
\]
such that on a determination of $D_i$ the class $\pi^*\alpha_i-D_i$ is
semiample, and such that
$\psi_{C_i}=-D_i(\cdot)\in\cH_{\mathrm{hom}}(L_i)$ decreases pointwise to
$\psi$ on $Y^{\An}$.  The three displayed properties are
\Cref{cor:ch38-gf-regularization-bdivisor}\,\cref{itm:ch38-pkg:mono},
\cref{itm:ch38-pkg:class} and \cref{itm:ch38-pkg:neg}.

\medskip
\noindent\emph{Step 2: correct the classes without losing the order.}
Apply \Cref{lem:ch39-drift-window} with $\gamma=\beta^*\theta$ and the net
$(\alpha_i)$, obtaining a directed set $J$, a cofinal order-preserving
map $J\to I$, and nonincreasing rationals
$\epsilon_j=\delta_j^2\downarrow0$ with
\eqref{eq:bj-reflection-class-correction} ample; refine $J$ to a tail so
that $\delta_j\le c$.  Passing to this Willard subnet changes no pointwise
limit, by \Cref{def:appE-willard-subnet}, so
$\psi_{C_j}\searrow\psi$ still holds along $J$.  Form
$P_j$, $t_j$, $D'_j$ and $E_j$ as in \eqref{eq:bj-reflection-Pi},
\eqref{eq:bj-reflection-delta} and \eqref{eq:bj-reflection-Ei}.  By
\Cref{lem:ch39-interpolation}, the b-divisors $-E_j$ are antieffective
real Cartier b-divisors over $X$, they decrease as $j$ increases, and
each is nef-total in the class $\omega$.

\medskip
\noindent\emph{Step 3: realize downstairs.}
For each $j$, \Cref{thm:bj-real-cartier-realization}, applied on the
normal integral projective variety $X$ to the antieffective real Cartier
b-divisor $-E_j$ with nef total class in $\omega$, produces a unique
\[
 \chi_j\in\PSH_{\mathrm{hom},X}(\omega)
\]
whose restriction to $X^{\divis}$ is $\psi_{-E_j}=-E_j(\cdot)$.  No
envelope property is used, by the last sentence of that theorem.  By
\Cref{lem:ch39-interpolation}\,\cref{itm:ch39-int:order} the traces
$-E_j(\cdot)$ decrease on $X^{\divis}$, so
\Cref{cor:ch35-divisorial-determination}\,(1) gives
$\chi_{j'}\le\chi_j$ on all of $X^{\An}$ for $j\le j'$; the net
$(\chi_j)$ is decreasing.  Each $\chi_j$ is homogeneous psh, so
$\chi_j(v_{\triv})=0$ by
\Cref{prop:ch35-homogeneous-normalization}; hence the pointwise limit
\[
                    \chi:=\lim_j\chi_j
\]
satisfies $\chi(v_{\triv})=0$ and is in particular generically finite.
The class $\omega$ is the same at every stage, so PSH2, that is
\Cref{lem:ch39-closure-any-class}\,(4), gives
\[
                    \chi\in\PSH_{\mathrm{hom},X}(\omega),
\tag{39.18}\label{eq:bj-reflection-chi}
\]
homogeneity being preserved because the scaling action commutes with
decreasing limits, \Cref{lem:ch35-action-basic}\,(4).

\medskip
\noindent\emph{Step 4: compute the divisorial trace of $\chi$.}
Let $v\in X^{\divis}=Y^{\divis}$ and put $D(v):=\sup_jD_j(v)\in[0,+\infty]$.
Since $\psi_{C_j}(v)=-D_j(v)$ and $\psi_{C_j}\searrow\psi$,
\[
                    \psi(v)=-D(v)\in[-\infty,0].
\]
By \Cref{lem:ch39-interpolation}\,\cref{itm:ch39-int:limit},
$E_j(v)\to D(v)+v(F)$, so
\[
 \chi(v)=\lim_j\bigl(-E_j(v)\bigr)=-D(v)-v(F)
       =\psi(v)+\log|s_F|(v),
\tag{39.19}\label{eq:bj-reflection-divisorial-identity}
\]
using $\log|s_F|(v)=-v(F)$ from \eqref{eq:divisor-section-pole}.  Both
summands on the right lie in $[-\infty,0]$, so the sum is defined without
appeal to any convention, by ledger entry~(2) of
\Cref{sec:extended-arithmetic}; the only possibly infinite value is
$\psi(v)$, and it is never subtracted.

\medskip
\noindent\emph{Step 5: upgrade to a global identity.}
Consider the two functions
\[
 \beta^*\chi
 \qquad\text{and}\qquad
 \psi+\log|s_{\beta^*F}|
\]
on $Y^{\An}$.  The first is homogeneous $\beta^*\omega$-psh by
\Cref{thm:ch39-pullback}\,\cref{itm:ch39-pb:psh} applied to
\eqref{eq:bj-reflection-chi}.  For the second,
$\log|s_{\beta^*F}|=\beta^*\log|s_F|$ by
\Cref{lem:ch39-beta-an}\,\cref{itm:ch39-ban:sections}, and this function lies
in $\cH_{\mathrm{hom}}(\beta^*F)\subseteq
\PSH_{\mathrm{hom},Y}(\beta^*[F])$; adding $\psi$, which is homogeneous
$\beta^*\theta$-psh, gives a homogeneous
$(\beta^*\theta+\beta^*[F])$-psh function by
\Cref{lem:ch39-closure-any-class}\,(3), the sum being generically finite
because both summands vanish at $v_{Y,\triv}$.  Since
$\beta^*\theta+\beta^*[F]=\beta^*\omega$, the two functions are
homogeneous psh in the same class.  By
\eqref{eq:bj-reflection-divisorial-identity} and
\Cref{lem:ch39-beta-an}\,\cref{itm:ch39-ban:val}, they agree at every point
of $Y^{\divis}=X^{\divis}$.  Therefore
\Cref{cor:ch35-divisorial-determination}\,(2) gives
\[
             \beta^*\chi
             =\psi+\log|s_{\beta^*F}|
             \qquad\text{on }Y^{\An}.
\tag{39.20}\label{eq:bj-reflection-global-identity}
\]
This is the step which the audit of \Cref{sec:ch39-determination} isolates;
its two hypotheses --- both sides psh, and agreement on divisorial points
--- have just been verified, and \Cref{ex:ch39-usc-not-determined-birational}
shows that the first cannot be dropped.

\medskip
\noindent\emph{Step 6: divide out the fixed pole.}
Every homogeneous psh function is nonpositive
(\Cref{prop:ch35-homogeneous-normalization}), so $\psi\le0$ and
\eqref{eq:bj-reflection-global-identity} gives
\[
 \beta^*\chi\le\log|s_{\beta^*F}|=\beta^*\log|s_F|
 \qquad\text{on }Y^{\An}.
\]
By \Cref{thm:ch39-pullback}\,\cref{itm:ch39-pb:order} this descends to
$\chi\le\log|s_F|$ on $X^{\An}$.  Since $X$ is normal and $F$ is a
nonzero effective Cartier divisor,
\Cref{thm:ch35-homogeneous-siu-general} --- which carries no envelope
hypothesis --- yields a unique
\[
 u\in\PSH_{\mathrm{hom},X}(\omega-[F])=\PSH_{\mathrm{hom},X}(\theta)
 \qquad\text{with}\qquad
 \chi=\log|s_F|+u .
\tag{39.21}\label{eq:bj-reflection-siu}
\]
The identity \eqref{eq:bj-reflection-siu} is an identity of sums of
functions with values in $[-\infty,0]$; it is never rearranged into a
difference, and \Cref{warn:ch39-no-subtraction} records what would go wrong
if it were.

\medskip
\noindent\emph{Step 7: conclude.}
Pulling \eqref{eq:bj-reflection-siu} back along $\beta$ --- legitimate
because composition with a map commutes with a pointwise sum of functions
bounded above --- gives
$\beta^*\chi=\log|s_{\beta^*F}|+\beta^*u$.  Comparing with
\eqref{eq:bj-reflection-global-identity},
\[
 \psi+\log|s_{\beta^*F}|=\beta^*u+\log|s_{\beta^*F}|
 \qquad\text{on }Y^{\An}.
\]
At $v\in Y^{\divis}$ the number $\log|s_{\beta^*F}|(v)=-v(F)$ is a real
number, so it may be cancelled: adding a finite constant is injective on
$[-\infty,+\infty)$.  Hence $\psi=\beta^*u$ on $Y^{\divis}$.  Both
functions are homogeneous $\beta^*\theta$-psh --- the second by
\Cref{thm:ch39-pullback}\,\cref{itm:ch39-pb:psh} --- so
\Cref{cor:ch35-divisorial-determination}\,(2) gives $\psi=\beta^*u$ on all
of $Y^{\An}$.  This proves surjectivity, and with it
\eqref{eq:bj-homogeneous-reflection}.
\end{proof}

\begin{warning}{The fixed pole is removed by division, never by subtraction}
\label{warn:ch39-no-subtraction}
In Step~6 one is tempted to define $u:=\chi-\log|s_F|$.  At a
semivaluation $x$ whose support is contained in $\Supp F$ both terms are
$-\infty$, and the difference is undefined by ledger entry~(4) of
\Cref{sec:extended-arithmetic}.  No convention repairs this, because the
correct value of $u$ at two such points can differ: in
\Cref{ex:ch39-blowup-reflection-numbers} the reflected function
$u=\log|\mathfrak m_p|$ takes the value $-\infty$ at $v_{p,\triv}$ and
the value $0$ at $v_{r,\triv}$ for a point $r\in F\smallsetminus\{p\}$,
while $\chi$ and $\log|s_F|$ are both $-\infty$ at both points.  The
same phenomenon in the b-divisor setting is
\Cref{warn:ch38-fixed-pole} and \Cref{ex:ch38-naive-realization-fails}.  The
division of Step~6 is performed by
\Cref{thm:ch35-homogeneous-siu-general}, which produces the identity
$\chi=\log|s_F|+u$ directly and never forms a difference.
\end{warning}

\begin{remark}[Where each hypothesis enters]
\label{rem:ch39-hypothesis-map}
Normality of $X$ is used twice: in
\Cref{thm:bj-real-cartier-realization} in Step~3, and in
\Cref{thm:ch35-homogeneous-siu-general} in Step~6, whose proof needs the
regularity criterion for a rational section with effective divisor on a
normal variety.  Normality of $Y$ is used to have the model category and
to apply \Cref{cor:ch38-gf-regularization-bdivisor} in Step~1.  Projectivity
of $\beta$ is used in \Cref{lem:ch39-fixed-pole}, through
\Cref{thm:appB-admissible-domination}, and birationality is used
everywhere: it is what makes $X^{\divis}=Y^{\divis}$ and what makes
Cartier b-divisors over the two varieties the same objects.  No hypothesis
on $\operatorname{char}k$, on $\dim X$, on the singularities of $Y$
beyond normality, or on the positivity of $\theta$ is used anywhere.
Homogeneity, on the other hand, is used at three separate points: to know
$\psi\le0$ in Step~6, to know $\chi_j(v_{\triv})=0$ in Step~3, and to be
able to apply the homogeneous, envelope-free, form of Siu division rather
than its inhomogeneous form, which would be circular
(\Cref{warn:ch39-firewall}).
\end{remark}

\begin{remark}[What the proof does not do]
\label{rem:ch39-proof-negative}
Three shortcuts are unavailable and are not used.  One cannot define
$\chi$ as the function $x\mapsto\psi(w)$ for a preimage $w$ of $x$,
because fibrewise constancy of $\psi$ is not known in advance; it is a
\emph{consequence} of the theorem, \Cref{cor:ch39-fibre-constant}.  One
cannot define $\chi$ as the upper-semicontinuous regularization on
$X^{\An}$ of the function $v\mapsto\psi(v)$ on $X^{\divis}$: that
function exists, but nothing identifies it as psh, and
\Cref{ex:ch38-usc-not-psh} shows that in the analogous b-divisor situation
it need not be.  And one cannot dispense with the fixed pole $F$ by
applying \Cref{thm:bj-real-cartier-realization} directly to $-D_j$ in the
class $\theta$: the class $\beta^*\theta-D_j$ is not nef in general,
because $\theta$ itself need not be nef
(\Cref{warn:ch38-theta-not-nef}) and because the drift
$\alpha_j-\beta^*\theta$ has no sign.
\end{remark}

\subsection{Consequences}
\label{sec:ch39-homogeneous-consequences}

\begin{corollary}
\label{cor:ch39-fibre-constant}
Every $\psi\in\PSH_{\mathrm{hom},Y}(\beta^*\theta)$ is constant on every
fibre of $\beta^{\An}$.  In particular, for a closed point $p\in X$ over
which $\beta$ has positive dimensional fibre $F_p=\beta^{-1}(p)$,
\[
 \psi\ \text{is constant on}\ F_p^{\An},
 \qquad\text{with value}\ \chi(v_{p,\triv})\in[-\infty,0],
\]
where $\chi$ is the reflection of $\psi$.
\end{corollary}

\begin{proof}
By \Cref{thm:bj-homogeneous-reflection}, $\psi=\beta^*\chi$ for some
$\chi\in\PSH_{\mathrm{hom},X}(\theta)$, and a pullback is constant on
fibres by definition.  The value on
$(\beta^{\An})^{-1}(v_{p,\triv})=F_p^{\An}$
(\Cref{prop:ch39-fibres}\,\cref{itm:ch39-fib:point}) is
$\chi(v_{p,\triv})$, which is $\le0$ by
\Cref{prop:ch35-homogeneous-normalization}.
\end{proof}

\begin{corollary}
\label{cor:ch39-hom-structure}
The bijection \eqref{eq:bj-homogeneous-reflection} is an isomorphism of
ordered convex sets, compatible with the natural operations at their correct
class level:
\begin{enumerate}
\item\label{itm:ch39-hs:order} $\chi_1\le\chi_2$ if and only if
 $\beta^*\chi_1\le\beta^*\chi_2$;
\item\label{itm:ch39-hs:ops} $\beta^*\max\{\chi_1,\chi_2\}
 =\max\{\beta^*\chi_1,\beta^*\chi_2\}$,
 $\beta^*(s\chi_1+t\chi_2)=s\beta^*\chi_1+t\beta^*\chi_2$ for
 $s,t\ge0$ with $s+t=1$, and
 $\beta^*(t\cdot\chi)=t\cdot\beta^*\chi=\beta^*\chi$, the last because
 homogeneous functions are fixed by the scaling action;
\item\label{itm:ch39-hs:limits} $\beta^*$ carries decreasing nets to
 decreasing nets and commutes with their limits, in both directions;
\item\label{itm:ch39-hs:trace} $\chi$ and $\beta^*\chi$ have the same
 restriction to $X^{\divis}=Y^{\divis}$, and this restriction determines
 each of them;
\item\label{itm:ch39-hs:bdiv} $\beta^*$ matches the realization functors
 of \Cref{thm:bj-real-cartier-realization}: for every antieffective real
 Cartier b-divisor $B$ over $X$ with nef total class in $\theta$, one
 has $\beta^*\chi^X_B=\chi^Y_B$.
\end{enumerate}
\end{corollary}

\begin{proof}
\cref{itm:ch39-hs:order} is
\Cref{thm:ch39-pullback}\,\cref{itm:ch39-pb:order}, and
\cref{itm:ch39-hs:ops}, \cref{itm:ch39-hs:limits} are
\Cref{thm:ch39-pullback}\,\cref{itm:ch39-pb:ops}; finite maxima and convex
combinations remain in the fixed class by
\Cref{lem:ch39-closure-any-class}\,(2)--(3), and homogeneity is preserved by
\Cref{prop:ch35-homogeneous-normalization}.  More generally, positive linear
combinations are compatible with pullback in the graded family as the class
varies from $\theta$ to $(s+t)\theta$; they are not operations on the
single fixed-class space unless $s+t=1$.
\cref{itm:ch39-hs:trace} is
\Cref{lem:ch39-beta-an}\,\cref{itm:ch39-ban:val} together with
\Cref{cor:ch35-divisorial-determination}\,(2).  \cref{itm:ch39-hs:bdiv} is
\Cref{prop:ch39-cartier-reflection}, once one notes that a determination of
$B$ may be replaced by one factoring through $\beta$, by
\Cref{lem:appB-common-domination} and
\Cref{lem:ch38-determination-equivalence}\,\cref{itm:ch38-de:family}.
\end{proof}

\begin{remark}[The gap left open in \Cref{chap:bj-cartier-realization},
and what actually closes it]
\label{rem:ch39-closing-the-gap}
\Cref{cor:ch38-realization-pullback} proved
$\chi^{Y}_B\le\beta^*\chi^X_B$ on $Y^{\An}$ with equality on
$Y^{\val}$, and \Cref{rem:ch38-pullback-forward} announced that equality
extends to $Y^{\An}$ by \Cref{thm:bj-homogeneous-reflection}.  It is worth
recording that the extension is in fact cheaper than announced.  What is
needed is only that $\beta^*\chi^X_B$ be $\beta^*\theta$-psh, and that
is \Cref{thm:ch39-pullback}\,\cref{itm:ch39-pb:psh}, whose proof consists of
pulling back a Fubini--Study presentation and invoking
\Cref{def:bj-psh} with the classes $\beta^*c_1(L_i)\to\beta^*\theta$.
This is available because \Cref{def:bj-psh} allows the classes to vary; had
the definition demanded a fixed class, one would have needed the hypothesis
of \Cref{lem:ch38-pullback-stages}\,\cref{itm:ch38-pb:limit} that
$\omega-c_1(L_n)$ be nef, which is exactly what
\Cref{chap:bj-cartier-realization} had to arrange by hand.  The full
reflection theorem is needed for the \emph{surjectivity} of $\beta^*$, not
for the identity \eqref{eq:ch39-cartier-reflection}.
\end{remark}

\begin{corollary}
\label{cor:ch39-model-independence}
Let $Y_1,Y_2$ be normal projective models over $X$.  Then
$\PSH_{\mathrm{hom},Y_1}(\beta_1^*\theta)$ and
$\PSH_{\mathrm{hom},Y_2}(\beta_2^*\theta)$ are canonically isomorphic,
compatibly with the identification of both with
$\PSH_{\mathrm{hom},X}(\theta)$; the isomorphism is the identity on
divisorial traces.  Consequently the homogeneous psh space in a pulled-back
class is a birational invariant of the pair $(X,\theta)$.
\end{corollary}

\begin{proof}
Both are identified with $\PSH_{\mathrm{hom},X}(\theta)$ by
\Cref{thm:bj-homogeneous-reflection}, and the identifications preserve
divisorial traces by
\Cref{cor:ch39-hom-structure}\,\cref{itm:ch39-hs:trace}.  For a common
domination $Y_3\to Y_1,Y_2$ --- which exists by
\Cref{lem:appB-common-domination} --- the three identifications are
compatible because pullback is functorial:
$(\beta_1\rho_1)^*=\rho_1^*\beta_1^*$.
\end{proof}

\subsection{Worked examples}
\label{sec:ch39-homogeneous-examples}

\begin{example}[Blowing up one point of the projective plane]
\label{ex:ch39-blowup-point}
Let $X=\mathbf P^2_k$, choose a closed point $p$, and let
$\beta:Y=\operatorname{Bl}_pX\to X$ have exceptional divisor $E$.
For $L=\cO_X(1)$, the sheaf $L\otimes\mathfrak m_p$ is generated by the
two linear forms vanishing at $p$.  Hence
\[
             \psi_X:=\log|\mathfrak m_p|
             \in\cH_{\mathrm{hom},X}(L)
\]
by \Cref{prop:hom-fs-ideal}.  Since
$\mathfrak m_p\cO_Y=\cO_Y(-E)$, pullback gives
\[
       \beta^*\psi_X
          =\log|\cO_Y(-E)|=:\psi_Y,
\]
by \Cref{lem:ch39-beta-an}\,\cref{itm:ch39-ban:sections}.
At $v_c=c\,\ord_E$, one has $v_c(\mathfrak m_p)=c$, and therefore
\[
                  \psi_Y(v_c)=-c,
       \qquad \psi_Y(v_{\triv})=0.
\]
On the exceptional fibre $E^{\An}=(\beta^{\An})^{-1}(v_{p,\triv})$ of
\Cref{ex:ch39-fibre-plane}, every point $w$ has $s(w)\subseteq E$, so
$w(\cO_Y(-E))=+\infty$ and $\psi_Y\equiv-\infty$ there; correspondingly
$\psi_X(v_{p,\triv})=-\infty$, and \Cref{cor:ch39-fibre-constant} is
visible.

Thus the exceptional pole upstairs is exactly the pullback of the point
ideal downstairs; it is not an additional psh function created by the
blowup.  This is the basic local picture behind
\Cref{thm:bj-homogeneous-reflection}.
\end{example}

\begin{example}[A birational reflection computed on a blown-up surface]
\label{ex:ch39-blowup-reflection-numbers}
\index{reflection!computed on a surface}
We run the entire proof of \Cref{thm:bj-homogeneous-reflection} with
explicit numbers, on the simplest example in which the class drift of
Step~2 is genuinely nonzero.

\emph{The geometry.}  Keep $X=\mathbf P^2_k$, $p$, $\beta:Y\to X$ and
$E$ as in \Cref{ex:ch39-blowup-point}, and write
$H:=\beta^*c_1(\cO_X(1))$, $e:=[E]$, so that
\[
 H^2=1,\qquad H\cdot e=0,\qquad e^2=-1 .
\tag{39.22}\label{eq:ch39-surface-intersection}
\]
The Mori cone of $Y$ is generated by $e$ and by the class $H-e$ of the
strict transform of a line through $p$, these being the only irreducible
curves of nonpositive self-intersection.  By the Nakai--Moishezon criterion
\cite[Theorem~1.2.23]{PAG1}, for $a,s\in\mathbf R$,
\[
 aH-se\ \text{is ample}\iff 0<s<a,
 \qquad
 aH-se\ \text{is nef}\iff 0\le s\le a .
\tag{39.23}\label{eq:ch39-surface-ample-test}
\]
Indeed $(aH-se)\cdot e=s$, $(aH-se)\cdot(H-e)=a-s$, and
$(aH-se)^2=a^2-s^2$.

\emph{The datum upstairs.}  Take $\theta:=c_1(\cO_X(1))$, so that
$\beta^*\theta=H$, and let
\[
 \psi:=\log|\mathfrak m_p\cO_Y|=\log|\cO_Y(-E)|
 \in\PSH_{\mathrm{hom},Y}(\beta^*\theta),
\]
whose Cartier b-divisor is $-E$, determined on $Y$, with total class
$H-e$ nef by \eqref{eq:ch39-surface-ample-test}.  We pretend not to know
the answer and produce it by the algorithm.

\emph{Step 1: an approximating net with drift.}  For an integer $m\ge2$
set
\[
 D_m:=(1-m^{-1})E,
 \qquad
 L_m:=\beta^*\cO_X(1)-m^{-1}E,
 \qquad
 \alpha_m:=c_1(L_m)=H-m^{-1}e .
\tag{39.24}\label{eq:ch39-surface-stages}
\]
On the determination $Y$ itself the total class is
$\alpha_m-[D_m]=H-e$, which is semiample: the pencil of strict transforms
of lines through $p$ is base point free.  Hence
$\psi_{-D_m}=(1-m^{-1})\log|\cO_Y(-E)|\in\cH_{\mathrm{hom}}(L_m)$ by
\Cref{prop:ch38-model-fs-dictionary}\,\cref{itm:ch38-dict:backward}.  The
divisors $D_m$ increase, $\psi_{-D_m}\searrow\psi$, and
$\alpha_m\to H=\beta^*\theta$; this is exactly the data
\eqref{eq:bj-homogeneous-regularization-data}.  Note that
$\alpha_m\notin\beta^*N^1(X)_{\mathbf R}$ for any $m$: its coefficient
along $e$ is $-m^{-1}\ne0$.  So no stage is the pullback of anything,
and \Cref{prop:ch39-cartier-reflection} cannot be applied stagewise.

\emph{Step 0: the fixed-pole data.}  Let $F=L\subset X$ be a line through
$p$, with strict transform $L'$ on $Y$, so that
$\beta^*F=L'+E$ and $[\beta^*F]=H$.  Then $[F]=c_1(\cO_X(1))$ is ample
and $\omega=\theta+[F]=c_1(\cO_X(2))$ is ample.  Put
\[
 G:=\tfrac12E,\qquad
 R:=\beta^*F-G=L'+\tfrac12E\ge0,\qquad
 Q:=\beta^*\omega-G=2H-\tfrac12e .
\tag{39.25}\label{eq:ch39-surface-fixed-pole}
\]
Then $2G=E\le L'+E=\beta^*F$, and $Q$ is ample by
\eqref{eq:ch39-surface-ample-test} with $a=2$, $s=\frac12$.  Take
$A:=2H-e$, ample by \eqref{eq:ch39-surface-ample-test}, and $c:=\frac12$:
\[
 Q-cA=(2-2c)H+(c-\tfrac12)e=H,
\]
which is nef.  Thus \eqref{eq:bj-reflection-fixed-pole} holds with no model
replacement, $\beta$ being already a blowup.

\emph{Step 2: the drift window.}  By
\eqref{eq:ch39-surface-ample-test}, the class correction
\eqref{eq:bj-reflection-class-correction} is
\[
 \beta^*\theta+\epsilon A-\alpha_m
 =H+\epsilon(2H-e)-H+m^{-1}e
 =2\epsilon H-(\epsilon-m^{-1})e,
\]
which is ample if and only if $0<\epsilon-m^{-1}<2\epsilon$, that is, if
and only if $\epsilon>1/m$.  This is the whole point of
\Cref{lem:ch39-drift-window}: the auxiliary parameter must \emph{dominate}
the drift, and it must still tend to $0$.  Choose the cofinal subnet
$m=n^4$, $n\ge2$, with
\[
 \epsilon_n:=n^{-2}>n^{-4}=1/m,
 \qquad
 \delta_n=\sqrt{\epsilon_n}=n^{-1},
 \qquad
 t_n=(1-n^{-1})n^{-1}\le\tfrac14\le c .
\tag{39.26}\label{eq:ch39-surface-window}
\]

\emph{Step 2, continued: the interpolated b-divisors.}  With
$D'_n=D_{n^4}+\overline F$ and
$E_n=(1-\delta_n)D'_n+\delta_n\overline G$, the trace on $Y$ is
\[
 E_n=(1-n^{-1})\,L'+f(n)\,E,
 \qquad
 f(n):=(1-n^{-1})(2-n^{-4})+\tfrac1{2n}.
\tag{39.27}\label{eq:ch39-surface-En}
\]
Numerically $f(2)=\frac{39}{32}=1.21875$,
$f(3)=\frac{322}{243}+\frac16\approx1.49177$,
$f(4)=\frac{1533}{1024}+\frac18\approx1.62207$, increasing to $2$; and
the coefficient $1-n^{-1}$ of $L'$ increases to $1$.  Thus
$E_n\uparrow L'+2E=\beta^*F+E$, in accordance with
\eqref{eq:bj-reflection-Ei-limit}.

\emph{The exact total class.}  From \eqref{eq:ch39-surface-En},
\[
 \beta^*\omega-[E_n]
 =2H-(1-n^{-1})(H-e)-f(n)e
 =(1+n^{-1})H-\bigl(f(n)-1+n^{-1}\bigr)e .
\tag{39.28}\label{eq:ch39-surface-total}
\]
At $n=2$ this is $\frac32H-\frac{23}{32}e$, which is ample by
\eqref{eq:ch39-surface-ample-test}.  The identity
\eqref{eq:bj-reflection-total} is verified at $n=2$ as follows:
\[
 P_2=\beta^*\theta+\epsilon_2A-D_{16}
 =H+\tfrac14(2H-e)-\tfrac{15}{16}e=\tfrac32H-\tfrac{19}{16}e,
 \qquad
 Q-t_2A=\tfrac32H-\tfrac14e,
\]
and
\[
 (1-\delta_2)P_2+\delta_2(Q-t_2A)
 =\tfrac12\Bigl(\tfrac32H-\tfrac{19}{16}e\Bigr)
 +\tfrac12\Bigl(\tfrac32H-\tfrac14e\Bigr)
 =\tfrac32H-\tfrac{23}{32}e,
\]
which agrees with \eqref{eq:ch39-surface-total}.  Both summands are nef:
$P_2=\frac32H-\frac{19}{16}e$ has $0\le\frac{19}{16}\le\frac32$, and
$Q-t_2A=\frac32H-\frac14e$ has $0\le\frac14\le\frac32$.

\emph{Step 3: realization downstairs.}  Each $-E_n$ is antieffective,
Cartier over $X$, determined on $Y$, with nef total class in
$\omega=c_1(\cO_X(2))$.  \Cref{thm:bj-real-cartier-realization} gives
$\chi_n\in\PSH_{\mathrm{hom},X}(\omega)$ with divisorial trace
\[
 \chi_n(\ord_L)=-(1-n^{-1}),\qquad
 \chi_n(\ord_E)=-f(n),\qquad
 \chi_n(\ord_C)=0
\]
for every prime divisor $C\subset X$ other than $L$.  These decrease in
$n$, and in the limit
\[
 \chi(\ord_L)=-1,\qquad \chi(\ord_E)=-2,\qquad \chi(\ord_C)=0 .
\tag{39.29}\label{eq:ch39-surface-chi-trace}
\]
More invariantly, every $E_n$ is determined on the fixed model $Y$ and
its trace there converges coefficientwise to $L'+2E$.  Hence for every
$v\in X^{\divis}=Y^{\divis}$, including divisors on models strictly above
$Y$,
\[
                         \chi(v)=-(L'+2E)(v).
\]

\emph{Steps 5--7: identification.}  Let $\ell$ be the linear form cutting
out $L$.  The ideal $\mathcal J:=\ell\cdot\mathfrak m_p\subset\cO_X$
satisfies: $\cO_X(2)\otimes\mathcal J$ is generated by the two conics
$\ell x$ and $\ell y$, where $x,y$ are linear forms vanishing at
$p$; hence $\log|\mathcal J|\in\cH_{\mathrm{hom},X}(\cO_X(2))$ by
\Cref{prop:hom-fs-ideal}.  On $Y$ one has
\[
 \mathcal J\cO_Y
 =\cO_Y(-L'-E)\,\cO_Y(-E)
 =\cO_Y(-L'-2E),
\]
so the full divisorial trace of $\log|\mathcal J|$ is
$v\mapsto-(L'+2E)(v)$, not merely the three displayed sample values in
\eqref{eq:ch39-surface-chi-trace}.  By divisorial determination,
$\chi=\log|\mathcal J|=\log|s_F|+\log|\mathfrak m_p|$, which is
\eqref{eq:bj-reflection-siu} with
\[
 u=\log|\mathfrak m_p|\in\PSH_{\mathrm{hom},X}(\theta),
 \qquad
 \beta^*u=\log|\mathfrak m_p\cO_Y|=\psi .
\]
The algorithm has returned the expected answer, and every intermediate
object was an explicit divisor or an explicit class.
\end{example}

\begin{example}[The Siu division in
\Cref{ex:ch39-blowup-reflection-numbers} is borderline]
\label{ex:ch39-reflection-borderline}
Continue with the data of \Cref{ex:ch39-blowup-reflection-numbers} and let
\[
 K_F=\{v\in X^{\An}:\log|s_F|(v)=-1\}
\]
be the compact set \eqref{eq:ch35-siu-slice} attached to $F=L$.  The
divisorial point $\ord_L$ lies in $K_F$, since
$\ord_L(\ell)=1$, and there
\[
 \chi(\ord_L)=\log|s_F|(\ord_L)+\log|\mathfrak m_p|(\ord_L)=-1+0=-1 ,
\]
because $\mathfrak m_p$ is the unit ideal at the generic point of $L$.
Since $\chi\le\log|s_F|$ everywhere, this gives
\[
 \sup_{K_F}\chi=-1 .
\]
Thus the strict hypothesis \eqref{eq:ch35-siu-strict-hypothesis} of
\Cref{thm:ch35-homogeneous-siu} \emph{fails}, by
\Cref{lem:ch35-slice-to-cone}, and the borderline form
\Cref{thm:ch35-homogeneous-siu-general} is the statement actually used in
Step~6 of the proof.  This is the phenomenon predicted in
\Cref{ex:ch38-strict-room-fails}: the activation divisor $F$ is required
to contain the whole non-isomorphism locus, so its support meets the locus
where the b-divisor has coefficient zero, and enlarging $F$ enlarges the
pole by the same amount.  Since
\Cref{thm:ch35-homogeneous-siu-general} carries no envelope hypothesis, the
firewall \Cref{warn:ch39-firewall} is respected.
\end{example}

\begin{example}[A homogeneous psh function that is not fibrewise constant]
\label{ex:ch39-not-fibre-constant}
Keep $X=\mathbf P^2$, $\beta:Y=\operatorname{Bl}_pX\to X$ and $E$, and
choose a closed point $q\in E$.  Let $A$ be an ample line bundle on
$Y$ so positive that $A\otimes\mathcal I_q$ is globally generated, which
is possible by Serre's theorem.  Then
\[
 \psi':=\log|\mathcal I_q|\in\cH_{\mathrm{hom}}(A)
 \subseteq\PSH_{\mathrm{hom},Y}(c_1(A))
\]
by \Cref{prop:hom-fs-ideal}.  Evaluate $\psi'$ at two points of the fibre
$E^{\An}=(\beta^{\An})^{-1}(v_{p,\triv})$ of \Cref{ex:ch39-fibre-plane}:
\[
 \psi'(v_{q,\triv})=-v_{q,\triv}(\mathcal I_q)=-\infty,
 \qquad
 \psi'(v_{E,\triv})=-v_{E,\triv}(\mathcal I_q)=0,
\]
the second because $\mathcal I_q$ is the unit ideal at the generic point
of $E$.  So $\psi'$ is \emph{not} constant on this fibre, and therefore
is not the pullback of any function on $X^{\An}$.

This is consistent with \Cref{cor:ch39-fibre-constant} only because
$c_1(A)\notin\beta^*N^1(X)_{\mathbf R}$: a pullback class is trivial on
the contracted curve $E$, by
\Cref{lem:ch38-numerical-birational}\,\cref{itm:ch38-num:notample}, whereas
$c_1(A)\cdot E>0$.  The example therefore isolates what
\Cref{thm:bj-homogeneous-reflection} really asserts.  It is not that
homogeneous psh functions on $Y$ descend --- most do not --- but that
those in a \emph{pulled-back class} do, and the mechanism producing the
descent is the realization of an interpolated b-divisor downstairs, not any
visible property of the function upstairs.
\end{example}

\section{From equality on divisorial points to equality everywhere}
\label{sec:ch39-determination}

Both proofs of this chapter end with the same manoeuvre: an identity
verified at every point of $X^{\divis}=Y^{\divis}$ is promoted to an
identity on the whole analytification.  The proof map of
\Cref{chap:bj-envelope-statement} singles this passage out for audit,
because it is the point at which a statement about a countable-looking set
of valuations becomes a statement about a compact space containing points of
proper support, at which psh functions may take the value $-\infty$ and at
which evaluation is not continuous
(\Cref{ex:ch35-evaluation-discontinuous}).  We therefore isolate the step,
display its hypotheses, and exhibit the failures that occur when they are
weakened.

\subsection{The package}
\label{sec:ch39-determination-package}

\begin{proposition}
\label{prop:ch39-determination-package}
Let $X$ be an integral projective $k$-variety.
\begin{enumerate}
\item\label{itm:ch39-det:compare}
 \emph{(Comparison.)}  If $\phi\in\PSH_X(\theta)$ for some
 $\theta\in N^1(X)_{\mathbf R}$, and if
 $\eta:X^{\An}\to\mathbf R\cup\{-\infty\}$ is upper semicontinuous, then
 $\phi\le\eta$ on $X^{\divis}$ implies $\phi\le\eta$ on $X^{\An}$.
\item\label{itm:ch39-det:equal}
 \emph{(Determination.)}  If $\phi\in\PSH_X(\theta)$ and
 $\phi'\in\PSH_X(\theta')$, with $\theta,\theta'\in N^1(X)_{\mathbf R}$
 possibly different, and if $\phi=\phi'$ on $X^{\divis}$, then
 $\phi=\phi'$ on $X^{\An}$.
\item\label{itm:ch39-det:minimal}
 \emph{(Extremality.)}  A $\theta$-psh $\phi$ is the smallest
 upper-semicontinuous extension of $\phi|_{X^{\divis}}$ to $X^{\An}$.
\item\label{itm:ch39-det:birational}
 \emph{(Birational form.)}  Let $\beta:Y\to X$ be as in
 \eqref{eq:ch39-standing} and identify $X^{\divis}=Y^{\divis}$.  If
 $u\in\PSH_Y(\eta_1)$ and $u'\in\PSH_Y(\eta_2)$ satisfy $u=u'$ on
 $X^{\divis}$, then $u=u'$ on $Y^{\An}$.  In particular, for
 $\chi\in\PSH_X(\theta)$ and $u\in\PSH_Y(\eta)$,
\[
 \beta^*\chi=u\ \text{on }X^{\divis}
 \iff
 \beta^*\chi=u\ \text{on }Y^{\An},
\tag{39.30}\label{eq:ch39-determination-birational}
\]
 the implication $\Rightarrow$ using
 \Cref{thm:ch39-pullback}\,\cref{itm:ch39-pb:psh} to know that
 $\beta^*\chi$ is psh.
\end{enumerate}
\end{proposition}

\begin{proof}
\cref{itm:ch39-det:compare} is \Cref{thm:ch35-divisorial-comparison},
\cref{itm:ch39-det:equal} and \cref{itm:ch39-det:minimal} are
\Cref{cor:ch35-divisorial-determination}\,(2) and the final sentence of
\Cref{thm:ch35-divisorial-comparison}.  For
\cref{itm:ch39-det:birational}, apply \cref{itm:ch39-det:equal} on $Y$,
using $X^{\divis}=Y^{\divis}$ from
\Cref{lem:ch39-beta-an}\,\cref{itm:ch39-ban:val}; and for
\eqref{eq:ch39-determination-birational} note that $\beta^*\chi$ is
$\beta^*\theta$-psh, so that \cref{itm:ch39-det:equal} applies to the pair
$(\beta^*\chi,u)$, which lies in $\PSH_Y(\beta^*\theta)\times
\PSH_Y(\eta)$.
\end{proof}

Three features of \Cref{prop:ch39-determination-package} deserve emphasis,
because each is used and each is easy to misremember.

\emph{The two classes may differ.}  In \cref{itm:ch39-det:equal} nothing
relates $\theta$ and $\theta'$.  This is what allows Step~5 of the proof
of \Cref{thm:bj-homogeneous-reflection} to compare $\beta^*\chi$, a
priori only known to be $\beta^*\omega$-psh, with
$\psi+\log|s_{\beta^*F}|$, and what allows Step~7 to compare
$\psi$, which is $\beta^*\theta$-psh, with $\beta^*u$.  Had the
statement required a common class, each comparison would have needed a
separate verification that the classes agree --- which they do, but only
after the very identity being proved.

\emph{The inequality is one-sided in the hypotheses.}  In
\cref{itm:ch39-det:compare} the \emph{smaller} function must be psh and the
larger one only upper semicontinuous.  Reversing the roles makes the
statement false, by \Cref{ex:ch35-usc-not-determined}.  Every application in
this chapter is checked against this asymmetry; \Cref{rem:ch35-determination-uses}
lists the applications in \Cref{chap:bj-homogeneity-bdivisors} and the same
discipline is followed here.

\emph{Density is not the mechanism.}  $X^{\divis}$ is dense in
$X^{\An}$ by \Cref{thm:triv-divisorial-density}, but density gives only
that a \emph{lower} semicontinuous function is dominated by an
upper-semicontinuous one, \Cref{lem:ch35-density-gives}, and both functions
here are upper semicontinuous.  The actual mechanism is the attainment of
suprema on a finite set of divisorial points,
\Cref{lem:ch35-finite-attainment} and \Cref{lem:ch35-attainment-psh}, which
is a genuinely plurisubharmonic input resting on Rees valuations.

\begin{warning}{What ``determined by divisorial values'' does not mean}
\label{warn:ch39-determination-hypotheses}
\Cref{prop:ch39-determination-package} asserts that the restriction map
\[
 \PSH_X(\theta)\longrightarrow
 \bigl(\mathbf R\cup\{-\infty\}\bigr)^{X^{\divis}}
\]
is injective.  It does not assert that the inverse is continuous, that a
given function on $X^{\divis}$ extends, or that pointwise convergence on
$X^{\divis}$ implies pointwise convergence on $X^{\An}$.  All three
fail: continuity of evaluation at a proper-support point fails by
\Cref{ex:ch35-evaluation-discontinuous}, and the extension problem is
precisely what \Cref{thm:bj-real-cartier-realization} solves, and only under
a numerical hypothesis.  The proper-support locus
$\{x\in X^{\An}:s(x)\subsetneq X\}$ of
\Cref{def:ch32-proper-support} supplies the explicit failure examples used
here.  In general \Cref{thm:bj-basic-closure}\,(7) guarantees finiteness
only on $X^{\lin}$, not on the whole full-support locus
$X^{\val}$; no stronger determination statement is used in this
chapter.
\end{warning}

\subsection{Two functions with the same divisorial trace}
\label{sec:ch39-determination-examples}

\begin{example}[An upper-semicontinuous counterexample]
\label{ex:ch39-usc-not-determined-birational}
Keep $X=\mathbf P^2$, $\beta:Y=\operatorname{Bl}_pX\to X$ with
exceptional prime $E$, and let
\[
 f_1:=\log|\cO_Y(-E)|=\beta^*\log|\mathfrak m_p|
 \in\PSH_{\mathrm{hom},Y}(\beta^*\theta),
 \qquad\theta=c_1(\cO_X(1)),
\]
as in \Cref{ex:ch39-blowup-point}.  Fix a closed point $q\in E$ and set
\[
 h:=\begin{cases}-1,&\text{at }v_{q,\triv},\\ -\infty,&\text{elsewhere},
 \end{cases}
 \qquad
 f_2:=\max\{f_1,h\}.
\tag{39.31}\label{eq:ch39-two-usc}
\]
\emph{$h$ is upper semicontinuous.}  The singleton $\{v_{q,\triv}\}$ is
the analytification of the closed subscheme $\{q\}\subset Y$, hence closed
in $Y^{\An}$ by
\Cref{prop:ch32-support-partition}\,\cref{itm:ch32-sp:closed}; the sets
$\{h\ge c\}$ are $\varnothing$ for $c>-1$, $\{v_{q,\triv}\}$ for
$-\infty<c\le-1$, and $Y^{\An}$ for $c=-\infty$, all closed.  A finite
maximum of upper-semicontinuous functions is upper semicontinuous, so
$f_2$ is upper semicontinuous.

\emph{The two functions have the same divisorial trace.}  Every point of
$Y^{\divis}$ has full support, whereas $s(v_{q,\triv})=\{q\}$; so
$h\equiv-\infty$ on $Y^{\divis}$ and $f_2=f_1$ there.

\emph{They differ.}  Since $s(v_{q,\triv})=\{q\}\subseteq E$, the section
$s_E$ vanishes identically on the support, so
$v_{q,\triv}(\cO_Y(-E))=+\infty$ and
$f_1(v_{q,\triv})=-\infty$, while $f_2(v_{q,\triv})=-1$.

\emph{Consequences.}  First, $f_2\notin\PSH_Y(\eta)$ for any
$\eta\in N^1(Y)_{\mathbf R}$: otherwise
\Cref{prop:ch39-determination-package}\,\cref{itm:ch39-det:equal} applied to
the pair $(f_1,f_2)$ would force $f_1=f_2$.  Second, $f_2$ is not the
pullback of any function on $X^{\An}$: it is not constant on the fibre
$E^{\An}=(\beta^{\An})^{-1}(v_{p,\triv})$, taking the value $-1$ at
$v_{q,\triv}$ and $-\infty$ at $v_{E,\triv}$.  Third,
$f_1\le f_2$ with equality on $Y^{\divis}$, which is consistent with
\Cref{prop:ch39-determination-package}\,\cref{itm:ch39-det:minimal}: the psh
function is the \emph{smallest} upper-semicontinuous extension of its
divisorial trace, and $f_2$ is a larger one.

\emph{Which property of psh functions fails.}  The function $f_2$ is not
decreasing for the order \eqref{eq:semivaluation-order}.  Indeed
$t\,\ord_E\le v_{q,\triv}$ for every $t>0$: a function $g$ with
$\ord_E(g)=a>0$ lies in the $a$-th power of the ideal of $E$, hence
vanishes at $q$, so $v_{q,\triv}(g)=+\infty\ge ta$; and a function with
$\ord_E(g)=0$ has $v_{q,\triv}(g)\ge0$.  Since
$f_2(t\,\ord_E)=f_1(t\,\ord_E)=-t\to-\infty$ while
$f_2(v_{q,\triv})=-1$, monotonicity fails.  Every psh function is
decreasing, by
\Cref{prop:ch34-generic-finiteness}\,\cref{itm:ch34-gf:decreasing} and
passage to decreasing limits, which is why this particular defect cannot
occur inside the class.  It should not be concluded that upper
semicontinuity plus monotonicity is the hypothesis of
\Cref{thm:ch35-divisorial-comparison}: that theorem is proved through the
finite-attainment lemmas, and its hypothesis is plurisubharmonicity of the
smaller function.
\end{example}

\begin{example}[A one-parameter family of upper-semicontinuous extensions]
\label{ex:ch39-two-usc-extensions}
Keep the notation of \Cref{ex:ch39-usc-not-determined-birational} and, for
$c\in[-\infty,0]$, define
\[
 f^{(c)}:=\max\{f_1,\ c\cdot\one_{E^{\An}}\}
 =\begin{cases}
   c,&\text{on }E^{\An},\\
   f_1,&\text{elsewhere},
 \end{cases}
\]
where $c\cdot\one_{E^{\An}}$ denotes the function equal to $c$ on
$E^{\An}$ and $-\infty$ off it.  Since $E^{\An}$ is closed in
$Y^{\An}$ by
\Cref{prop:ch32-support-partition}\,\cref{itm:ch32-sp:closed} and
$f_1\equiv-\infty$ on $E^{\An}$, each $f^{(c)}$ is upper
semicontinuous, and all of them restrict to the same function on
$Y^{\divis}$, namely $f_1|_{Y^{\divis}}$, because
$E^{\An}\cap Y^{\divis}=\varnothing$.  So a single divisorial datum has a
continuum of upper-semicontinuous extensions, ordered by $c$, of which
exactly one --- the smallest, $c=-\infty$ --- is psh.

The example also shows that fibrewise constancy is not the missing
condition: every $f^{(c)}$ \emph{is} constant on the fibre $E^{\An}$,
and $f^{(c)}=\beta^*\chi^{(c)}$ for the function $\chi^{(c)}$ on
$X^{\An}$ equal to $c$ at $v_{p,\triv}$ and to
$\log|\mathfrak m_p|$ elsewhere.  Thus $\chi^{(c)}$ is an
upper-semicontinuous function on $X^{\An}$ with the same divisorial trace
as $\log|\mathfrak m_p|$ and different values, and the reflection
statement \eqref{eq:bj-homogeneous-reflection} would be false if
$\PSH_{\mathrm{hom}}$ were replaced by ``homogeneous upper-semicontinuous
with the given trace''.  What distinguishes $c=-\infty$ is not the
birational geometry but plurisubharmonicity.
\end{example}

\begin{remark}[The audit entry]
\label{rem:ch39-audit-divisorial}
The proof-audit list of \Cref{chap:bj-envelope-statement} contains the item
``passage from equality on divisorial points to equality everywhere''.  In
this chapter that passage occurs exactly four times, and in each case both
functions are psh and the citation is
\Cref{prop:ch39-determination-package}\,\cref{itm:ch39-det:equal}:
\begin{enumerate}
\item in \Cref{prop:ch39-cartier-reflection}, to identify
 $\beta^*\chi^X_B$ with $\chi^Y_B$;
\item in Step~5 of the proof of \Cref{thm:bj-homogeneous-reflection}, to
 obtain \eqref{eq:bj-reflection-global-identity};
\item in Step~7 of the same proof, to obtain $\psi=\beta^*u$;
\item in the uniqueness clause of \Cref{thm:bj-carrier-reflection}.
\end{enumerate}
There are two further occurrences of the one-sided comparison
\cref{itm:ch39-det:compare}, both with a psh function on the small side:
in Step~3 of the proof of \Cref{thm:bj-homogeneous-reflection}, to convert
the ordering of divisorial traces into an ordering of the functions
$\chi_j$; and in \Cref{cor:ch39-hom-structure}\,\cref{itm:ch39-hs:order}.
No occurrence compares two merely upper-semicontinuous functions, and
\Cref{ex:ch39-usc-not-determined-birational} shows why none may.
\end{remark}

\section{Legendre slices and carriers}
\label{sec:ch39-legendre}

\subsection{The slices, recalled}
\label{sec:ch39-slices}

For $t>0$, the scaling action of \eqref{eq:function-scaling-action} on
functions is
\[
                   (t\mathbin\cdot u)(v):=t\,u(t^{-1}v),
\tag{39.32}\label{eq:bj-psh-scaling}
\]
and, for $\lambda\in\mathbf R$, the homogeneous Legendre slice of $u$ is
\[
 \widehat u_\lambda
 :=\inf_{t>0}\{t\mathbin\cdot u-t\lambda\},
\tag{39.33}\label{eq:bj-legendre-slice}
\]
which is the largest homogeneous function $\psi$ with
$\psi+\lambda\le u$, by
\Cref{lem:ch35-slice-basic}\,(1).  If $u\in\PSH_X(\theta)$ and
$b:=\sup u=u(v_{\triv})$, then $\widehat u_\lambda$ is homogeneous
$\theta$-psh for every $\lambda\le b$, including the endpoint
$\lambda=b$, and it is identically $-\infty$ for $\lambda>b$; this is
\Cref{prop:ch35-slice-trichotomy}.  Moreover
\[
                   u=\sup_{\lambda\le b}
                       (\widehat u_\lambda+\lambda)
\tag{39.34}\label{eq:bj-legendre-reconstruction}
\]
pointwise on all of $X^{\An}$, by
\Cref{thm:homogeneous-legendre-decomposition}\,(3); and
$\lambda\mapsto\widehat u_\lambda(v)$ is nonincreasing and concave, by
\Cref{lem:ch35-slice-basic}\,(6).  All of this is unconditional and uses no
envelope property.

\begin{definition}[Homogeneous carrier]
\label{def:bj-homogeneous-carrier}
A function
\[
              p\in\PSH_{\mathrm{hom},Y}(\beta^*\theta)
\]
is a \emph{carrier} for $u\in\PSH_Y(\beta^*\theta)$ if
\[
                              u\le p+C
\tag{39.35}\label{eq:bj-carrier-domination}
\]
for some finite $C\in\mathbf R$.
\end{definition}

The constant belongs to the inequality, not to $p$: adding a nonzero
constant to $p$ would destroy homogeneity, by
\Cref{lem:ch35-homogeneity-reformulations}.  This is the case
$\eta=\theta$ of the general notion \Cref{def:ch35-carrier}, in which the
class of the carrier is allowed to differ; the specialization is what makes
homogeneous reflection applicable to $p$ itself, and
\Cref{cor:ch39-carrier-other-class} records the extent to which the class
may be relaxed.

Inequality \eqref{eq:bj-carrier-domination} is an inequality of functions
into $[-\infty,+\infty)$ and is meaningful at proper-support points: where
$p=-\infty$ it forces $u=-\infty$.  It is never rewritten as
$u-p\le C$, which would be undefined there.

The only consequence of \eqref{eq:bj-carrier-domination} used below is the
tail estimate of \Cref{lem:ch35-carrier-tail}: for $\lambda\le C$,
$\widehat u_\lambda\le p$, and hence, for
$\lambda\le\min\{C,b\}$,
$\widehat u_\lambda+\lambda\le p+\lambda$.  Its role is to truncate the
unbounded parameter interval in
\eqref{eq:bj-legendre-reconstruction} to a compact one; see
\Cref{rem:ch35-carrier-purpose}.

\subsection{The compact window}
\label{sec:ch39-window}

Truncating the interval is not by itself enough.  Even over a compact
parameter interval, \eqref{eq:bj-legendre-reconstruction} presents $u$ as
a \emph{supremum}, and \Cref{warn:ch39-firewall} forbids the use of any
closure property under suprema.  The following lemma converts an increasing
family of finite maxima into a decreasing net, at the cost of a summable
correction.  It is the technical heart of carrier reflection, and it is
purely a statement about a monotone one-parameter family of homogeneous psh
functions; no birational geometry enters.

\begin{lemma}
\label{lem:bj-compact-legendre-window}
Let
\[
 q_\lambda\in\PSH_{\mathrm{hom},X}(\theta),
 \qquad -M\le\lambda\le b,
\]
be decreasing as $\lambda$ increases, with $-M\le b$ real, and let
$r\in\PSH_X(\theta)$.  Then
\[
 \max\left\{r,
   \sup_{-M\le\lambda\le b}(q_\lambda+\lambda)\right\}
 \in\PSH_X(\theta).
\tag{39.36}\label{eq:bj-compact-legendre-window}
\]
\end{lemma}

\begin{proof}
Write $h_n:=(b+M)2^{-n}$ and let
\[
 \Gamma_n:=\{-M+kh_n\ :\ 0\le k\le2^n\}\subset[-M,b]
\]
be the dyadic grid of mesh $h_n$.  The grids are nested,
$\Gamma_n\subset\Gamma_{n+1}$, each contains both endpoints, and
$h_{n+1}=h_n/2$.  If $b=-M$ the interval is a point and
\eqref{eq:bj-compact-legendre-window} is the finite maximum
$\max\{r,q_{-M}-M\}$, which is $\theta$-psh by
\Cref{lem:ch39-closure-any-class}\,(1)--(2); so assume $b>-M$, whence
$h_n>0$.

Put
\[
 G_n:=\max\left\{r,\ \max_{a\in\Gamma_n}(q_a+a)\right\}.
\]
Each $q_a+a$ lies in $\PSH_X(\theta)$ by
\Cref{lem:ch39-closure-any-class}\,(1), and $\Gamma_n$ is finite, so
$G_n\in\PSH_X(\theta)$ by
\Cref{lem:ch39-closure-any-class}\,(2).

\emph{The grid error.}  Let $a\in\Gamma_{n+1}$ and let
$a^-:=\max\{x\in\Gamma_n:x\le a\}$, which exists because
$-M\in\Gamma_n$ and satisfies $0\le a-a^-\le h_n$.  Since
$\lambda\mapsto q_\lambda$ is decreasing and $a\ge a^-$,
\[
                 q_a+a\le q_{a^-}+a^-+h_n\le G_n+h_n,
\]
an inequality in $[-\infty,+\infty)$ in which only finite constants are
added.  Also $r\le G_n\le G_n+h_n$.  Taking the maximum over
$a\in\Gamma_{n+1}$,
\[
                         G_n\le G_{n+1}\le G_n+h_n,
\tag{39.37}\label{eq:bj-grid-error}
\]
the first inequality because $\Gamma_n\subseteq\Gamma_{n+1}$.

\emph{The correction.}  Set $c_n:=\sum_{k\ge n}h_k=2h_n<\infty$, so that
$c_n=h_n+c_{n+1}$.  By \eqref{eq:bj-grid-error},
\[
                         G_{n+1}+c_{n+1}\le G_n+h_n+c_{n+1}=G_n+c_n .
\]
Thus $(G_n+c_n)_{n\ge0}$ is a decreasing sequence in $\PSH_X(\theta)$,
again by \Cref{lem:ch39-closure-any-class}\,(1).

\emph{The limit.}  Since $c_n\to0$ and $(G_n)$ is nondecreasing,
\[
 \lim_n(G_n+c_n)=\lim_nG_n=\sup_nG_n
 =\max\left\{r,\ \sup_{a\in\Gamma_\infty}(q_a+a)\right\},
 \qquad
 \Gamma_\infty:=\bigcup_n\Gamma_n,
\]
all limits being pointwise in $\Rbar$.  We claim that this equals the
right-hand side of \eqref{eq:bj-compact-legendre-window}.  One inequality is
clear, $\Gamma_\infty\subseteq[-M,b]$.  For the other, fix
$\lambda\in[-M,b]$ and $v\in X^{\An}$, and let
$a_n:=\max\{x\in\Gamma_n:x\le\lambda\}$, so that
$\lambda-h_n\le a_n\le\lambda$.  Monotonicity in $\lambda$ gives
$q_{a_n}\ge q_\lambda$, whence
\[
 q_{a_n}(v)+a_n\ \ge\ q_\lambda(v)+\lambda-h_n .
\]
If $q_\lambda(v)=-\infty$ the desired inequality
$\sup_{\Gamma_\infty}(q_a+a)(v)\ge q_\lambda(v)+\lambda$ is trivial;
otherwise let $n\to\infty$.  This proves the claim.

\emph{Generic finiteness.}  At $v_{\triv}$ one has
$q_a(v_{\triv})=0$ for every $a$, by
\Cref{prop:ch35-homogeneous-normalization}, and $r(v_{\triv})=\sup r$ is
finite by \Cref{prop:ch32-maximum-principle}\,\cref{itm:ch32-mp:gf}.  Hence
$G_n(v_{\triv})=\max\{r(v_{\triv}),b\}$ for every $n$, and the limit is
finite there.  Since $X$ is integral, the limit is generically finite.

Now PSH2, that is \Cref{lem:ch39-closure-any-class}\,(4) applied to the
decreasing sequence $(G_n+c_n)$ in the fixed class $\theta$, gives
\eqref{eq:bj-compact-legendre-window}.  Only finite maxima, additions of
finite constants and one decreasing limit were used, so the argument is
literally valid at points where all entries are $-\infty$.
\end{proof}

\begin{remark}[Why the correction is needed, and why it is summable]
\label{rem:ch39-why-correction}
The family $(G_n)$ increases, and an increasing limit of psh functions is
exactly what is unavailable before the envelope theorem
(\Cref{warn:ch39-firewall}).  The device of
\eqref{eq:bj-grid-error} is to add a decreasing sequence of constants large
enough to reverse the monotonicity: since the grid refines by halving, the
error committed at stage $n$ is at most $h_n$, and
$\sum_nh_n<\infty$, so the constants $c_n$ can absorb all future errors
at once.  Two features of the hypotheses are used and neither can be
dropped.  The parameter interval must be \emph{compact}, or the mesh would
not be finite; this is what the carrier supplies, through
\Cref{lem:ch35-carrier-tail}.  And the family $\lambda\mapsto q_\lambda$
must be \emph{monotone}, or the estimate $q_a\le q_{a^-}$ would fail and
the grid error would not be controlled by $h_n$; monotonicity is
\Cref{lem:ch35-slice-basic}\,(2).  Continuity of
$\lambda\mapsto q_\lambda$, which does hold on divisorial points by
\Cref{lem:ch35-slice-continuity}, is \emph{not} used, and this is
deliberate: continuity there is proved only on $X^{\divis}$, whereas the
conclusion is needed on all of $X^{\An}$.
\end{remark}

\begin{example}[An exact grid, and the harmless correction]
\label{ex:ch39-grid-correction-needed}
Take $X=\mathbf P^1_k$, $\theta=c_1(\cO(1))$, $p=[0:1]$, and
$q_\lambda:=\max\{0,1+\lambda\}\log|\mathfrak m_p|$ for
$-1\le\lambda\le0$, which by \Cref{ex:ch35-legendre-pair-p1} is the family
of Legendre slices of $u=\max\{\log|x|,\log|y|-1\}$; it is a decreasing
family in $\PSH_{\mathrm{hom},X}(\theta)$.  With $M=1$, $b=0$ and
$r=q_0-1$, the functions $G_n$ of the proof are the finite maxima
\[
 G_n=\max\Bigl\{\log|\mathfrak m_p|-1,\
 \max_{a\in\Gamma_n}\bigl((1+a)\log|\mathfrak m_p|+a\bigr)\Bigr\},
\]
each a Fubini--Study function.  On the ray through $p$, writing
$s=-\log|\mathfrak m_p|(v)\ge0$, one has
$G_n(v)=\max\{-s-1,\max_{a\in\Gamma_n}(a-(1+a)s)\}$.  Here the grid is
\emph{exact} at every stage: $a\mapsto a-(1+a)s=a(1-s)-s$ is affine in
$a$, and every $\Gamma_n$ contains both endpoints $-1$ and $0$ of
the window, so its maximum over $\Gamma_n$ equals its maximum over
$[-1,0]$, namely $-s$ if $s\le1$ and $-1$ if $s\ge1$.  Hence
\[
 G_n\equiv\max\{-s-1,\ \max\{-s,-1\}\}=\max\{-s,-1\}=u
\]
for every $n$: the sequence $(G_n)$ is constant, the grid error
\eqref{eq:bj-grid-error} is zero, and $G_n$ is already visibly psh, being
a Fubini--Study function.  A strictly increasing $(G_n)$ requires
$\lambda\mapsto q_\lambda(v)$ to be strictly concave at some $v$;
a family that is affine in $\lambda$, as this one is, is invisible to
refinement.  What the example does show is that the correction costs
nothing: $G_n+c_n$ with $c_n=2^{1-n}$ decreases, and it decreases to
$u$, so PSH2 applies to it and returns the right answer.  At $s=0$ one
has $G_n\equiv0$ and $c_n\downarrow0$, so the corrected sequence
decreases to the correct value $0$.
\end{example}

\subsection{Carrier reflection}
\label{sec:ch39-carrier-reflection}

\begin{theorem}
\label{thm:bj-carrier-reflection}
Suppose
\[
 u\in\PSH_Y(\beta^*\theta),\qquad
 p\in\PSH_{\mathrm{hom},Y}(\beta^*\theta),\qquad
 u\le p+C
\tag{39.38}\label{eq:bj-carrier-reflection-hypotheses}
\]
for some $C\in\mathbf R$.  Then there is a unique
$v\in\PSH_X(\theta)$ satisfying $\beta^*v=u$.
\end{theorem}

\begin{proof}
Put $b:=\sup_{Y^{\An}}u$.  This is a real number: it equals
$u(v_{Y,\triv})$ by
\Cref{prop:ch32-maximum-principle}\,\cref{itm:ch32-mp:sup}, it is
$>-\infty$ by generic finiteness in \Cref{def:bj-psh}, and it is
$<+\infty$ because a psh function is bounded above by any member of an
approximating net, which is upper semicontinuous on the compact space
$Y^{\An}$.

\medskip
\noindent\emph{Step 1: control the Legendre tail.}
For $\lambda\le C$, homogeneity of $p$ and carrier domination give
\[
 (t\mathbin\cdot u-t\lambda)(v)
 \le p(v)+t(C-\lambda)
\tag{39.39}\label{eq:bj-carrier-slice-estimate}
\]
for every $t>0$.  Indeed $t\cdot u\le t\cdot(p+C)=t\cdot p+tC=p+tC$ by
\Cref{lem:ch35-action-basic}\,(2)--(3) and homogeneity of $p$; subtracting
the finite number $t\lambda$ is legitimate by ledger entry~(2) of
\Cref{sec:extended-arithmetic}.  Taking the infimum over $t>0$ and then
letting $t\downarrow0$ yields
\[
                         \widehat u_\lambda\le p
                         \qquad(\lambda\le C).
\tag{39.40}\label{eq:bj-slice-below-carrier}
\]
If $p(v)=-\infty$, carrier domination forces
$u(t^{-1}v)=-\infty$ for every $t$, so both sides of
\eqref{eq:bj-carrier-slice-estimate} are $-\infty$ and no infinite
subtraction occurs.  This is \Cref{lem:ch35-carrier-tail}, restated here
because the constant $C$ is the one that fixes the window.

Choose $M\in\mathbf N$ with $-M\le\min\{C,b\}$.  For $\lambda<-M$, the
inequality \eqref{eq:bj-slice-below-carrier} applies, since
$\lambda<-M\le C$, and gives
\[
                         \widehat u_\lambda+\lambda\le p-M .
\]
Combining with \eqref{eq:bj-legendre-reconstruction}, and splitting the
parameter range at $-M$,
\[
 \max\{u,p-M\}
 =\max\left\{p-M,
   \sup_{-M\le\lambda\le b}
        (\widehat u_\lambda+\lambda)\right\}.
\tag{39.41}\label{eq:bj-carrier-window}
\]
The inequality $\ge$ holds because each competitor on the right is
$\le\max\{u,p-M\}$, using $\widehat u_\lambda+\lambda\le u$; the
inequality $\le$ holds because
$u=\sup_{\lambda\le b}(\widehat u_\lambda+\lambda)$ and the part of the
supremum with $\lambda<-M$ is dominated by $p-M$.  On the locus
$\{p=-\infty\}$ carrier domination forces every displayed term to equal
$-\infty$, so \eqref{eq:bj-carrier-window} holds there as the identity
$-\infty=-\infty$.

\medskip
\noindent\emph{Step 2: reflect the compact window.}
By \Cref{prop:ch35-slice-trichotomy}\,(2), each
$\widehat u_\lambda$ with $\lambda\le b$ lies in
$\PSH_{\mathrm{hom},Y}(\beta^*\theta)$.  By
\Cref{thm:bj-homogeneous-reflection} there are unique
\[
                  q,\ q_\lambda\ \in
                  \PSH_{\mathrm{hom},X}(\theta)
                  \qquad(-M\le\lambda\le b)
\]
with
\[
                  \beta^*q=p,\qquad
                  \beta^*q_\lambda=\widehat u_\lambda .
\]
The family $(q_\lambda)$ is decreasing in $\lambda$: the slices
$\widehat u_\lambda$ are, by \Cref{lem:ch35-slice-basic}\,(2), and
$\beta^*$ reflects the order by
\Cref{thm:ch39-pullback}\,\cref{itm:ch39-pb:order}.  Apply
\Cref{lem:bj-compact-legendre-window} with $r:=q-M\in\PSH_X(\theta)$,
which is $\theta$-psh by \Cref{lem:ch39-closure-any-class}\,(1), and set
\[
 v_M:=\max\left\{q-M,
   \sup_{-M\le\lambda\le b}(q_\lambda+\lambda)\right\}
 \in\PSH_X(\theta).
\tag{39.42}\label{eq:bj-reflected-window}
\]
Pullback commutes with finite maxima and with arbitrary suprema, by
\Cref{thm:ch39-pullback}\,\cref{itm:ch39-pb:ops}; applying $\beta^*$ to
\eqref{eq:bj-reflected-window} and comparing with
\eqref{eq:bj-carrier-window},
\[
                         \beta^*v_M=\max\{u,p-M\}.
\tag{39.43}\label{eq:bj-reflected-window-pullback}
\]

\medskip
\noindent\emph{Step 3: remove the truncation.}
Let $M$ range over the integers $\ge M_0$, where $M_0$ is the least
integer with $-M_0\le\min\{C,b\}$.  The right-hand side of
\eqref{eq:bj-reflected-window-pullback} decreases pointwise as $M$
increases, and
\[
 \lim_{M\to\infty}\max\{u,p-M\}=u
\]
pointwise on $Y^{\An}$: where $p$ is finite this is clear, and where
$p=-\infty$ carrier domination gives $u=-\infty$ and both sides equal
$-\infty$.  By \Cref{thm:ch39-pullback}\,\cref{itm:ch39-pb:order} the
sequence $(v_M)$ is therefore decreasing.  Moreover, by
\Cref{thm:ch39-pullback}\,\cref{itm:ch39-pb:sup} and
$p(v_{\triv})=0$,
\[
                         v_M(v_{\triv})=\max\{b,-M\}=b
\tag{39.44}\label{eq:bj-reflected-trivial-value}
\]
for every $M\ge M_0\ge-b$.  Hence the pointwise limit
$v:=\lim_Mv_M$ is finite at $v_{\triv}$, so generically finite, and
PSH2 --- \Cref{lem:ch39-closure-any-class}\,(4), with the single fixed class
$\theta$ --- gives $v\in\PSH_X(\theta)$.  Passing to the limit in
\eqref{eq:bj-reflected-window-pullback} gives $\beta^*v=u$.

\medskip
\noindent\emph{Uniqueness.}  If $\beta^*v_1=\beta^*v_2=u$ with
$v_1,v_2\in\PSH_X(\theta)$, then $v_1=v_2$ by
\Cref{thm:ch39-pullback}\,\cref{itm:ch39-pb:order}; alternatively, the two
agree on $X^{\divis}=Y^{\divis}$ and
\Cref{prop:ch39-determination-package}\,\cref{itm:ch39-det:equal} applies.
\end{proof}

\begin{corollary}
\label{cor:ch39-carrier-other-class}
Let $u\in\PSH_Y(\beta^*\theta)$ and suppose there are
$\eta\in N^1(X)_{\mathbf R}$ with $\theta-\eta$ nef,
$p\in\PSH_{\mathrm{hom},Y}(\beta^*\eta)$ and $C\in\mathbf R$ with
$u\le p+C$.  Then $u=\beta^*v$ for a unique $v\in\PSH_X(\theta)$.
\end{corollary}

\begin{proof}
By \Cref{lem:ch39-closure-any-class}\,(5) applied on $Y$ with the nef
class $\beta^*(\theta-\eta)$ --- nef by
\Cref{lem:ch38-numerical-birational}\,\cref{itm:ch38-num:inj} --- one has
$p\in\PSH_Y(\beta^*\theta)$.  It is still homogeneous, homogeneity being a
property of the function and not of the class, so
$p\in\PSH_{\mathrm{hom},Y}(\beta^*\theta)$ and
\Cref{thm:bj-carrier-reflection} applies verbatim.
\end{proof}

\begin{remark}[Exactly where the carrier is used]
\label{rem:ch39-carrier-where}
The carrier enters the proof of \Cref{thm:bj-carrier-reflection} at one
point only: it makes the parameter window compact, in
\eqref{eq:bj-carrier-window}.  Everything after that is
\Cref{lem:bj-compact-legendre-window} and homogeneous reflection.  In
particular the carrier is \emph{not} used to construct $v$, and $v$ is
not obtained as an envelope; it is obtained as a decreasing limit of the
explicitly constructed functions \eqref{eq:bj-reflected-window}.  A bounded
$u$ has the tautological carrier $0$ whenever $\beta^*\theta$ is nef
(\Cref{ex:ch35-bounded-carrier}), which is why bounded psh functions reflect
without any extra input; the content of the theorem is entirely in the
unbounded case, and \Cref{ex:ch35-divisorial-carrier} shows that a
plausible-looking candidate carrier can fail to dominate the slices.
\end{remark}

\begin{corollary}
\label{cor:ch39-nef-reflection}
Suppose $\theta\in\Nef(X)$, equivalently $\beta^*\theta\in\Nef(Y)$.
Then
\[
 \beta^*:\PSH_X(\theta)\longrightarrow\PSH_Y(\beta^*\theta)
\]
is a bijection, with no further hypothesis.
\end{corollary}

\begin{proof}
The two nefness statements are equivalent: $\beta^*\Nef(X)\subseteq\Nef(Y)$
by \Cref{lem:ch38-numerical-birational}\,\cref{itm:ch38-num:inj}, and
conversely $\beta^*\theta$ nef implies $\theta$ nef by the argument
proving \Cref{lem:ch38-nef-total-intrinsic}\,\cref{itm:ch38-nt:indep}, which
replaces a strict transform by an arbitrary integral curve dominating the
given curve on $X$.  One cannot simply write
$\theta\cdot C=\beta^*\theta\cdot\widetilde C$ for every integral
$C\subset X$: a curve contained in $\beta(\operatorname{Exc}\beta)$, which
exists as soon as $\dim X\geq3$, has no strict transform at all.

Assume them.  By \Cref{lem:ch34-psh-nonempty} the constant $0$ is
$\beta^*\theta$-psh, and it is homogeneous, so
$0\in\PSH_{\mathrm{hom},Y}(\beta^*\theta)$.  For any
$u\in\PSH_Y(\beta^*\theta)$, the number $b=\sup u$ is finite and
$u\le0+b$; thus $0$ is a carrier for $u$ in the sense of
\Cref{def:bj-homogeneous-carrier}, with $C=b$.  Now
\Cref{thm:bj-carrier-reflection} gives surjectivity, and injectivity is
\Cref{thm:ch39-pullback}\,\cref{itm:ch39-pb:order}.
\end{proof}

\begin{remark}[Where the carrier hypothesis actually bites]
\label{rem:ch39-carrier-bites}
By \Cref{cor:ch39-nef-reflection}, the carrier hypothesis of
\Cref{thm:bj-carrier-reflection} is vacuous when $\theta$ is nef.  Its
content is therefore in the non-nef range where
$\PSH_X(\theta)\ne\emptyset$; such a class is pseudoeffective and contains
no bounded psh function by \Cref{lem:ch34-psh-nonempty}.  This is one more reason
why the correction of \Cref{def:bj-psh} matters for this chapter: under the
older convention restricted to bounded PL approximants in a fixed class,
$\PSH_X(\theta)$ would have been empty outside the nef cone and
\Cref{thm:bj-carrier-reflection} would have been a restatement of
\Cref{cor:ch39-nef-reflection}.  The generically finite, varying-class
convention admits non-nef pseudoeffective examples, and with it the
canonical carrier of \Cref{sec:ch39-descent} becomes the working hypothesis
rather than a formality.
\end{remark}

\subsection{A Legendre carrier pair, computed}
\label{sec:ch39-carrier-example}

\begin{example}[Carrier, slices, window and reflection, in numbers]
\label{ex:ch39-carrier-pair}
\index{carrier!computed}
Keep $X=\mathbf P^2_k$, $p\in X$, $\beta:Y=\operatorname{Bl}_pX\to X$
and $E$, and set
\[
 \theta:=3\,c_1(\cO_X(1)),\qquad
 \beta^*\theta=3H,\qquad
 \psi_Y:=\log|\mathfrak m_p\cO_Y|=\log|\cO_Y(-E)| .
\]
Write $s(w):=-\psi_Y(w)\in[0,+\infty]$, so that $s(t\,\ord_E)=t$,
$s\equiv0$ on every other divisorial ray and at $v_{Y,\triv}$, and
$s\equiv+\infty$ on the fibre $E^{\An}$ of \Cref{ex:ch39-fibre-plane}.
Define
\[
 u:=\max\{3\psi_Y,\ \psi_Y-1\}
 \qquad\text{so that}\qquad
 u=\max\{-3s,\ -s-1\}.
\tag{39.45}\label{eq:ch39-carrier-u}
\]

\emph{$u$ is $3H$-psh.}  Indeed $3\psi_Y\in\PSH_{\mathrm{hom},Y}(3H)$
by \Cref{prop:ch35-homogeneous-normalization}, and
$\psi_Y\in\PSH_{\mathrm{hom},Y}(H)\subseteq\PSH_Y(3H)$ by
\Cref{lem:ch39-closure-any-class}\,(5), $2H$ being nef; adding the
constant $-1$ and taking the maximum stays in $\PSH_Y(3H)$ by
\Cref{lem:ch39-closure-any-class}\,(1)--(2).  Its supremum is
$b=u(v_{Y,\triv})=0$, and $u$ is unbounded below along the exceptional
ray.

\emph{The carrier.}  To avoid a clash with the blown-up point $p$, we
write $\varrho$ for the carrier of \Cref{def:bj-homogeneous-carrier} in
this example.  Take $\varrho:=\psi_Y$, which lies in
$\PSH_{\mathrm{hom},Y}(3H)=\PSH_{\mathrm{hom},Y}(\beta^*\theta)$ by the
same monotonicity.  Then $u\le\varrho$, that is,
\eqref{eq:bj-carrier-domination} holds with $C=0$: on the exceptional ray
one checks $\max\{-3s,-s-1\}\le-s$ separately on $s\le\frac12$, where
the maximum is $-3s\le-s$, and on $s\ge\frac12$, where it is
$-s-1\le-s$; off that ray both sides are $0$; and on $E^{\An}$ both
are $-\infty$.

\emph{The slices.}  For $\lambda\le0$ and $w$ on the exceptional ray
with $s=s(w)$,
\[
 \widehat u_\lambda(w)
 =\inf_{t>0}\bigl\{\max\{-3s,\,-s-t\}-t\lambda\bigr\}.
\]
The function of $t$ inside the infimum equals $-s-t(1+\lambda)$ for
$t\le2s$ and $-3s-t\lambda$ for $t\ge2s$.  If $-1\le\lambda\le0$
both pieces are minimized at $t=2s$, with common value
$-3s-2s\lambda$; if $\lambda\le-1$ the first piece increases in $t$
and the infimum is approached as $t\downarrow0$, with value $-s$.  Off
the exceptional ray the infimum is $\inf_{t>0}(-t\lambda)=0$ for
$\lambda\le0$.  Hence
\[
 \widehat u_\lambda=c(\lambda)\,\psi_Y,
 \qquad
 c(\lambda):=\max\{1,\ 3+2\lambda\}
 \qquad(\lambda\le0),
\tag{39.46}\label{eq:ch39-carrier-slices}
\]
and $\widehat u_\lambda\equiv-\infty$ for $\lambda>0$, in accordance
with \Cref{prop:ch35-slice-trichotomy}.  Since $1\le c(\lambda)\le3$, each
slice lies in $\PSH_{\mathrm{hom},Y}(c(\lambda)H)\subseteq
\PSH_{\mathrm{hom},Y}(3H)$; the map $\lambda\mapsto c(\lambda)$ is
nondecreasing, so $\lambda\mapsto\widehat u_\lambda$ is nonincreasing, and
$-c(\lambda)$ is concave, as
\Cref{lem:ch35-slice-basic}\,(6) requires.  The carrier bound
\eqref{eq:bj-slice-below-carrier} reads $c(\lambda)\ge1$, which is visible
in \eqref{eq:ch39-carrier-slices}.

\emph{Reconstruction.}  On the exceptional ray,
\[
 \sup_{\lambda\le0}\bigl\{\widehat u_\lambda+\lambda\bigr\}(w)
 =\sup_{\lambda\le0}\bigl\{-c(\lambda)s+\lambda\bigr\}
 =\max\Bigl\{\sup_{\lambda\le-1}(-s+\lambda),\
 \sup_{-1\le\lambda\le0}\bigl(-3s+\lambda(1-2s)\bigr)\Bigr\}.
\]
The first supremum is $-s-1$, attained at $\lambda=-1$.  The second is a
supremum of an affine function of $\lambda$ with slope $1-2s$, attained
at $\lambda=0$ with value $-3s$ if $s\le\frac12$ and at
$\lambda=-1$ with value $-s-1$ if $s\ge\frac12$.  The total is
$\max\{-3s,-s-1\}=u(w)$, confirming
\eqref{eq:bj-legendre-reconstruction}.

\emph{The window.}  Take $M=1$, which satisfies
$-M\le\min\{C,b\}=0$.  For $\lambda<-1$ one has
$\widehat u_\lambda+\lambda=\psi_Y+\lambda\le\psi_Y-1=\varrho-M$, so
\eqref{eq:bj-carrier-window} holds:
\[
 \max\{u,\ \psi_Y-1\}=u
 =\max\Bigl\{\psi_Y-1,\ \sup_{-1\le\lambda\le0}
 \bigl(\widehat u_\lambda+\lambda\bigr)\Bigr\},
\]
both sides being $\max\{-3s,-s-1\}$ by the previous paragraph.

\emph{Reflection.}  The homogeneous reflections are
\[
 q=\log|\mathfrak m_p|\in\PSH_{\mathrm{hom},X}(\theta),
 \qquad
 q_\lambda=c(\lambda)\log|\mathfrak m_p|,
\]
with $\beta^*q=\psi_Y$ and $\beta^*q_\lambda=\widehat u_\lambda$, by
\Cref{ex:ch39-blowup-point} and \Cref{thm:bj-homogeneous-reflection}.
Formula \eqref{eq:bj-reflected-window} gives
\[
 \begin{aligned}
 v_1&=\max\Bigl\{\log|\mathfrak m_p|-1,\
 \sup_{-1\le\lambda\le0}
 \bigl(c(\lambda)\log|\mathfrak m_p|+\lambda\bigr)\Bigr\}\\
 &=\max\bigl\{3\log|\mathfrak m_p|,\ \log|\mathfrak m_p|-1\bigr\},
 \end{aligned}
\]
and $v_M=v_1$ for every $M\ge1$, because for $\lambda\le-1$ the
competitor $q_\lambda+\lambda=\log|\mathfrak m_p|+\lambda$ is dominated by
the term $\log|\mathfrak m_p|-1$ already present.  So the truncation
removal of Step~3 is instantaneous here, and
\[
 v=\max\bigl\{3\log|\mathfrak m_p|,\ \log|\mathfrak m_p|-1\bigr\}
 \in\PSH_X(\theta),
 \qquad
 \beta^*v=u .
\]

\emph{The proper-support values.}  At $v_{p,\triv}\in X^{\An}$ one has
$\log|\mathfrak m_p|=-\infty$, hence $v(v_{p,\triv})=-\infty$; and
$u\equiv-\infty$ on the fibre $E^{\An}=(\beta^{\An})^{-1}(v_{p,\triv})$.
The two are consistent, as \Cref{cor:ch39-fibre-constant} requires.  Note
that the identity $\beta^*v=u$ is an identity of extended-real functions
including at these points; it was obtained without ever writing a difference
of two infinite values, and the intermediate objects
$v_M$ were finite wherever $\log|\mathfrak m_p|$ is.
\end{example}

\begin{remark}[A carrier that fails]
\label{rem:ch39-carrier-fails}
It is instructive to see the hypothesis fail.  Replace
\eqref{eq:ch39-carrier-u} by $u':=\max\{\psi_Y,\ -1\}$, a bounded
$H$-psh function, and ask whether $\psi_Y$ is a carrier for it.  It is
not: $u'\le\psi_Y+C$ would require $-1\le-s+C$ for every $s$, which
fails for $s>C+1$.  The slices bear this out: by the computation of
\Cref{ex:ch35-legendre-pair-p1}, transported through
\Cref{ex:ch39-blowup-point}, one finds
$\widehat{u'}_\lambda=\max\{0,1+\lambda\}\psi_Y$, whose coefficient is
$<1$ for $\lambda<0$, so that
$\widehat{u'}_\lambda\not\le\psi_Y$ off the vertex, contradicting
\eqref{eq:bj-slice-below-carrier}.  Of course $u'$ does have the
tautological carrier $0$, $H$ being nef, and
\Cref{cor:ch39-nef-reflection} applies; the point is only that not every
homogeneous psh function dominates a given $u$, and that
\Cref{lem:ch35-carrier-tail} is sensitive to which one is chosen.
\end{remark}

\section{The canonical carrier and descent of the envelope property}
\label{sec:ch39-descent}

Up to this point no envelope hypothesis has been made.  In this section one
is made, explicitly and on the source $Y$, and its two uses are kept
separate.  First it manufactures a single homogeneous function carrying
every $\beta^*\theta$-psh function at once.  Then its increasing-net form
produces the upstairs psh limit that is reflected downstairs.  The output,
\Cref{thm:bj-birational-ep-descent}, carries no envelope hypothesis on
$X$.

\subsection{The canonical carrier}
\label{sec:ch39-canonical-carrier}

\begin{lemma}
\label{lem:bj-canonical-carrier}
Let $Y$ be normal projective and let
$\eta\in N^1(Y)_{\mathbf R}$.  Assume
$\PSH_Y(\eta)\ne\varnothing$ and
$\operatorname{EP}_Y(\eta)$.  Then
\[
 P_Y:=\left(\sup\{w\in\PSH_Y(\eta):\sup w\le0\}\right)^*
\tag{39.47}\label{eq:bj-canonical-carrier}
\]
is homogeneous psh.  Every $u\in\PSH_Y(\eta)$ satisfies
\[
                              u\le P_Y+\sup u.
\tag{39.48}\label{eq:bj-canonical-carrier-bound}
\]
\end{lemma}

\begin{proof}
Write $\mathcal F:=\{w\in\PSH_Y(\eta):\sup w\le0\}$.

\emph{$\mathcal F$ is nonempty and bounded above.}  Choose
$u\in\PSH_Y(\eta)$.  Its supremum $\sup u$ is a real number, by
\Cref{prop:ch32-maximum-principle}\,\cref{itm:ch32-mp:sup} and generic
finiteness, so $u-\sup u\in\PSH_Y(\eta)$ by
\Cref{lem:ch39-closure-any-class}\,(1) and it has supremum $0$; hence
$u-\sup u\in\mathcal F$.  Every member of $\mathcal F$ is $\le0$.

\emph{$P_Y$ is $\eta$-psh.}  This is the hypothesis
$\operatorname{EP}_Y(\eta)$ of \Cref{def:bj-envelope-property}, applied to
the nonempty bounded-above family $\mathcal F$.  This is the first use of
the source envelope hypothesis; the second is the increasing-net limit in
the proof of \Cref{thm:bj-birational-ep-descent}.

\emph{$P_Y$ is homogeneous.}  Let $t>0$ and let $\sigma_t:Y^{\An}\to
Y^{\An}$, $v\mapsto tv$, which is a homeomorphism by
\Cref{prop:ch32-scaling-action}\,\cref{itm:ch32-sa:cont}.  The scaling action
of \eqref{eq:bj-psh-scaling} is $t\cdot w=t\,(w\circ\sigma_{t^{-1}})$.
It preserves $\PSH_Y(\eta)$, by
\Cref{lem:ch35-action-basic}\,(6) applied to an approximating net together
with \Cref{def:bj-psh}, and it scales suprema:
$\sup(t\cdot w)=t\sup w$, by
\Cref{lem:ch35-action-basic}\,(4).  Hence $w\mapsto t\cdot w$ maps
$\mathcal F$ bijectively onto itself, with inverse
$w\mapsto t^{-1}\cdot w$.  Consequently
\[
 t\cdot\Bigl(\sup_{w\in\mathcal F}w\Bigr)
 =\sup_{w\in\mathcal F}(t\cdot w)=\sup_{w\in\mathcal F}w,
\]
by \Cref{lem:ch35-action-basic}\,(4).  Finally the operations
$f\mapsto f\circ\sigma_{t^{-1}}$ and $f\mapsto tf$ both commute with
upper-semicontinuous regularization --- the first because $\sigma_{t^{-1}}$
is a homeomorphism and \eqref{eq:appE-usc-regularization} is defined by
neighbourhood filters, the second because multiplication by $t>0$ is an
increasing homeomorphism of $\Rbar$, by
\Cref{lem:appE-usc-calculus}.  Therefore
$t\cdot P_Y=(t\cdot\sup\mathcal F)^*=(\sup\mathcal F)^*=P_Y$, which is
homogeneity in the sense of \Cref{def:homogeneous-function}.  Since
$P_Y$ is $\eta$-psh and homogeneous,
$P_Y\in\PSH_{\mathrm{hom},Y}(\eta)$.

\emph{The bound.}  For $u\in\PSH_Y(\eta)$ the function $u-\sup u$ lies
in $\mathcal F$, hence $u-\sup u\le\sup\mathcal F\le P_Y$, which is
\eqref{eq:bj-canonical-carrier-bound} after adding the finite constant
$\sup u$.
\end{proof}

\begin{proposition}
\label{prop:ch39-canonical-carrier-props}
Under the hypotheses of \Cref{lem:bj-canonical-carrier}:
\begin{enumerate}
\item\label{itm:ch39-cc:max} $P_Y$ is the largest element of
 $\PSH_{\mathrm{hom},Y}(\eta)$; in particular
 $P_Y(v_{\triv})=0$ and $P_Y\le0$;
\item\label{itm:ch39-cc:nef} $P_Y=0$ if and only if $\eta$ is nef;
\item\label{itm:ch39-cc:idem} $P_Y$ is a carrier for itself, with
 $C=0$, and $t\cdot P_Y=P_Y$ for every $t>0$;
\item\label{itm:ch39-cc:pull} if $\eta=\beta^*\theta$ then
 $P_Y=\beta^*P_X$ for a unique
 $P_X\in\PSH_{\mathrm{hom},X}(\theta)$, and $P_X$ is the largest element
 of $\PSH_{\mathrm{hom},X}(\theta)$.
\end{enumerate}
\end{proposition}

\begin{proof}
\cref{itm:ch39-cc:max}  Every $\psi\in\PSH_{\mathrm{hom},Y}(\eta)$ has
$\sup\psi=\psi(v_{\triv})=0$ by
\Cref{prop:ch35-homogeneous-normalization}, so $\psi\in\mathcal F$ and
$\psi\le P_Y$; and $P_Y$ itself lies in
$\PSH_{\mathrm{hom},Y}(\eta)$ by \Cref{lem:bj-canonical-carrier}.  Taking
$\psi=u-\sup u$ for any $u$ shows $P_Y(v_{\triv})\ge0$, while
$P_Y\le0$ gives the reverse.

\cref{itm:ch39-cc:nef}  If $\eta$ is nef then $0\in\PSH_Y(\eta)$ by
\Cref{lem:ch34-psh-nonempty} and $0$ is homogeneous, so $0\le P_Y\le0$.
Conversely if $P_Y=0$ then the constant $0$ is $\eta$-psh, hence
bounded and $\eta$-psh, so $\eta$ is nef by
\Cref{lem:ch34-psh-nonempty}.

\cref{itm:ch39-cc:idem} is \eqref{eq:bj-canonical-carrier-bound} with
$u=P_Y$, since $\sup P_Y=0$; homogeneity was proved above.

\cref{itm:ch39-cc:pull} is \Cref{thm:bj-homogeneous-reflection} together with
\Cref{cor:ch39-hom-structure}\,\cref{itm:ch39-hs:order}, which transports the
maximality.
\end{proof}

\begin{warning}{The canonical carrier is an assumption, not a construction}
\label{warn:ch39-canonical-carrier}
Formula \eqref{eq:bj-canonical-carrier} makes sense for every
$\eta$, and the resulting function is always homogeneous, upper
semicontinuous and $\le0$.  What requires
$\operatorname{EP}_Y(\eta)$ is the assertion that it is \emph{psh}.
Without that, $P_Y$ is one more upper-semicontinuous function with the
right divisorial trace, and
\Cref{ex:ch39-two-usc-extensions} shows how little that determines.  The
temptation to regard \eqref{eq:bj-canonical-carrier} as a construction, and
hence \Cref{thm:bj-carrier-reflection} as unconditional, is exactly the
circular step that \Cref{warn:ch39-firewall} forbids: the envelope property
for $X$ is what this chapter is helping to prove.
\end{warning}

\subsection{Descent}
\label{sec:ch39-ep-descent}

\begin{theorem}
\label{thm:bj-birational-ep-descent}
For \eqref{eq:bj-reflection-map},
\[
 \operatorname{EP}_Y(\beta^*\theta)
 \quad\Longrightarrow\quad
 \operatorname{EP}_X(\theta).
\tag{39.49}\label{eq:bj-birational-ep-descent}
\]
If source EP holds and the upstairs psh space is nonempty, pullback is a
bijection
\[
 \beta^*:\PSH_X(\theta)
       \xrightarrow{\ \sim\ }\PSH_Y(\beta^*\theta).
\tag{39.50}\label{eq:bj-ordinary-reflection-under-ep}
\]
\end{theorem}

\begin{proof}
Suppose first that $\PSH_Y(\beta^*\theta)=\varnothing$.  Then
$\PSH_X(\theta)=\varnothing$ as well, because
\Cref{thm:ch39-pullback}\,\cref{itm:ch39-pb:psh} maps the latter into the
former; so $\operatorname{EP}_X(\theta)$ holds vacuously, there being no
nonempty family to test, and \eqref{eq:bj-ordinary-reflection-under-ep} is
an identity of empty sets.  Assume from now on that
$\PSH_Y(\beta^*\theta)\ne\varnothing$.

\emph{The bijection.}  By \Cref{lem:bj-canonical-carrier}, applied on $Y$
with $\eta=\beta^*\theta$, the function $P_Y$ of
\eqref{eq:bj-canonical-carrier} lies in
$\PSH_{\mathrm{hom},Y}(\beta^*\theta)$ and satisfies
$u\le P_Y+\sup u$ for every $u\in\PSH_Y(\beta^*\theta)$.  Thus every
such $u$ has a homogeneous carrier in the sense of
\Cref{def:bj-homogeneous-carrier}, with the finite constant $C=\sup u$,
and \Cref{thm:bj-carrier-reflection} produces a unique
$v\in\PSH_X(\theta)$ with $\beta^*v=u$.  This is surjectivity in
\eqref{eq:bj-ordinary-reflection-under-ep}; injectivity is
\Cref{thm:ch39-pullback}\,\cref{itm:ch39-pb:order}.

\emph{The envelope property downstairs.}  By
\Cref{prop:bj-increasing-net-criterion} it suffices to verify the
increasing-net form: for every bounded-above upward-directed net
$(v_j)_{j\in J}$ in $\PSH_X(\theta)$, the pointwise supremum
$\sup_jv_j$ on $X^{\divis}$ is the divisorial trace of an element of
$\PSH_X(\theta)$.  If $\PSH_X(\theta)=\varnothing$ there is no such net
and the criterion holds vacuously; assume there is one.

Pull back.  By \Cref{thm:ch39-pullback}\,\cref{itm:ch39-pb:psh} and
\cref{itm:ch39-pb:order}, the family $(\beta^*v_j)$ is an upward-directed
net in $\PSH_Y(\beta^*\theta)$, bounded above by the same constant.  The
hypothesis $\operatorname{EP}_Y(\beta^*\theta)$, through the equivalence
\Cref{prop:bj-increasing-net-criterion}, supplies
\[
 U_Y\in\PSH_Y(\beta^*\theta),\qquad
 U_Y(w)=\sup_j\beta^*v_j(w)
 \quad(w\in Y^{\divis}).
\tag{39.51}\label{eq:bj-source-net-limit}
\]
Reflect $U_Y$ using its canonical carrier: by the previous paragraph there
is a unique $U_X\in\PSH_X(\theta)$ with $\beta^*U_X=U_Y$.

It remains to compute the divisorial trace of $U_X$.  The common function
field identifies $X^{\divis}$ and $Y^{\divis}$, and $\beta^{\An}$
restricts to the identity there, by
\Cref{lem:ch39-beta-an}\,\cref{itm:ch39-ban:val}.  Hence for
$v\in X^{\divis}$,
\[
 U_X(v)=\beta^*U_X(v)=U_Y(v)
 =\sup_j\beta^*v_j(v)=\sup_jv_j(v).
\]
So $\sup_jv_j$ on $X^{\divis}$ is the divisorial trace of
$U_X\in\PSH_X(\theta)$, which is the increasing-net criterion; and
\Cref{prop:bj-increasing-net-criterion} then also identifies
$U_X=(\sup_jv_j)^*$.  This proves
$\operatorname{EP}_X(\theta)$.
\end{proof}

\begin{corollary}
\label{cor:ch39-ep-tower}
Let $X$ be normal integral projective and $\theta\in N^1(X)_{\mathbf R}$.
\begin{enumerate}
\item\label{itm:ch39-tow:comp} If $Z\xrightarrow{\ \gamma\ }Y
 \xrightarrow{\ \beta\ }X$ are projective birational morphisms of normal
 integral projective varieties, then
 $\operatorname{EP}_Z((\beta\gamma)^*\theta)$ implies
 $\operatorname{EP}_Y(\beta^*\theta)$ and hence
 $\operatorname{EP}_X(\theta)$.
\item\label{itm:ch39-tow:some} If $\operatorname{EP}_Y(\beta^*\theta)$
 holds for \emph{some} normal projective model $\beta:Y\to X$, then
 $\operatorname{EP}_X(\theta)$ holds.
\item\label{itm:ch39-tow:all} If $\operatorname{EP}_Y(\beta^*\theta)$
 holds for some model $\beta:Y\to X$, it holds for every model dominated
 by $Y$.
\end{enumerate}
\end{corollary}

\begin{proof}
\cref{itm:ch39-tow:comp} is \Cref{thm:bj-birational-ep-descent} applied
twice, to $\gamma$ with the class $\beta^*\theta$ and then to
$\beta$ with the class $\theta$; the composite class matches because
$(\beta\gamma)^*=\gamma^*\beta^*$.  \cref{itm:ch39-tow:some} is the case
$Z=Y$, and \cref{itm:ch39-tow:all} is \cref{itm:ch39-tow:comp} with the
intermediate model in the middle.
\end{proof}

\begin{remark}[The implication is one-way, as proved here]
\label{rem:ch39-one-way}
\Cref{thm:bj-birational-ep-descent} proves
$\operatorname{EP}_Y(\beta^*\theta)\Rightarrow\operatorname{EP}_X(\theta)$
and says nothing about the converse.  The obstruction is visible: to run the
argument in the other direction one would need to reflect a family upstairs
to a family downstairs, that is, to know that every
$u\in\PSH_Y(\beta^*\theta)$ is a pullback; and that is
\eqref{eq:bj-ordinary-reflection-under-ep}, whose proof needs a carrier for
each $u$, that is, needs $\operatorname{EP}_Y$ itself.  After
\Cref{thm:bj-final-envelope} has been proved, both properties hold for all
normal projective $X$ and all real classes, so the converse is true; but
it is true as a consequence of the theorem, not as an input, and using it
here would be circular.  The one-way arrow is drawn as such in
\eqref{eq:bj-ep-reflection-diagram}.
\end{remark}

\begin{remark}[How \Cref{chap:bj-alterations} and
\Cref{chap:bj-finite-descent} consume this section]
\label{rem:ch39-downstream}
\Cref{chap:bj-alterations} factors an alteration $f:Y\to X$ as
$Y\xrightarrow{g}T\xrightarrow{h}X$ with $g$ birational and $h$
finite, and applies \Cref{thm:bj-birational-ep-descent} to $g$ with the
class $h^*\theta$ on $T$; the class $h^*\theta$ is an arbitrary real
class on a normal projective variety, which is why the present chapter is
stated without any positivity hypothesis.
\Cref{chap:bj-finite-descent} uses instead the two ingredients separately:
\Cref{lem:bj-canonical-carrier} to produce a homogeneous carrier on the
source of a finite map, and \Cref{thm:bj-carrier-reflection} to reflect
along a birational modification appearing in the flattening construction.
Neither chapter uses \Cref{thm:bj-homogeneous-reflection} directly; they use
it through \Cref{thm:bj-carrier-reflection}.
\end{remark}

\section{Proper support and the extended-real audit}
\label{sec:ch39-safeguards}

Every function produced in this chapter takes values in
$[-\infty,+\infty)$.  Every psh function is finite on
$X^{\lin}$, hence on $X^{\divis}$, by
\Cref{thm:bj-basic-closure}\,\cref{itm:bj-closure:determination}; however,
it may equal $-\infty$ at proper-support points and, for a decreasing
limit, possibly also at full-support points outside $X^{\lin}$.
The polar-locus formula
\Cref{prop:ch34-generic-finiteness}\,\cref{itm:ch34-gf:polar} applies to an
individual Fubini--Study stage, not automatically to its psh limit.  The
proof-audit list of \Cref{chap:bj-envelope-statement} contains the entries
``values at proper-support semivaluations'' and ``every use of subtraction
involving an infinite value''.  We discharge them here.

\subsection{The audit}
\label{sec:ch39-audit}

\begin{proposition}
\label{prop:ch39-safeguards}
In the proofs of \Cref{thm:bj-homogeneous-reflection},
\Cref{lem:bj-compact-legendre-window},
\Cref{thm:bj-carrier-reflection},
\Cref{lem:bj-canonical-carrier} and
\Cref{thm:bj-birational-ep-descent}, every extended-real operation is one of
the following, and each is legitimate for the reason stated.
\begin{enumerate}
\item\label{itm:ch39-safe:eval}
 \emph{Evaluation of a Cartier b-divisor.}  The numbers $D_j(v)$,
 $v(F)$, $v(G)$ and $E_j(v)$ of
 \Cref{lem:ch39-interpolation}\,\cref{itm:ch39-int:limit} are finite real
 numbers for $v\in X^{\divis}$, by
 \Cref{prop:ch38-evaluation-welldefined}\,\cref{itm:ch38-ev:finite}.  Only
 the limit $D(v)=\sup_jD_j(v)$ may equal $+\infty$, and it is never
 subtracted; the limit \eqref{eq:bj-reflection-Ei-limit} is computed in
 $[0,+\infty]$ as a monotone limit.
\item\label{itm:ch39-safe:sum}
 \emph{The divisorial identity.}  In
 \eqref{eq:bj-reflection-divisorial-identity} both summands lie in
 $[-\infty,0]$, so the sum is defined by ledger entry~(2) of
 \Cref{sec:extended-arithmetic}.  The identity is asserted as an identity of
 sums, never solved for $\psi(v)$ at a point where
 $\log|s_F|(v)=-\infty$; at divisorial points the latter is finite and the
 solution is legitimate.
\item\label{itm:ch39-safe:cancel}
 \emph{The cancellation in Step~7.}  From
 $\psi+\log|s_{\beta^*F}|=\beta^*u+\log|s_{\beta^*F}|$ one concludes
 $\psi=\beta^*u$ only at points of $Y^{\divis}$, where the common
 summand is a real number and addition of it is injective on
 $[-\infty,+\infty)$.  The extension to $Y^{\An}$ is then made by
 \Cref{prop:ch39-determination-package}\,\cref{itm:ch39-det:equal}, not by
 cancelling infinities.
\item\label{itm:ch39-safe:siu}
 \emph{Siu division.}  \Cref{thm:ch35-homogeneous-siu-general} outputs the
 identity $\chi=\log|s_F|+u$ directly.  It is never rearranged as
 $u=\chi-\log|s_F|$; \Cref{warn:ch39-no-subtraction} and
 \Cref{ex:ch39-subtraction-fails} show that no convention would make the
 rearrangement correct.
\item\label{itm:ch39-safe:carrier}
 \emph{The carrier estimate.}  In
 \eqref{eq:bj-carrier-slice-estimate} the quantity subtracted is the real
 number $t\lambda$, by ledger entry~(2).  Where $p(v)=-\infty$, carrier
 domination forces $u(t^{-1}v)=-\infty$ and both sides are $-\infty$; no
 indeterminate difference arises.  The inequality
 \eqref{eq:bj-carrier-domination} is never rewritten as $u-p\le C$.
\item\label{itm:ch39-safe:window}
 \emph{The window lemma.}  The proof of
 \Cref{lem:bj-compact-legendre-window} uses only finite maxima, additions of
 real constants, and one decreasing limit.  All three are defined on
 $[-\infty,+\infty)$ without convention, by
 \Cref{lem:extended-real-safe}, so the argument is literal at points where
 every entry is $-\infty$.
\item\label{itm:ch39-safe:truncation}
 \emph{Removing the truncation.}  The limit
 $\max\{u,p-M\}\downarrow u$ of Step~3 of
 \Cref{thm:bj-carrier-reflection} holds pointwise on all of $Y^{\An}$,
 including on $\{p=-\infty\}$, where both sides are $-\infty$ by
 \eqref{eq:bj-carrier-domination}.  It is a monotone limit in $\Rbar$.
\item\label{itm:ch39-safe:canonical}
 \emph{The canonical carrier.}  The bound
 \eqref{eq:bj-canonical-carrier-bound} is an inequality of functions into
 $[-\infty,0]$ plus a real constant.  The normalization $u-\sup u$
 subtracts the real number $\sup u$, finite by
 \Cref{prop:ch32-maximum-principle}\,\cref{itm:ch32-mp:sup}.
\item\label{itm:ch39-safe:value}
 \emph{Values at proper support.}  The value of a reflected function at a
 point of proper support is never computed by a formula.  It is determined
 by \Cref{prop:ch39-determination-package}\,\cref{itm:ch39-det:equal}, which
 fixes it once the divisorial trace and the psh property are known, and
 \Cref{ex:ch39-two-usc-extensions} shows that the divisorial trace alone
 does not fix it.
\end{enumerate}
\end{proposition}

\begin{proof}
Each item was verified where it occurs; the list is a record, and the
justifications are those given in the corresponding proofs, together with
\Cref{lem:extended-real-safe} and the ledger of
\Cref{sec:extended-arithmetic}.
\end{proof}

\subsection{What fails if a safeguard is dropped}
\label{sec:ch39-safeguard-failures}

\begin{example}[The forbidden subtraction, in numbers]
\label{ex:ch39-subtraction-fails}
Continue \Cref{ex:ch39-blowup-reflection-numbers}: $X=\mathbf P^2$,
$p\in X$, $F=L$ a line through $p$ with equation $\ell$, and
\[
 \chi=\log|\ell\cdot\mathfrak m_p|,\qquad
 \log|s_F|=\log|\ell|,\qquad
 u=\log|\mathfrak m_p| ,
\]
so that $\chi=\log|s_F|+u$ is \eqref{eq:bj-reflection-siu}.  Evaluate at
two points of proper support.

At $v_{p,\triv}$, the trivial semivaluation of the point $p$: both
$\ell$ and every element of $\mathfrak m_p$ vanish at $p$, so
\[
 \chi(v_{p,\triv})=-\infty,\quad
 \log|s_F|(v_{p,\triv})=-\infty,\quad
 u(v_{p,\triv})=-\infty .
\]

At $v_{r,\triv}$ for a closed point $r\in L\smallsetminus\{p\}$: the
form $\ell$ vanishes at $r$, so $v_{r,\triv}(\ell)=+\infty$ and
$\log|s_F|(v_{r,\triv})=-\infty$, whence also
$\chi(v_{r,\triv})=-\infty$; but $\mathfrak m_p$ is the unit ideal in
$\cO_{X,r}$, so
\[
 u(v_{r,\triv})=-v_{r,\triv}(\mathfrak m_p)=0 .
\]

Thus the pair $(\chi,\log|s_F|)$ takes the value
$(-\infty,-\infty)$ at both points, while the correct value of $u$ is
$-\infty$ at the first and $0$ at the second.  No rule assigning a value
to $(-\infty)-(-\infty)$ can produce both.  The subtraction
$u:=\chi-\log|s_F|$ is therefore not merely undefined by convention: it is
undefined because the information required to define it is absent from the
pair.  This is the exact analogue, on the reflection side, of
\Cref{ex:ch38-naive-realization-fails}, and it is why Step~6 of the proof of
\Cref{thm:bj-homogeneous-reflection} is performed by division rather than by
subtraction.
\end{example}

\begin{example}[A naive fibrewise definition is inconsistent]
\label{ex:ch39-naive-fibre-definition}
Suppose one tried to define the reflection of $\psi\in\PSH_Y(\eta)$ by
choosing, for each $x\in X^{\An}$, a preimage $w_x$ with
$\beta^{\An}(w_x)=x$ and setting $\chi(x):=\psi(w_x)$.  For
$\eta=\beta^*\theta$ this is consistent, by
\Cref{cor:ch39-fibre-constant} --- but that corollary is a consequence of
\Cref{thm:bj-homogeneous-reflection} and therefore unavailable at the
outset.  For a general $\eta$ the recipe is genuinely inconsistent: with
$Y=\operatorname{Bl}_p\mathbf P^2$, $q\in E$ and
$\psi'=\log|\mathcal I_q|$ as in \Cref{ex:ch39-not-fibre-constant}, the
two preimages $v_{q,\triv}$ and $v_{E,\triv}$ of $v_{p,\triv}$ give
\[
 \psi'(v_{q,\triv})=-\infty,
 \qquad
 \psi'(v_{E,\triv})=0 ,
\]
so the recipe would produce two different values at the single point
$v_{p,\triv}$.  The failure is invisible on $X^{\divis}$, where the
fibres are singletons by
\Cref{prop:ch39-fibres}\,\cref{itm:ch39-fib:val}; it is a proper-support
phenomenon exclusively.
\end{example}

\begin{warning}{Proper-support values are not recovered by density}
\label{warn:ch39-proper-support}
$X^{\divis}$ is dense in $X^{\An}$
(\Cref{thm:triv-divisorial-density}), and a first instinct is to define the
missing values by a limit along divisorial points.  This does not work, for
two independent reasons.  Evaluation at a proper-support point is not
continuous on $\PSH_X(\theta)$
(\Cref{ex:ch35-evaluation-discontinuous}); and a limit along a net of
divisorial points computes only an upper bound, since the functions involved
are upper semicontinuous, so it constrains the value from one side only
(\Cref{lem:ch35-density-gives} and the discussion after
\Cref{prop:ch39-determination-package}).  What does fix the value is
plurisubharmonicity together with the divisorial trace, by
\Cref{prop:ch39-determination-package}\,\cref{itm:ch39-det:equal}; and that
is a theorem about psh functions, not about dense subsets.
\end{warning}

\section{The exact functorial diagram}
\label{sec:ch39-diagram}

We collect the results as diagrams, and we label every arrow with what has
actually been proved about it: $\sim$ for a bijection,
$\hookrightarrow$ for an injection that is not asserted to be surjective,
$=$ for an identification of objects, and $\Rightarrow$ for an
implication of properties.  A dashed arrow marks an implication that is
\emph{not} proved in this chapter.

\subsection{Notation for the diagram}
\label{sec:ch39-diagram-notation}

Besides $\PSH_{\mathrm{hom}}$ and $\PSH$ we use two further objects.
First, the carrier-dominated subset
\[
 \PSH^{\mathrm{car}}_X(\theta)
 :=\bigl\{\phi\in\PSH_X(\theta)\ :\
 \exists\,p\in\PSH_{\mathrm{hom},X}(\theta),\ \exists\,C\in\mathbf R,\
 \phi\le p+C\bigr\},
\tag{39.52}\label{eq:ch39-carrier-subset}
\]
which contains $\PSH_{\mathrm{hom},X}(\theta)$ --- take $p=\phi$,
$C=0$ --- and which equals $\PSH_X(\theta)$ when $\theta$ is nef, by
\Cref{cor:ch39-nef-reflection}, or when $\operatorname{EP}_X(\theta)$
holds, by \Cref{lem:bj-canonical-carrier}.  Second, the set
\[
 \mathrm{Cart}^{\le0}_{\nef}(X,\theta)
 :=\bigl\{B\ \text{real Cartier b-divisor over }X\ :\
 B\le0,\ B\ \text{nef-total in }\theta\bigr\}
\]
of \Cref{def:bj-real-cartier-bdivisor} and \Cref{def:bj-nef-total}.

\begin{proposition}
\label{prop:ch39-carrier-subset}
$\beta^*$ restricts to a bijection
\[
 \beta^*:\PSH^{\mathrm{car}}_X(\theta)
 \xrightarrow{\ \sim\ }\PSH^{\mathrm{car}}_Y(\beta^*\theta),
\]
and $\PSH^{\mathrm{car}}_Y(\beta^*\theta)$ is contained in the image of
$\beta^*:\PSH_X(\theta)\to\PSH_Y(\beta^*\theta)$.
\end{proposition}

\begin{proof}
If $\phi\le p+C$ on $X^{\An}$ with
$p\in\PSH_{\mathrm{hom},X}(\theta)$, then
$\beta^*\phi\le\beta^*p+C$ with
$\beta^*p\in\PSH_{\mathrm{hom},Y}(\beta^*\theta)$ by
\Cref{thm:ch39-pullback}\,\cref{itm:ch39-pb:psh}; so $\beta^*$ maps
$\PSH^{\mathrm{car}}_X(\theta)$ into
$\PSH^{\mathrm{car}}_Y(\beta^*\theta)$, injectively.

Conversely let $u\in\PSH^{\mathrm{car}}_Y(\beta^*\theta)$, say
$u\le p_Y+C$ with $p_Y\in\PSH_{\mathrm{hom},Y}(\beta^*\theta)$.  By
\Cref{thm:bj-homogeneous-reflection} there is
$q\in\PSH_{\mathrm{hom},X}(\theta)$ with $\beta^*q=p_Y$, and by
\Cref{thm:bj-carrier-reflection} there is $v\in\PSH_X(\theta)$ with
$\beta^*v=u$.  Then $\beta^*v\le\beta^*(q+C)$, so $v\le q+C$ by
\Cref{thm:ch39-pullback}\,\cref{itm:ch39-pb:order}, and
$v\in\PSH^{\mathrm{car}}_X(\theta)$.  This proves surjectivity and the
last assertion.
\end{proof}

\subsection{The unconditional diagrams}
\label{sec:ch39-diagram-unconditional}

The b-divisor level is the source of everything, and it is where the two
varieties are literally identified:
\[
\begin{tikzcd}[column sep=huge,row sep=large]
\mathrm{Cart}^{\le0}_{\nef}(X,\theta)
 \arrow[r,equal,"\text{same objects}"]
 \arrow[d,hook,"B\mapsto\chi^X_B"']
&
\mathrm{Cart}^{\le0}_{\nef}(Y,\beta^*\theta)
 \arrow[d,hook,"B\mapsto\chi^Y_B"]
\\
\PSH_{\mathrm{hom},X}(\theta)
 \arrow[r,"\beta^*"," \sim"']
&
\PSH_{\mathrm{hom},Y}(\beta^*\theta)
\end{tikzcd}
\tag{39.53}\label{eq:ch39-bdivisor-square}
\]
The top identification is \Cref{lem:ch39-total-class}: antieffectivity and
nef-totality of a Cartier b-divisor are computed on a determination and do
not remember whether the base is $X$ or $Y$.  The vertical arrows are
the realization functors of \Cref{thm:bj-real-cartier-realization}; they are
injective, because a realization is determined by its divisorial trace and
distinct b-divisors have distinct traces
(\Cref{prop:ch38-evaluation-welldefined}\,\cref{itm:ch38-ev:linear}), and
they are \emph{not} surjective, by \Cref{rem:ch38-cartier-not-weil}.  The
square commutes by \Cref{prop:ch39-cartier-reflection}.  The bottom arrow is
\Cref{thm:bj-homogeneous-reflection}.

The psh level is the square announced in the introduction:
\[
\begin{tikzcd}[column sep=huge,row sep=large]
\PSH_{\mathrm{hom},X}(\theta)
 \arrow[r,"\beta^*"," \sim"']
 \arrow[d,hook]
&
\PSH_{\mathrm{hom},Y}(\beta^*\theta)
 \arrow[d,hook]
\\
\PSH^{\mathrm{car}}_X(\theta)
 \arrow[r,"\beta^*"," \sim"']
 \arrow[d,hook]
&
\PSH^{\mathrm{car}}_Y(\beta^*\theta)
 \arrow[d,hook]
\\
\PSH_X(\theta)
 \arrow[r,"\beta^*",hook]
&
\PSH_Y(\beta^*\theta).
\end{tikzcd}
\tag{39.54}\label{eq:bj-reflection-square}
\]
The two upper horizontal arrows are bijections, by
\Cref{thm:bj-homogeneous-reflection} and
\Cref{prop:ch39-carrier-subset}.  The bottom horizontal arrow is an
injection, by \Cref{thm:ch39-pullback}\,\cref{itm:ch39-pb:order}, and it is
\emph{not} asserted to be surjective: a point of the lower right corner is
known to lie in the image as soon as it has a homogeneous carrier, which is
the middle row, and \Cref{ex:ch39-not-fibre-constant} shows that psh
functions on $Y$ in other classes may fail to be pullbacks.  The vertical
arrows are the inclusions of \eqref{eq:ch39-carrier-subset}.  The top pair
is strict whenever $\theta$ is nef: the nonzero constants are then
$\theta$-psh by \Cref{lem:ch34-psh-nonempty}, they are carrier-dominated
by $0$, and they are not homogeneous, since $t\cdot c=tc$.  The bottom
pair is an equality whenever $\theta$ is nef --- again because $0$ is
then a homogeneous carrier for every psh function, as in the proof of
\Cref{cor:ch39-nef-reflection} --- and whenever
$\operatorname{EP}_Y(\beta^*\theta)$ holds, by
\Cref{lem:bj-canonical-carrier} and \Cref{prop:ch39-carrier-subset}.

\subsection{The conditional diagram}
\label{sec:ch39-diagram-conditional}

Abbreviate by $(\mathrm C)$ the statement ``the function
\eqref{eq:bj-canonical-carrier} is a homogeneous $\beta^*\theta$-psh
carrier for every $u\in\PSH_Y(\beta^*\theta)$'' and by $(\mathrm B)$ the
statement ``$\beta^*:\PSH_X(\theta)\to\PSH_Y(\beta^*\theta)$ is
a bijection''.  The logical ledger is
\[
\begin{gathered}
 \operatorname{EP}_Y(\beta^*\theta)
 \overset{\mathrm{(a)}}{\Longrightarrow}(\mathrm C)
 \overset{\mathrm{(b)}}{\Longrightarrow}(\mathrm B),\\
 \operatorname{EP}_Y(\beta^*\theta)\ \text{and}\ (\mathrm B)
 \overset{\mathrm{(c)}}{\Longrightarrow}
 \operatorname{EP}_X(\theta),\\
 \operatorname{EP}_X(\theta)
 \dashrightarrow \operatorname{EP}_Y(\beta^*\theta)
 \qquad\text{not proved here}.
\end{gathered}
\tag{39.55}\label{eq:bj-ep-reflection-diagram}
\]
Implication $\mathrm{(a)}$ is \Cref{lem:bj-canonical-carrier}; it is the
first use of the source envelope hypothesis, and it is a use of an
\emph{assumption}.  Implication $\mathrm{(b)}$ is
\Cref{thm:bj-carrier-reflection}, applied to every $u$ at once, together
with the injectivity of \Cref{thm:ch39-pullback}.  Bijection $(\mathrm B)$
alone does \emph{not} imply $\operatorname{EP}_X(\theta)$.  Implication
$\mathrm{(c)}$ uses the source envelope hypothesis a second time: it
supplies the increasing-net limit upstairs through
\Cref{prop:bj-increasing-net-criterion}, after which $(\mathrm B)$ reflects
that limit and its divisorial trace downstairs.  The dashed arrow is
discussed in \Cref{rem:ch39-one-way}: it becomes true once
\Cref{thm:bj-final-envelope} is available, and assuming it here would be
circular.

\begin{proposition}
\label{prop:ch39-diagram}
With the notation above, and for $\beta:Y\to X$ and
$\theta\in N^1(X)_{\mathbf R}$ as in \eqref{eq:ch39-standing}:
\begin{enumerate}
\item\label{itm:ch39-diag:hom} the map $\beta^*$ on homogeneous psh
 classes is a bijection, unconditionally, and is compatible with order,
 finite maxima, convex combinations, scaling, Legendre slices and
 divisorial traces
 (\Cref{cor:ch39-hom-structure});
\item\label{itm:ch39-diag:car} the map $\beta^*$ on carrier-dominated psh
 classes is a bijection, unconditionally
 (\Cref{prop:ch39-carrier-subset});
\item\label{itm:ch39-diag:full} the map $\beta^*$ on full psh classes is
 an injection whose image contains
 $\PSH^{\mathrm{car}}_Y(\beta^*\theta)$; it is a bijection if
 $\theta\in\Nef(X)$ or if $\operatorname{EP}_Y(\beta^*\theta)$ holds;
\item\label{itm:ch39-diag:ep} $\operatorname{EP}_Y(\beta^*\theta)$
 implies $\operatorname{EP}_X(\theta)$, and the converse is not proved
 here;
\item\label{itm:ch39-diag:bdiv} the realization functors intertwine
 $\beta^*$ with the identity on Cartier b-divisors, and are injective but
 not surjective.
\end{enumerate}
\end{proposition}

\begin{proof}
Each item has been proved above:
\cref{itm:ch39-diag:hom} is \Cref{thm:bj-homogeneous-reflection} and
\Cref{cor:ch39-hom-structure}; \cref{itm:ch39-diag:car} is
\Cref{prop:ch39-carrier-subset}; \cref{itm:ch39-diag:full} combines
\Cref{thm:ch39-pullback}, \Cref{prop:ch39-carrier-subset},
\Cref{cor:ch39-nef-reflection} and
\Cref{thm:bj-birational-ep-descent}; \cref{itm:ch39-diag:ep} is
\Cref{thm:bj-birational-ep-descent} and \Cref{rem:ch39-one-way}; and
\cref{itm:ch39-diag:bdiv} is \Cref{lem:ch39-total-class},
\Cref{prop:ch39-cartier-reflection} and
\Cref{rem:ch38-cartier-not-weil}.
\end{proof}

\paragraph{Proof checkpoint.}
Homogeneous reflection comes from real Cartier realization and the exact
identity \eqref{eq:bj-reflection-total}, not from an envelope.  This is the
logical fact that makes \eqref{eq:bj-ep-reflection-diagram} noncircular: the
horizontal arrows of \eqref{eq:ch39-bdivisor-square} and of the top two rows
of \eqref{eq:bj-reflection-square} are proved before any envelope hypothesis
is written down.  The source envelope hypothesis first appears in
implication $\mathrm{(a)}$ to create carriers and is used again in
implication $\mathrm{(c)}$ to create the increasing-net limit; neither use
imports the envelope theorem as a previously established result.

\paragraph{Proof checkpoint.}
Carrier reflection uses $u\le p+C$, not a bound by a homogeneous constant
function.  The carrier controls the far Legendre tail, and the summable grid
correction turns the compact-window supremum into a decreasing psh limit.
When $\theta$ is nef the constant function $0$ is available as a
carrier and the hypothesis is vacuous
(\Cref{cor:ch39-nef-reflection}); the content of
\Cref{thm:bj-carrier-reflection} lies in the non-nef range where
$\PSH_X(\theta)$ is nonempty; there the class is pseudoeffective but
contains no bounded psh function.

\section{Chapter dictionary}
\label{sec:ch39-dictionary}

The correspondences established or used in this chapter are the following.
Each row records a statement together with the hypotheses under which it was
proved.  Rows one and two are the topological dictionary of
\Cref{sec:ch39-analytifications}; rows three to six are the reflection
statements; rows seven and eight record what governs the passage to proper
support and the direction of the envelope implication.

\begin{center}
\begin{tabular}{p{.29\textwidth}p{.31\textwidth}p{.29\textwidth}}
\toprule
analytic datum & birational datum & hypothesis\\
\midrule
$\beta^{\An}$ on $Y^{\divis}$ &
the identity of $X^{\divis}$ with $Y^{\divis}$, via $k(X)=k(Y)$ &
$\beta$ projective birational
(\Cref{lem:ch39-beta-an}\,\cref{itm:ch39-ban:val})\\
the fibre of $\beta^{\An}$ over $v_{p,\triv}$, $p\in X$ a closed
point &
$\bigl(\beta^{-1}(p)\bigr)^{\An}$, all of whose points have proper
support &
$\beta$ proper birational
(\Cref{prop:ch39-fibres}\,\cref{itm:ch39-fib:point})\\
$\beta^*$ on homogeneous psh classes, a bijection &
the same Cartier b-divisors at every realized stage; general functions
come from the interpolated decreasing limit &
$X,Y$ normal integral projective; no positivity of $\theta$; no
envelope property
(\Cref{thm:bj-homogeneous-reflection},
\Cref{rem:ch39-cartier-not-enough})\\
$\beta^*$ on carrier-dominated psh classes, a bijection &
the compact Legendre window plus homogeneous reflection &
a homogeneous carrier in the same class
(\Cref{thm:bj-carrier-reflection}, \Cref{prop:ch39-carrier-subset})\\
$\beta^*$ on all psh classes, an injection &
image contains the carrier-dominated part &
none
(\Cref{thm:ch39-pullback}, \Cref{prop:ch39-carrier-subset})\\
$\beta^*$ on all psh classes, a bijection &
--- &
$\theta\in\Nef(X)$ (\Cref{cor:ch39-nef-reflection}), or
$\operatorname{EP}_Y(\beta^*\theta)$
(\Cref{thm:bj-birational-ep-descent})\\
equality on $X^{\divis}$ &
equality on $X^{\An}$, including at proper support &
\emph{both} functions psh, in possibly different classes
(\Cref{prop:ch39-determination-package}\,\cref{itm:ch39-det:equal})\\
$\operatorname{EP}_Y(\beta^*\theta)$ &
$\operatorname{EP}_X(\theta)$ &
one-way as proved here
(\Cref{thm:bj-birational-ep-descent}, \Cref{rem:ch39-one-way})\\
\bottomrule
\end{tabular}
\end{center}

Five separations should be carried forward.

\emph{Divisorial point versus exceptional fibre.}  The exceptional divisor
of $\beta$ contributes one ray of $X^{\divis}=Y^{\divis}$, on which
$\beta^{\An}$ is a bijection; the exceptional fibre of $\beta^{\An}$
over the trivial semivaluation of the blown-up centre is a different object,
lying entirely in the proper-support locus
(\Cref{ex:ch39-fibre-plane}).  Statements proved on $X^{\divis}$ say
nothing directly about the latter, and the passage between them is
\Cref{prop:ch39-determination-package}.

\emph{Pulled-back class versus arbitrary class.}  Homogeneous psh functions
in a pulled-back class are constant on the fibres of $\beta^{\An}$
(\Cref{cor:ch39-fibre-constant}); in other classes they need not be
(\Cref{ex:ch39-not-fibre-constant}).  Fibrewise constancy is a conclusion,
never a hypothesis, and a proof that begins by defining the reflection
fibrewise is circular (\Cref{ex:ch39-naive-fibre-definition}).

\emph{Orientation of the drift versus its size.}  The approximating classes
on $Y$ converge to $\beta^*\theta$, so the drift is arbitrarily small;
that is not why it is harmless.  It is harmless because the correction
\eqref{eq:bj-reflection-class-correction} is arranged to be \emph{ample},
which lets \eqref{eq:bj-reflection-total} exhibit the total class as a sum
of two nef classes, and because $\omega$ is a single fixed class in which
every stage lives.  The distinction is
\Cref{warn:ch38-drift-not-absorbed}.

\emph{Division versus subtraction.}  The fixed pole is removed by
homogeneous Siu division, \Cref{thm:ch35-homogeneous-siu-general}, which
produces an identity of sums.  The corresponding difference is undefined at
proper support and cannot be repaired by any convention
(\Cref{warn:ch39-no-subtraction}, \Cref{ex:ch39-subtraction-fails}).

\emph{Envelope property as hypothesis versus as theorem.}  Every occurrence
of $\operatorname{EP}$ in this chapter is a displayed hypothesis on the
\emph{source} of the morphism.  The conclusion,
$\operatorname{EP}_X(\theta)$, carries no hypothesis on $X$ beyond
normality and projectivity.  Reading the final theorem
\Cref{thm:bj-final-envelope} back into
\Cref{lem:bj-canonical-carrier} would make the whole of
\Crefrange{chap:bj-cartier-realization}{chap:bj-subintegral-normalization}
circular (\Cref{warn:ch39-firewall}, \Cref{warn:ch39-canonical-carrier}).

\section{Source note}

This chapter follows the homogeneous reflection, carrier reflection, and
canonical-carrier sections of the supplied envelope manuscript.  The
regularization, comparison, Legendre reconstruction, and homogeneous Siu
division interfaces are from \cite{BJ22}.  We retained the rational-square
class correction, the exact drift cancellation, the lower-grid argument,
and the global treatment of proper-support poles.

The precise external attributions are as follows.  The comparison theorem
used at every promotion from $X^{\divis}$ to $X^{\An}$ is
\cite[Theorem~4.22]{BJ22}, with its determination corollary
\cite[Corollary~4.23]{BJ22}; both statements are proved in full in
\Cref{chap:bj-homogeneity-bdivisors} and are quoted here through
\Cref{prop:ch39-determination-package}.  The Legendre decomposition
\eqref{eq:bj-legendre-reconstruction} and the behaviour of the slices in the
parameter are \cite[Lemma~6.23 and Theorem~6.24]{BJ22}.  The homogeneous
regularization behind \Cref{cor:ch38-gf-regularization-bdivisor} is
\cite[Corollary~6.17]{BJ22}, and homogeneous Siu division is
\cite[Theorem~6.21]{BJ22}; the warning after
\Cref{thm:ch35-homogeneous-siu-general} distinguishes it from the
inhomogeneous statement \cite[Theorem~6.31]{BJ22}, which does require the
envelope property and which is therefore not used here.  The envelope
property was formulated in \cite{BJ22} as Boucksom--Jonsson's
Conjecture~5.14; \cite{BJ22b} records partial results available in that
literature.  The proof developed in this book is
\Cref{thm:bj-final-envelope}, with historical context in
\Cref{chap:bj-envelope-history}.  The domination of a modification by a blowup, used inside
\Cref{lem:ch39-fixed-pole} through \Cref{lem:ch38-activation}, is
\cite[\href{https://stacks.math.columbia.edu/tag/081T}{Tag 081T}]{stacks-project},
imported as \Cref{thm:appB-admissible-domination}; the twist theorem
converting relative into absolute ampleness is \cite[4.6.13]{EGAII}, used
through \Cref{lem:ch38-numerical-birational}\,\cref{itm:ch38-num:relative};
Zariski's main theorem in the form used for
\Cref{lem:ch38-numerical-birational}\,\cref{itm:ch38-num:notample} is
\cite[Corollary~III.11.4]{Har}; and the Nakai--Moishezon criterion used in
the surface computations of \Cref{ex:ch39-blowup-reflection-numbers} and
\Cref{ex:ch39-carrier-pair} is \cite[Theorem~1.2.23]{PAG1}.  The b-divisor
formalism is that of \Cref{app:intersection-bdivisors}, ultimately
\cite{BFJ09}.

Six things in this chapter are not in the sources in this form and are
supplied here.  The closure calculus is supplemented by the
monotonicity in the class as \Cref{lem:ch39-closure-any-class}, on the basis
of the
elementary monotonicity \Cref{lem:ch39-hgf-monotone} of the Fubini--Study
spaces in the line bundle; this is needed because the classes to which
\Cref{chap:bj-alterations} applies the descent theorem are arbitrary real
classes, and because, after the correction of \Cref{def:bj-psh}, a
homogeneous psh function can exist in a class that is pseudoeffective but
not nef.
The pullback statement \Cref{thm:ch39-pullback} is isolated and proved
directly from the varying-class form of \Cref{def:bj-psh}, which is what
makes \Cref{rem:ch39-closing-the-gap} possible.  The fixed-pole data is
packaged for a birational morphism as a free-standing lemma,
\Cref{lem:ch39-fixed-pole}, quoting only those items of
\Cref{lem:ch38-activation} whose proof is independent of the total class.
The choice of the interpolation parameter is isolated as
\Cref{lem:ch39-drift-window}, with the subnet formalism of
\Cref{def:appE-willard-subnet} made explicit.  The
carrier-dominated subspace \eqref{eq:ch39-carrier-subset} and the
observation \Cref{cor:ch39-nef-reflection} that the carrier hypothesis is
vacuous in a nef class do not appear in the sources and are what make the
diagram \eqref{eq:bj-reflection-square} exact rather than merely suggestive.
Finally the worked computations
\Cref{ex:ch39-blowup-reflection-numbers},
\Cref{ex:ch39-carrier-pair},
\Cref{ex:ch39-usc-not-determined-birational} and
\Cref{ex:ch39-subtraction-fails} are new, and are included because the
audit items attached to this chapter --- the passage from divisorial
equality to global equality, and the values at proper-support
semivaluations --- are precisely the points at which a plausible argument
can be wrong without being visibly wrong.
